\documentclass[12pt, a4paper]{article}

\usepackage[utf8]{inputenc}
\usepackage[ngerman,english]{babel} % English is the main language; ngerman for the statutory declaration
\usepackage[a4paper, width=160mm, top=35mm, bottom=30mm, bindingoffset=0mm]{geometry}
\usepackage{graphicx}
\usepackage[table]{xcolor} % table: row colours in the included GUIDE-LLM checklist (app:guide_llm)
\usepackage{amsmath,amssymb,amsthm}
\usepackage{ragged2e}
\usepackage{fancyhdr}
\usepackage{sectsty}
\allsectionsfont{\raggedright}
\usepackage{etoolbox}
\preto{\section}{\clearpage}
\usepackage{needspace} % keeps the Results block of the table of contents on one page (see \addtocontents before \section{Results})
\usepackage{csquotes}
\usepackage[backend=biber, style=apa, sorting=nyt]{biblatex}
\DeclareLanguageMapping{english}{english-apa}
\addbibresource{bibliography.bib}
\addbibresource{literature.bib}
\usepackage{listings}
\usepackage{booktabs}
\usepackage{array}
\usepackage{longtable}
\usepackage{threeparttable}
\usepackage{chngcntr}
\usepackage{subcaption}
\usepackage{tikz}
\usetikzlibrary{arrows.meta,positioning,fit,patterns,backgrounds}
\usepackage{pgfplots}
\pgfplotsset{compat=1.18}
\usepgfplotslibrary{dateplot}
\usepackage[scaled=0.92]{helvet}
\usepackage[eulergreek]{sansmath}
\captionsetup{labelsep=period, labelfont=bf}
\definecolor{figInk}{HTML}{000000}
\definecolor{figModel}{HTML}{0072B2}
\definecolor{figAbove}{HTML}{E69F00}
\definecolor{figChanged}{HTML}{D55E00}
\definecolor{figAdjusted}{HTML}{009E73}
\definecolor{figModelLight}{HTML}{56B4E9}
\definecolor{figGerman}{HTML}{CC79A7}
\definecolor{figEnglish}{HTML}{F0E442}
\definecolor{figGreyA}{HTML}{666666}
\definecolor{figGreyB}{HTML}{A6A6A6}
\definecolor{figGreyC}{HTML}{E0E0E0}
\colorlet{figBlue}{figModel}
\colorlet{figVermillion}{figChanged}
\colorlet{figBelow}{figModel}
\colorlet{figHighlight}{figChanged}
\colorlet{figContext}{figGreyB}
\newcommand{\figfont}{\sffamily\sansmath\fontsize{9.8}{12}\selectfont\frenchspacing}
\newcommand{\figfontsmall}{\sffamily\sansmath\fontsize{8.7}{10.5}\selectfont\frenchspacing}
\pgfplotsset{
  thesisfig/.style={
    font=\figfont, label style={font=\figfont}, title style={font=\figfont},
    every tick label/.append style={font=\figfontsmall},
    legend style={font=\figfontsmall, draw=none, fill=none, at={(0.5,0)}, anchor=north, yshift=-13mm,
      /tikz/every even column/.append style={column sep=10pt}},
    legend cell align=left,
    axis lines*=left, axis line style={line width=0.5pt}, grid=none,
    tick style={line width=0.5pt, black}, tick align=outside, major tick length=3pt,
    every axis plot/.append style={line width=0.8pt},
    error bars/error bar style={line width=0.8pt},
    error bars/error mark options={rotate=90, mark size=1.6pt, line width=0.8pt},
    mark size=2.2pt,
    /pgf/number format/use period, /pgf/number format/1000 sep={,},
    clip=false,
  },
  rowaxis/.style={y axis line style={draw=none}, ytick style={draw=none}},
  pooled/.style={only marks, mark=*, draw=figModel, fill=figModel, mark options={draw=figModel, fill=figModel}},
  german/.style={only marks, mark=square*, draw=figGerman, mark options={draw=figGerman, fill=figGerman}, mark size=2.0pt,
    error bars/error mark options={rotate=90, mark size=1.6pt, line width=0.8pt, draw=figGerman}},
  english/.style={only marks, mark=triangle*, draw=figGreyA, mark options={draw=figGreyA, fill=figEnglish, line width=0.5pt}, mark size=2.7pt,
    error bars/error mark options={rotate=90, mark size=1.6pt, line width=0.8pt, draw=figGreyA}},
  registered/.style={only marks, mark=diamond*, draw=figModel, mark options={draw=figModel, fill=figModel}, mark size=3.0pt},
  highlight/.style={only marks, mark=pentagon*, draw=figHighlight, mark options={draw=figHighlight, fill=figHighlight}, mark size=2.5pt},
  highlightb/.style={only marks, mark=triangle*, draw=figHighlight, mark options={draw=figHighlight, fill=figHighlight}, mark size=2.9pt},
  context/.style={only marks, mark=*, draw=figContext, mark options={draw=figContext, fill=figContext}},
  open/.append style={mark options={fill=white, line width=0.8pt}},
  precisionmark/.style={only marks, mark=square*, mark size=1.9pt, draw=figModel, mark options={draw=figModel, fill=white, line width=0.7pt}},
  recallmark/.style={only marks, mark=triangle*, mark size=2.4pt, draw=figModel, mark options={draw=figModel, fill=white, line width=0.7pt, yshift=0.6pt, rotate=180}},
}
\tikzset{
  refline/.style={figGreyB, dashed, line width=0.5pt},
  guideline/.style={figGreyC, line width=0.5pt},
  panellabel/.style={anchor=base west, inner sep=0pt, font=\sffamily\bfseries\fontsize{11}{13}\selectfont},
  figtext/.style={font=\figfontsmall},
}
\usepackage{pifont}
\usepackage{docmute}
\providecommand{\checked}{\ding{110}}   % filled box
\providecommand{\unchecked}{\ding{111}} % empty box
\newcommand{\inthesis}{}
\usepackage{hyperref}
\DeclareUrlCommand\modelid{\urlstyle{tt}\def\UrlBreaks{\do\.\do\-\do\:\do\/\do\_}}
\makeatletter\renewcommand{\@pnumwidth}{4.4em}\makeatother
\setcounter{biburlnumpenalty}{9000}
\setcounter{biburlucpenalty}{9000}
\setcounter{biburllcpenalty}{9000}
\DeclareSourcemap{
  \maps[datatype=bibtex]{
    \map[overwrite]{
      \pertype{misc}
      \step[fieldsource=note, matchi=\regexp{preprint|peer.reviewed|arxiv}, final]
      \step[fieldsource=eprinttype, matchi=\regexp{\Aarxiv\Z}, final]
      \step[fieldset=organization, fieldvalue={arXiv}]
    }
    \map[overwrite]{
      \pertype{misc}
      \step[fieldsource=note, matchi=\regexp{preprint|peer.reviewed|arxiv}, final]
      \step[fieldset=titleaddon, fieldvalue={Preprint}]
      \step[fieldset=eprint, null]
      \step[fieldset=eprinttype, null]
      \step[fieldset=eprintclass, null]
      \step[typesource=misc, typetarget=online]
    }
    \map[overwrite]{
      \step[fieldsource=note, final]
      \step[fieldset=note, null]
    }
  }
}
\renewcommand{\topfraction}{.85}
\renewcommand{\bottomfraction}{.6}
\renewcommand{\textfraction}{.12}
\renewcommand{\floatpagefraction}{.8}
\makeatletter\setlength{\@fptop}{0pt}\setlength{\@fpsep}{\floatsep}\makeatother

\graphicspath{{images/}}

\newcommand{\mytitle}{Extracting Technology Frames from Renewable Energy News with Large Language Models}
\newcommand{\myname}{Nils Wagner}
\newcommand{\authorfix}[1]{\textcolor{red}{\textbf{#1}}}
\newcommand{\myexaminer}{Prof.\ Dr.\ Stefan Feuerriegel}
\newcommand{\mysupervisor}{Kerstin Forster}
\newcommand{\mymatriculation}{12554615}
\newcommand{\mycompletiondate}{\iflanguage{ngerman}{25.11.2026}{25 November 2026}}
\newcommand{\guideref}{\ifcsname r@app:guide_llm\endcsname Appendix~\ref{app:guide_llm}\else\authorfix{[TODO: GUIDE-LLM appendix not yet embedded]}\fi}

\newcommand{\changefont}{\fontsize{8}{8}\selectfont}
\pagestyle{fancy}
\fancyhead{}
\fancyhead[R]{\changefont{\mytitle}}
\fancyfoot{}
\fancyfoot[R]{\thepage}
\setlength{\headheight}{14.5pt}
\setlength{\parindent}{0pt}
\setlength{\parskip}{6pt}
\interfootnotelinepenalty=10000
\counterwithout{footnote}{section}

\makeatletter
\renewcommand\@makefntext[1]{\parindent 1em\noindent\scriptsize\@makefnmark#1}
\makeatother

\hypersetup{
    pdftitle={\mytitle},
    pdfauthor={\myname},
    pdfsubject={Bachelor's thesis, Munich School of Management, Ludwig-Maximilians-Universität München},
    pdfkeywords={technology frames, framing analysis, large language models, automated content analysis, validation, renewable energy, news media},
    colorlinks=true,
    linkcolor=black,
    urlcolor=blue,
    citecolor=blue,
    bookmarksdepth=2
}

\lstset{
    language=Python,
    basicstyle=\ttfamily\footnotesize,
    keywordstyle=\color{blue}\bfseries,
    stringstyle=\color{red},
    commentstyle=\color{gray}\itshape,
    numbers=left,
    numberstyle=\tiny\color{gray},
    stepnumber=1,
    numbersep=8pt,
    backgroundcolor=\color{gray!5},
    frame=single,
    rulecolor=\color{black!30},
    breaklines=true,
    breakatwhitespace=true,
    showstringspaces=false,
    tabsize=4,
    captionpos=b,
    aboveskip=10pt,
    belowskip=10pt
}

\theoremstyle{definition}
\newtheorem{definition}{Definition}[section]
\newtheorem{example}{Example}[section]

\newcommand{\code}[1]{\texttt{#1}}
\newcommand{\pkg}[1]{\textsf{#1}}
\newcommand{\file}[1]{\texttt{#1}}

\begin{document}

\pagenumbering{Roman} % Roman from the title page on: avoids a second page.1 anchor (the titlepage environment resets the counter)
\begin{titlepage}
\begin{center}

\rule{\textwidth}{1.5pt}\par
\vspace{0.3cm}
{\LARGE\textbf{\mytitle}\par}
\vspace{0.3cm}
\rule{\textwidth}{1.5pt}

\vspace{0.8cm}

{\Large\textbf{Bachelor's Thesis}\par}

\vspace{0.6cm}

{\large submitted for the degree of Bachelor of Science\\
at the Munich School of Management of\\
Ludwig-Maximilians-Universität München\par}

\vfill

\includegraphics[width=0.24\textwidth]{sigillum.png}

\vfill

\end{center}

{\large
\noindent
\begin{minipage}[t]{0.47\textwidth}
\textbf{Author:}\\[0.25cm]
\myname\\[0.25cm]
Student Number: \mymatriculation
\end{minipage}\hfill
\begin{minipage}[t]{0.47\textwidth}
\textbf{Examiner:}\\[0.25cm]
\myexaminer\\[0.4cm]
\textbf{Supervisor:}\\[0.25cm]
\mysupervisor
\end{minipage}

\vspace{0.7cm}

\noindent\hfill
\begin{minipage}[t]{0.47\textwidth}
Munich, \mycompletiondate
\end{minipage}\par}

\vfill

\begin{center}
\vspace*{-\topsep}\footnotesize
\textbf{Document Statistics:} This thesis contains 69,537 characters (including spaces, counted over Chapters 1--8: running text and headings, footnotes, tables, and figure and table captions).
\end{center}
\end{titlepage}

\pagenumbering{Roman}
\setcounter{page}{2}

\newpage
\begin{abstract}
\noindent
Firms and agencies monitoring energy-technology news could use language-model frame shares, but accuracy-validated labels can bias them. No reviewed study validates one such instrument at frame-element level, per language and on shares. Pre-specified rules compare a frozen-codebook pipeline with one coder on 120 German and English paragraphs, probability-sampled from a keyword-prefiltered, purposive corpus. Against the first coding, technology detection (micro-F1 0.935) and shared-frame evaluations (accuracy 0.874) agree closely, frame elements less (micro-F1 0.643). The model marks fewer elements. Most shares lie lower, investment's by 18.3 percentage points. Against 40 paragraphs re-coded before the results, later checks find this recall deficit smaller but persisting in point estimate. In exploratory checks against the full, mostly post-result re-code, investment's share differs by $-13.6$ [$-20.4$, $-7.6$] points. Neither coding is established as correct. Projects with their own frames can reuse the validation workflow. Shares need a human-coded probability sample and a defended coding convention.
\end{abstract}

\newpage
\makeatletter\renewcommand{\@tocrmarg}{2.55em plus1fil}\makeatother
\tableofcontents

\newpage
\phantomsection\addcontentsline{toc}{section}{List of Figures}
\listoffigures

\newpage
\phantomsection\addcontentsline{toc}{section}{List of Tables}
\listoftables

\newpage
\section*{List of Abbreviations}
\addcontentsline{toc}{section}{List of Abbreviations}
\begin{tabular}{@{}l>{\raggedright\arraybackslash}p{118mm}@{}}
API & application programming interface \\
AUC & area under the receiver operating characteristic curve \\
CI & confidence interval \\
DSL & design-based supervised learning \parencite{egami2023dsl} \\
EU & European Union \\
F1 & harmonic mean of precision and recall \\
FN, FP, TP & false negative, false positive, true positive \\
GDELT & Global Database of Events, Language, and Tone \\
GPT & generative pre-trained transformer (OpenAI model family) \\
GUIDE-LLM & guidelines for the use of LLMs in behavioural and social science, a reporting checklist \parencite{feuerriegel2026guidellm} \\
LLM & large language model \\
NLP & natural language processing \\
RQ & research question \\
SE & standard error \\
SHA-256 & Secure Hash Algorithm with a 256-bit digest (file fingerprints) \\
SemEval & International Workshop on Semantic Evaluation (shared tasks) \\
USD & United States dollar \\
\end{tabular}

\newpage
\newcounter{frontmatterpages}\setcounter{frontmatterpages}{\value{page}}
\pagenumbering{arabic}
\setcounter{page}{1}

\section{Introduction}\label{sec:introduction}
News coverage links energy technologies to costs, risks and promises \parencite[e.g.,][]{stephens2009windmedia,ali2025hydrogenframing}. The shares of such framings, or frame shares, describe coverage, not acceptance, but can inform firms, investors and agencies monitoring a technology's public debate. Policy needs such evidence, which is slow and costly to collect. Large language models (LLMs) could therefore help \parencite[pp.~1--2]{toetzke2023llmclimatetech}. Yet their frame shares are useful only if the instrument behind them is validated.

\paragraph{Motivation.}
Such validation is rarely complete. Many content-analysis validations omit human-label reliability \parencite{song2020validations}. Three computational framing studies translate multilingual material or divide it between models \parencite{toetzke2022energydiscourse,jumle2024findingframes,guo2023openframing}, and none validates one instrument in each language. Validations typically report accuracy or agreement \parencite[p.~2]{linder2026multicalibrated}. Yet labels used directly can substantially bias statistics computed from them, such as shares, even at 80 to 90\% accuracy \parencite{egami2023dsl}. No reviewed study validates one LLM instrument at frame-element level, per language and on aggregate shares (Chapter~\ref{sec:related_work}).

\paragraph{Research Questions.}
The guiding research question is how a multilingual pipeline based on an LLM can be designed and validated for extracting technology frames from international renewable energy news. A technology frame is a news frame attached to a named energy technology. It is measured through frame elements, concrete presence questions such as whether a paragraph links wind power to its landscape impact. The elements marked present give the frames, and each frame carries an evaluation (Chapter~\ref{sec:background}, Figure~\ref{fig:anatomy}). Four research questions (RQ1 to RQ4) operationalise the guiding question:
\begin{enumerate}
\item[RQ1] How should such a pipeline be designed, and why is each part the way it is?
\item[RQ2] How well do its labels agree with a human reference, per language and per label dimension?
\item[RQ3] How far do identical repeated executions of the pipeline differ from one another?
\item[RQ4] How far do the aggregate prevalences it would report differ from those of the human reference, and where does the sample size limit what can be reported?
\end{enumerate}

\paragraph{Objective and Approach.}
One frozen codebook of coding rules supplies definitions for Claude Sonnet 4.6's system prompt and output schema and for the coding sheets of a single human coder (the author). The pipeline submits each paragraph untranslated, in German or English, the two languages the coder reads. Evaluation uses a stratified probability sample of 120 paragraphs, 60 per language, design-weighted to the eligible frame. This sampling frame (not a news frame) holds 618 paragraphs from a purposive, keyword-prefiltered corpus (Figure~\ref{fig:overview}a). The first of three identical model executions is scored against the coder's first coding. With the first labels hidden, the coder re-coded 40 paragraphs before the model run, as registered. A supplementary re-code covered the other 80 after the results were known; analyses using it are exploratory. Together they form the second coding. \emph{Registered} refers to an internal pre-analysis specification recorded before evaluation coding and prediction, not a public registry (Appendix~\ref{app:protocol}).

% 2026-11-09, second slimming pass (G-70): caption articles and legend wording tightened; no key removed.
% Figure 1 (Introduction) -- the study in one picture. Panel a is drawn by hand in TikZ: the evaluation design, lean
% (restructuring plan v3, 2026-11-05, Fig 1a): the 618-paragraph eligible frame, the 120-paragraph probability sample,
% the model and the single coder applying the same frozen codebook, and the comparison; one dashed marker for the
% analyses added after the registered result (listed in Table tab:added_analyses, not here). The system architecture
% that implements this design is images/architecture.tex. Design counts only, no measured result; every count is
% checked against the records: 618 paragraphs and 120 sampled (tools/runs/tables/dsl/dsl_prevalence.json, N and
% sampled), 60 German and 60 English (sample manifest; tab_primary), three executions with the first scored
% (Section research design), 40 paragraphs re-coded blind, 20 per language (Section on the human reference).
% Panel b is generated by tools/pipeline/figures.py results from the registered metrics (images/fig_headline.tex);
% no number in panel b is typed by hand. Style from the thesis preamble (figure style block); no internal codes
% (Chapter 1). The two panels sit flush left in one column, so that each panel letter, the left edge of its picture,
% lines up with the other.
% Deliberate deviation from the figure style specification (2026-11-05): panel b shows no displacement row; the
% largest downward displacement (18 points) is stated in the Findings paragraph beside the figure, and the Chapter 1
% budget leaves no room for a fourth row with its caption text. The author decides whether to add it.
% Palette 2026-11-06: panel a carries the thesis colour convention (preamble palette block). The model box is outlined
% in figModel (blue, 0.8pt, white fill), the coder box stays figInk, the comparison box, which combines both, is
% outlined in figAdjusted (green, 0.8pt) over its grey fill, the "later analyses" box and its arrow are figChanged
% (vermillion), the arrow from the model into the comparison is figModel; the arrow from the coder stays figInk. The
% frame and sample boxes stay black and grey. Every colour keeps its non-colour cue: the direct label in each box, the
% heavier outline on the two coloured boxes, the grey fill on the comparison, the dashed line and italic text on the
% later-analyses marker, and the position of each arrow at its source box. No count, label or layout changed; the
% caption names no colour in the panel descriptions and closes with the thesis-wide colour key.
\begin{figure}[htbp]
\centering
\begin{tabular}{@{}l@{}}
\begin{tikzpicture}[
  font=\figfontsmall,
  node distance=5mm and 7mm,
  pstep/.style={draw=figInk, line width=0.5pt, rectangle, rounded corners=1.5pt, minimum height=11mm,
    execute at begin node={\hyphenpenalty=10000\exhyphenpenalty=10000},
    align=center, inner sep=3pt, fill=white},
  key/.style={pstep, fill=figGreyC},
  modelbox/.style={pstep, draw=figModel, line width=0.8pt},
  cmpbox/.style={key, draw=figAdjusted, line width=0.8pt},
  later/.style={draw=figChanged, dashed, line width=0.5pt, rectangle, rounded corners=1.5pt, align=center,
    inner sep=3pt, text=figChanged, font=\figfontsmall\itshape,
    execute at begin node={\hyphenpenalty=10000\exhyphenpenalty=10000}},
  arr/.style={-{Latex[length=1.8mm]}, line width=0.6pt},
  darr/.style={arr, figChanged, dashed},
]
\node[pstep, text width=23mm] (frame) {Eligible frame\\618 paragraphs};
\node[key, text width=27mm, right=6mm of frame] (sample) {Probability sample\\120 paragraphs\\(60 German, 60 English)};
\node[modelbox, text width=36mm, above right=1.5mm and 10mm of sample.east] (model) {Model, codebook as prompt:\\three identical executions,\\the first scored};
\node[pstep, text width=36mm, below right=1.5mm and 10mm of sample.east] (coder) {Single coder, same codebook:\\reference labels;\\blind re-code of 40};
\node[cmpbox, text width=26mm, anchor=west] (cmp) at ([xshift=10mm]sample.east -| model.east) {Comparison:\\agreement, stability,\\displacement};
\draw[arr] (frame) -- (sample);
\draw[arr] (sample.east) -- ++(4mm,0) |- (model.west);
\draw[arr] (sample.east) -- ++(4mm,0) |- (coder.west);
\draw[arr, draw=figModel] (model.east) -- ++(4mm,0) |- ([yshift=1.5mm]cmp.west);
\draw[arr] (coder.east) -- ++(4mm,0) |- ([yshift=-1.5mm]cmp.west);
\node[later, text width=26mm, below=7mm of cmp] (later) {later analyses (Table~\ref*{tab:added_analyses})};
\draw[darr] (cmp) -- (later);
\node[panellabel] at ([xshift=-4mm, yshift=2mm]frame.west |- model.north) {a};
\end{tikzpicture}
\\[5mm]
% Generated by tools/pipeline/figures.py from a2_sonnet_20261013/metrics.json (sha256 9d69e490475a885f7216ceb5c49d295b9d7e55662b7092ca6faa1636e46c57f4). Do not edit by hand.
% Panel b of Figure 1 (paper/images/fig_overview.tex inputs this tikzpicture); element precision and recall from
% tools/runs/tables/model_comparison/comparison.json, the registered row bound to this metrics file by digest.
% Palette 2026-11-06: the precision and recall labels in model blue (figModel), as their open marks; the score marks take model blue from the preamble style pooled. Shapes and positions unchanged.
\begin{tikzpicture}
\begin{axis}[thesisfig,
  width=96mm, height=30mm, scale only axis,
  xmin=0, xmax=1, xtick={0,0.5,1}, xticklabels={0,0.5,1}, minor xtick={0.1,0.2,0.3,0.4,0.6,0.7,0.8,0.9}, xlabel={Score},
  ymin=0.4, ymax=3.80, y dir=reverse, rowaxis,
  ytick={1,2,3},
  yticklabels={{Technologies, micro-F1},{Evaluation in shared frames, accuracy},{Frame elements, micro-F1}},
]
\draw[figGreyB, line width=0.8pt] (axis cs:0.5572,3.42) -- (axis cs:0.7591,3.42);
\addplot[precisionmark] coordinates {(0.7591,3.42)};
\addplot[recallmark] coordinates {(0.5572,3.42)};
\addplot[pooled, error bars/.cd, x dir=both, x explicit] coordinates {(0.9350,1) -= (0.0336,0) += (0.0276,0) (0.8741,2) -= (0.0617,0) += (0.0498,0) (0.6427,3) -= (0.0622,0) += (0.0623,0)};
\node[figtext, right=3pt] at (axis cs:0.9626,1) {0.935};
\node[figtext, right=3pt] at (axis cs:0.9239,2) {0.874};
\node[figtext, right=3pt] at (axis cs:0.7050,3) {0.643};
\node[figtext, left=3pt, text=figModel] at (axis cs:0.5572,3.42) {recall 0.56};
\node[figtext, right=3pt, text=figModel] at (axis cs:0.7591,3.42) {precision 0.76};
\end{axis}
\node[panellabel] at ([yshift=2mm]current axis.north -| current axis.outer west) {b};
\end{tikzpicture}

\end{tabular}
% Caption revised 2026-11-06: takeaway restricted to the eligible, prefiltered frame (screening was not evaluated); panel b names its resampling unit; row labels with their metrics are requested in the generator of panel b.
\caption[Evaluation design and headline result]{\textbf{Against the first coding in prefiltered paragraphs, the model agrees closely on technologies but marks fewer frame elements.} \textbf{a},~Evaluation design. \textbf{b},~Registered agreement with the first coding, design-weighted, with article-resampled intervals; open marks: precision and recall, added later. Colours throughout: blue, registered model; black, coder; vermillion, changed or added later; green, model and coder combined; orange, model surplus.}
\label{fig:overview}
\end{figure}


\paragraph{Findings.}
Against the first coding, agreement is high for technology detection in prefiltered paragraphs (micro-F1 0.935) and for evaluations of shared frames (accuracy 0.874; Figure~\ref{fig:overview}b). Micro-F1 pools all labels and takes the harmonic mean of precision (the share of model marks the reference holds) and recall (the share of reference marks the model finds). End-to-end evaluation F1, also counting missed and invented frames, is 0.617. Repeated executions vary in output but give similar scores. Frame elements agree less, at a micro-F1 of 0.643 with no clear language difference, and the model marks fewer of them. Its prevalence, the share of paragraphs where it marks an element, lies below the first coding's in point estimate for 11 of the 15 reported elements, for investment by 18.3 percentage points.

Later analyses, added after the registered result was known (Table~\ref{tab:added_analyses}), set the recall deficit beside the coder's second coding. In point estimate, it narrows but persists on the 40 paragraphs re-coded before the model run. In exploratory comparisons with the full second coding, it nearly closes, but investment's share still differs by $-13.6$ [$-20.4$, $-7.6$] points. Neither coding is established as correct, so the central finding is a recall deficit against one coder, its size depending on the reference reading.

\paragraph{Contributions.}
The contribution is a tested combination of known components, not a new method. One frozen codebook is the common source of prompt, output schema and coding sheets, as proposed for codebooks that coders and models both read \parencite[p.~192]{halterman2026codebookllms}. A registered validation against one coder links label agreement per language and at element level, output variation across executions, the coder's stability over time and the displacement of aggregate shares. Later analyses indicate that the recall deficit's size, and the shares a design-based adjustment yields, depend on which coding is the reference.

\paragraph{Structure of the Thesis.}
Chapter~\ref{sec:background} defines technology frames and elements. Chapter~\ref{sec:related_work} reviews the validation gap. To close it, Chapter~\ref{sec:system_design} designs a pipeline and Chapter~\ref{sec:methodology} its evaluation. Chapter~\ref{sec:results} reports registered and later results, which Chapter~\ref{sec:discussion} interprets, before Chapter~\ref{sec:conclusion} answers RQ1 to RQ4.


\section{Theoretical Background}\label{sec:background}
This chapter defines the \emph{technology frame}: a news frame attached to a named energy technology and measured through concrete frame elements. Section~\ref{subsec:framing} contrasts Entman's definition with the codebook's more permissive elements and traces the nine frames to their sources. Section~\ref{subsec:frame_elements} explains why seventeen presence questions replace holistic frame labels. Section~\ref{subsec:relation} defines the record tying a technology to its frames and their evaluation, shown and scored in Figure~\ref{fig:anatomy}.

\subsection{Framing}\label{subsec:framing}

Following \textcite[p.~52]{entman1993framing}, to frame is \enquote{to select some aspects of a perceived reality and make them more salient in a communicating text}. Such a selection can define a problem, interpret causes, evaluate morally or recommend a treatment \parencite[p.~52]{entman1993framing}. The codebook requires neither such a function nor a complete argument. Most elements are marked when a paragraph refers to their subject for a technology, without a stated mechanism. Only bare naming and purely procedural or factual reporting stay unframed, a boundary that is not sharp (Section~\ref{subsec:limitations_future}).

The evaluation records the direction in which a frame presents a technology. This operational codebook dimension is close to but not identical with Entman's moral evaluation. \textcite[p.~57]{entman1993framing} criticises coders who \enquote{simply tote up} positive and negative messages without measuring salience. Tying each evaluation to a frame answers this only in part: presence is counted, not salience, and only in the text.

The frame dimension operationalises issue framing in Entman's lineage: which aspects of an issue a naturally occurring text makes salient \parencite{card2015mediaframes}. Valence framing, a second tradition, presents the same information positively or negatively, and its attribute variant frames one characteristic of an object, such as beef described by its lean rather than its fat share \parencite[pp.~158--159]{levin1998allframes}. The evaluation dimension does not measure this. Like tone or aspect polarity, it records the direction of what a paragraph asserts within a codebook-defined frame.

Entman lists no issue-specific frames. The topic description supplies five: costs, jobs and the economy, climate and emissions, energy security, and reliability and system integration. Before calibration, these were related to the Policy Frames Codebook's issue-general dimensions \parencite{boydstun2020policyframes,card2015mediaframes} and the energy context labels of \textcite{toetzke2022energydiscourse} (Table~\ref{tab:frame_crosswalk}). Technical feasibility, siting and regulation were added, without a recorded mapping, after the first calibration session, in which 9 of 20 paragraphs fitted no element. A ninth, residual frame follows the Policy Frames Codebook and makes inventory misfit measurable.

\subsection{Frame Elements Instead of Frame Labels}\label{subsec:frame_elements}

Element-based coding is described as less subjective \parencite[p.~3]{ali2025hydrogenframing}, and \textcite[p.~264]{matthes2008contentanalysis} argue that single frame elements are more reliable than abstract, holistic frames. Yet their elements are variables of Entman's four functions, clustered into frames across texts \parencite[pp.~263--264]{matthes2008contentanalysis}. The argument therefore motivates decomposition but establishes no advantage for this study's mostly topical indicators. Two structurally close studies derive frames from binary indicators. \textcite[p.~7]{ali2025hydrogenframing} code sub-variables of Entman's four elements, and \textcite[p.~987]{liu2023strategyframing} mark a frame present when any of its items occurs.

As presence questions, elements let humans and models answer the same questions, on a coding sheet or in a prompt. Per-element agreement also keeps pooled figures from hiding weak elements. \textcite{matthes2008contentanalysis} and \textcite{ali2025hydrogenframing} derive frames from data. This thesis instead fixes the mapping of seventeen elements to nine frames beforehand as its own operationalisation: a frame is present if any of its elements is (Table~\ref{tab:codebook_elements}). The coder's two codings measure temporal stability, but whether derived frames agree with holistic frame judgements remains untested.

\subsection{The Technology--Frame--Evaluation Relation}\label{subsec:relation}\label{subsec:domain_background}

The pipeline and coder produce one \emph{technology record}, or row, per technology a paragraph makes a statement about. The closed inventory comprises renewable energy in general, solar, wind, hydrogen, other renewables, energy storage, and nuclear energy as a low-carbon comparator. Wind siting and, following the topic description, the hydrogen production pathway \parencite{sadatrazavi2024hydrogen} are exported as subtypes but validated only for execution stability, not against the human reference. A comparison target records what the technology is compared with. Elements that \emph{fire}, or are marked present, determine the record's frames by the fixed disjunction. Each frame carries one evaluation (Figure~\ref{fig:anatomy}).

% 2026-11-09, second slimming pass (G-54): caption syntax tightened; no fact removed.
% 2026-11-09: caption shortened (consensus F-08; the author's note: even shorter)
% Figure 2 (Chapter 2, reused for scoring in Chapter 5) -- anatomy of a technology frame and the scoring of a model
% output against it. Drawn by hand in TikZ in the thesis figure style (preamble); restructuring plan v3, Fig 2.
% THE PARAGRAPH IS INVENTED. It is not corpus text and paraphrases none: checked on 2026-11-05 against all 1,047
% paragraphs and headlines of tools/annotation/sheets/ALL_paragraphs.csv (lower-cased word 5-grams): no shared
% 5-word sequence (script and counts only, no corpus text printed). Do not replace it with a corpus paragraph.
% Element names, frames, the disjunctive derivation, the evaluation values and the comparison-target value
% "natural gas" (natural_gas) are those of tools/annotation/codebook/codebook-v1.0.yaml. The coding in panel a is the
% author of the figure's reading of the invented text under the codebook, for illustration only; the model output in
% panel b is hypothetical. The scoring follows tools/pipeline/evaluate.py: element tuples (paragraph x technology x
% element) scored as sets; evaluation score A on the reference (row, frame) pairs the model also predicted, score B on
% (paragraph, technology, frame, value) tuples; comparison target per row, score A on rows present in both;
% displacement on paragraph-level presence indicators (Equation eq:displacement). Arithmetic of panel b: elements TP 2,
% FP 1, FN 2 -> precision 2/3, recall 2/4, micro-F1 4/7 = 0.57; evaluation score A 2/2 = 1.00, score B TP 2, FP 1,
% FN 1 -> 4/6 = 0.67; comparison target: the model gives natural gas (wind) and none (solar), as the coder does in panel a,
% so 2 of 2 matched rows agree (shown in the footer since 2026-11-05, so that the reader can verify it).
% Colour doubles the written outcome only (missed in model blue, invented in orange, as "below" and "above" the
% reference elsewhere); every outcome is also spelled out.
% Palette 2026-11-06: outcome cells moved to the palette names (missed: figModel, invented: figAbove, matched and TP:
% figInk), the "Coder" and "Model" column heads of panel b marked figInk and figModel, the caption names the two
% colours; no value, label text or band fill changed, panel a untouched.
\begin{figure}[htbp]
\centering
\begin{tikzpicture}[
  font=\figfontsmall,
  para/.style={draw=figInk, line width=0.5pt, rectangle, rounded corners=1.5pt, fill=white, align=left,
    inner sep=4pt, text width=146mm, font=\figfont},
  band/.style={fill=figGreyC, rounded corners=2pt},
  el/.style={draw=figInk, line width=0.5pt, rectangle, rounded corners=1.5pt, fill=white, align=left,
    anchor=west, minimum height=5.6mm, text width=63mm, inner sep=2pt},
  fr/.style={draw=figInk, line width=0.8pt, rectangle, rounded corners=1.5pt, fill=white, align=center,
    minimum height=5.6mm, text width=19mm, inner sep=2pt},
  head/.style={anchor=base west, text=figGreyA, font=\figfontsmall},
  tech/.style={anchor=west, font=\figfont\bfseries},
  val/.style={anchor=west},
  arr/.style={-{Latex[length=1.6mm]}, line width=0.6pt},
  cue/.style={},
]
% ---------------------------------------------------------------- panel a: the paragraph
\node[para, anchor=north west] (p) at (0,0) {%
The new offshore wind farm will supply electricity \textbf{more cheaply than gas-fired power stations}\textsuperscript{1} and \textbf{needs no state subsidy}\textsuperscript{2}, the operator said. Residents on the coast fear that \textbf{the turbines will spoil their view of the sea}\textsuperscript{3}. Meanwhile, \textbf{planning law still blocks a solar park}\textsuperscript{4} next to the harbour.};
\node[anchor=south west, inner sep=0pt, font=\figfontsmall\itshape, text=figGreyA] at ([yshift=1.2mm]p.north west)
  {Constructed example, invented for this figure (not corpus text)};
\node[panellabel, anchor=south east] at ([xshift=-1.5mm, yshift=1.2mm]p.north west) {a};

% ---------------------------------------------------------------- panel a: the coding of that paragraph
\begin{scope}[shift={(0,-21mm)}]
  \node[head] at (0,0) {Technology};
  \node[head] at (21mm,0) {Fired elements (codebook names)};
  \node[head] at (95mm,0) {Frame};
  \node[head] at (118mm,0) {Evaluation};
  \node[head, align=left] at (135mm,0) {Comparison\\target};
  \node[head, font=\figfontsmall\itshape] at (95mm,-3.4mm) {any element fires};
  % wind row
  \fill[band] (0,-6mm) rectangle (152mm,-31mm);
  \node[tech] at (1.5mm,-18mm) {Wind};
  \node[el] (e1) at (21mm,-10mm) {\textsuperscript{1}\,Cost or price (F1a)};
  \node[el] (e2) at (21mm,-18mm) {\textsuperscript{2}\,Subsidy, levy or public funding (F1b)};
  \node[el] (e3) at (21mm,-26mm) {\textsuperscript{3}\,Landscape, visual impact and local amenity (F7b)};
  \node[fr] (f1) at (105mm,-14mm) {Costs};
  \node[fr] (f2) at (105mm,-26mm) {Siting};
  \draw[arr] (e1.east) -- ++(4mm,0) |- ([yshift=1mm]f1.west);
  \draw[arr] (e2.east) -- ++(4mm,0) |- ([yshift=-1mm]f1.west);
  \draw[arr] (e3.east) -- (f2.west);
  \node[val] at (118mm,-14mm) {positive};
  \node[val] at (118mm,-26mm) {negative};
  \node[val] at (135mm,-18mm) {natural gas};
  % solar row
  \fill[band] (0,-33mm) rectangle (152mm,-41mm);
  \node[tech] at (1.5mm,-37mm) {Solar};
  \node[el] (e4) at (21mm,-37mm) {\textsuperscript{4}\,Regulation, permitting and legal process (F8a)};
  \node[fr] (f3) at (105mm,-37mm) {Regulation};
  \draw[arr] (e4.east) -- (f3.west);
  \node[val] at (118mm,-37mm) {negative};
  \node[val] at (135mm,-37mm) {none};
\end{scope}

% ---------------------------------------------------------------- panel b: scoring a hypothetical model output
\begin{scope}[shift={(0,-71mm)}]
  \node[panellabel, anchor=south east] at (-1.5mm,1.2mm) {b};
  \node[anchor=south west, inner sep=0pt, font=\figfontsmall\itshape, text=figGreyA] at (0,1.2mm)
    {A hypothetical model output scored against the coding in panel a};
  % elements
  \node[anchor=north west, inner sep=0pt] (tb1) at (0,-1mm) {%
    \renewcommand{\arraystretch}{1.12}\setlength{\tabcolsep}{3pt}%
    \begin{tabular}{@{}l c c l@{}}
    \multicolumn{4}{@{}l}{\textbf{Frame elements}, tuples (row $\times$ element)} \\
    \toprule
    Row $\cdot$ element & \textcolor{figInk}{Coder} & \textcolor{figModel}{Model} & Scored as \\
    \midrule
    wind $\cdot$ cost or price (F1a) & 1 & 1 & \textcolor{figInk}{TP} \\
    wind $\cdot$ subsidy (F1b) & 1 & 0 & \textcolor{figModel}{FN, missed} \\
    wind $\cdot$ landscape (F7b) & 1 & 0 & \textcolor{figModel}{FN, missed} \\
    wind $\cdot$ emissions (F3a) & 0 & 1 & \textcolor{figAbove}{FP, invented} \\
    solar $\cdot$ regulation (F8a) & 1 & 1 & \textcolor{figInk}{TP} \\
    \midrule
    \multicolumn{4}{@{}l}{precision 2/3, recall 2/4, micro-F1 0.57} \\
    \bottomrule
    \end{tabular}};
  % evaluation and comparison target
  \node[anchor=north west, inner sep=0pt] (tb2) at (78mm,-1mm) {%
    \renewcommand{\arraystretch}{1.12}\setlength{\tabcolsep}{3pt}%
    \begin{tabular}{@{}l l l l@{}}
    \multicolumn{4}{@{}l}{\textbf{Evaluation}, one value per row $\times$ frame} \\
    \toprule
    Row $\cdot$ frame & \textcolor{figInk}{Coder} & \textcolor{figModel}{Model} & Scored as \\
    \midrule
    wind $\cdot$ costs & positive & positive & \textcolor{figInk}{matched, agrees} \\
    wind $\cdot$ siting & negative & -- & \textcolor{figModel}{FN, frame missed} \\
    wind $\cdot$ climate & -- & positive & \textcolor{figAbove}{FP, frame invented} \\
    solar $\cdot$ regulation & negative & negative & \textcolor{figInk}{matched, agrees} \\
    \midrule
    \multicolumn{4}{@{}l}{matched-frame accuracy 2/2 = 1.00} \\
    \multicolumn{4}{@{}l}{end-to-end F1 0.67 (TP 2, FP 1, FN 1)} \\
    \multicolumn{4}{@{}l}{comparison target, model: natural gas (wind), none (solar);} \\
    \multicolumn{4}{@{}l}{as coded, so 2 of 2 matched rows agree} \\
    \bottomrule
    \end{tabular}};
  % from tuples to a paragraph prevalence
  \node[anchor=north west, inner sep=0pt] at (0,-50mm) {\textbf{Paragraph indicator}: element present on any row};
  \node[anchor=north west, inner sep=0pt] (tb3) at (0,-54.5mm) {%
    \renewcommand{\arraystretch}{1.12}\setlength{\tabcolsep}{3pt}%
    \begin{tabular}{@{}l *{5}{>{\centering\arraybackslash}p{7mm}}@{}}
    \toprule
    & F1a & F1b & F3a & F7b & F8a \\
    \midrule
    \textcolor{figInk}{Coder}, $y_i$ & 1 & 1 & 0 & 1 & 1 \\
    \textcolor{figModel}{Model}, $\hat{y}_i$ & 1 & \textcolor{figModel}{0} & \textcolor{figAbove}{1} & \textcolor{figModel}{0} & 1 \\
    \bottomrule
    \end{tabular}};
  \node[anchor=north west, text width=72mm, align=left, inner sep=0pt] at (78mm,-54.5mm) {%
    Prevalence: the design-weighted share of sampled paragraphs with indicator 1; displacement: the model's share minus the coder's, in percentage points, as in Equation~\eqref{eq:displacement}. Here, the model's share falls for subsidy and landscape and rises for emissions.};
  \path (0,-75mm); % keeps a little space above the caption
\end{scope}
\end{tikzpicture}
\caption[Anatomy of a technology frame and scoring of a model output]{\textbf{Technology frames derive from fired elements; model output is scored at three levels.} \textbf{a},~An invented paragraph sharing no five-word sequence with the corpus. \textbf{b},~The hypothetical output misses subsidy, but cost or price keeps the costs frame present. TP, FP, FN: true positive, false positive, false negative. Blue: the model and its misses; orange: its inventions.}
\label{fig:anatomy}
\end{figure}


The evaluation is positive, negative, neutral or mixed. Its direction follows the Policy Frames Codebook's litmus test of whether a supporter or opponent would welcome the text \parencite[p.~35]{boydstun2020policyframes}, with the test moved within each frame from a policy to a technology. The result is a coding judgement, not observed audience response or outlet stance. When torn between neutral and a valenced label, the model is told to choose neutral, so unclear and neutral text can share a label. The topic description's fifth value, uncertain, becomes a separate coder-confidence rating, so uncertainty in the text is not recorded.

Each frame of a technology record is a \emph{technology frame}. Unlike the \emph{technological frames} that groups of people hold \parencite{orlikowski1994frames}, it is a property of a text, and a record can carry several or none. Technology, frame and evaluation correspond to the entity, the aspect discussed and its polarity in aspect-category sentiment analysis \parencite{pontiki2014semeval,pontiki2016semeval}. The value \emph{mixed} follows the 2014 Semantic Evaluation (SemEval) shared task's \emph{conflict} label, which forced no dominant polarity \parencite{pontiki2014semeval}.


\section{Related Work}\label{sec:related_work}
The pipeline draws on four lines of research (Sections~\ref{subsec:stream_framing} to~\ref{subsec:stream_validation}, Table~\ref{tab:literature_short}): energy-framing studies for precedents of its frames, computational framing analysis for its scale, language-model annotation for its coding procedure and validation research for its standard. Table~\ref{tab:literature} in Appendix~\ref{app:literature} compares the studies in full, by material, method, finding and what was taken. Section~\ref{subsec:research_gap} names the gap they leave.

% 2026-11-11: rewritten for the fourth review's plan (M16); earlier versions and number sources in git (state d5a2bfb).
% 2026-11-13 (final optimisation pass, plan item TAB1): two columns; the column "Finding relied on" dropped (findings by study remain in Table tab:literature, Appendix A.1); the same citation keys moved into the design column, now headed "Taken into this design"; caption reworded to match; table note unchanged.
% 2026-11-16: table note removed (author comment C14); the preprint caveat remains in Table tab:literature; caption semicolon replaced.
\begin{table}[htbp]
\centering
\scriptsize
\setlength{\tabcolsep}{4pt}
\renewcommand{\arraystretch}{1.15}
\begin{threeparttable}
\caption[Studies that shaped the design (abridged)]{Literature streams and what the design takes from each (findings by study: Table~\ref{tab:literature}, Appendix~\ref{app:literature}).}
\label{tab:literature_short}
\begin{tabular}{@{}>{\raggedright\arraybackslash}p{38mm}>{\raggedright\arraybackslash}p{112mm}@{}}
\toprule
\textbf{Stream} & \textbf{Taken into this design} \\
\midrule
Media framing of energy technologies & evaluation separate from frame; per-element reliability; fixed element-to-frame mapping, not factor analysis \parencite{stephens2009windmedia,ali2025hydrogenframing} \\
\addlinespace
Computational analysis of energy discourse & paragraph as coding unit; keyword matching to technologies \parencite{toetzke2022energydiscourse} \\
\addlinespace
Large language models as annotators & full codebook definitions as prompt; one joint call per paragraph; agreement beyond raw percentages; configuration fixed in advance; human reference; one codebook for coder and model \parencite{gilardi2023chatgpt,lupo2024humanlevelcoding,toetzke2023llmclimatetech,baumann2025llmhacking,halterman2026codebookllms} \\
\addlinespace
Validation and aggregates & reference's temporal stability; sample-size limit; registered displacement; design-weighted probability sample; declared exploratory adjustment \parencite{song2020validations,linder2026multicalibrated,egami2023dsl} \\
\bottomrule
\end{tabular}
\end{threeparttable}
\end{table}


\subsection{Media Framing of Energy Technologies}\label{subsec:stream_framing}

Manual studies show that energy frames and their valence vary between places \parencites[p.~6]{rochyadireetz2019framing}[pp.~177, 179]{stephens2009windmedia} and that coverage need not track deployment \parencite[p.~179]{stephens2009windmedia}. \textcite[pp.~173--174]{stephens2009windmedia} code risks and benefits within each frame. The closest published predecessor, the hydrogen study by \textcite[p.~7]{ali2025hydrogenframing}, instead codes evaluation on a separate scale, as the codebook does. The three frames added in calibration recur in this coverage \parencites[pp.~6, 9]{ali2025hydrogenframing}[pp.~177--181]{stephens2009windmedia}. Each study coded at most a few thousand articles by hand.

\subsection{Computational Framing Analysis across Languages}\label{subsec:stream_computational}\label{subsec:stream_multilingual}

Computational studies can scale further, but those reviewed mostly translate inputs or use one model per language. \textcite{toetzke2022energydiscourse}, the only energy-discourse study among them, plan to translate about 40 million paragraphs into English. \textcite{jumle2024findingframes} fine-tune on machine-translated paragraphs. \textcite{guo2023openframing} train separate English and German headline classifiers without a bilingual test set. News-frame benchmarks cover six \parencite{piskorski2023multifaceted} and nine languages, the latter scored per language on original-language test sets \parencite{piskorski2023semeval}. Headline- and sentence-framing studies compare translation with cross-lingual transfer in multilingual encoders \parencite{akyurek2020multilingualframing,akter2024scalingframing}. All score frame classification.

Translation can change labels, as it did for human quality labels of arguments \parencite[p.~310]{toledoronen2020multilingual}. Designs varying only the language of fixed content find that model behaviour differs by language \parencite{xing2024crosslingualinconsistency,hasan2026englishrewrites,xuan2025mmluprox,liang2025multilingualbehaviour}. With one untranslated prompt in eleven countries, \textcite[p.~1189]{toernberg2025outperform} finds lower accuracy outside the United States. Most evidence on language-model annotation comes from English text and the United States \parencite[p.~1183]{toernberg2025outperform}. Agreement should therefore be measured on source text, per language.

\subsection{Large Language Models as Annotators}\label{subsec:stream_llm}

Across 27 replicated annotation tasks, agreement varied widely, so every application needs its own validation \parencite[p.~1472]{pangakis2025humansintheloop}. In the closest methodological analogue, exhaustive label definitions improved coding, and adding labelled examples worked best \parencite[p.~9]{lupo2024humanlevelcoding}. From the chair's extraction pipeline, this thesis adopts temperature-zero decoding, structured output, an explicit abstention instruction and a stage-by-stage attrition account \parencite{forster2026corporatesustainability}.

Several failure modes limit model coding. Thirteen models rarely beat fine-tuned classifiers, and many errors in four analysed tasks favoured the most generic label \parencite[pp.~259--260]{ziems2024css}. Open-weight models show label bias \parencite[p.~6784]{reif2024labelbias}. List-generating models tended to concentrate probability on one label per step \parencite[p.~2472]{ma2025multilabel}. Minor prompt changes can shift label distributions \parencite{atreja2025whatsinaprompt}, and configuration choices yielded incorrect conclusions for 31 to 50\% of tested hypotheses \parencite{baumann2025llmhacking}. Repeated runs disagree, less so at lower temperatures \parencite{gilardi2023chatgpt}. Hence, the configuration is fixed before evaluation and stability is measured across executions.

\subsection{Validating Automated Content Analysis}\label{subsec:stream_validation}

A validation is only as good as its human labels and sample. In simulation, less reliable labels distorted performance estimates, yet only 14 of 37 studies validating against human coding reported intercoder reliability \parencite[pp.~554, 560]{song2020validations}. For a simulated 130{,}000-article corpus, \textcite[p.~564]{song2020validations} advise over 1{,}300 validation items with intercoder reliability of at least 0.7. Small test sets detect only large differences \parencite{card2020littlepower}. Capturing rare labels requires sufficient sample size and heterogeneity \parencite[p.~106]{feuerriegel2025nlptextdata}.

Classifying each document and counting the labels can bias category proportions even when individual labels are accurate \parencite[p.~229]{hopkins2010method}, so this thesis evaluates prevalences separately from label agreement. Even at per-language areas under the receiver operating characteristic curve (AUC) of 0.983 to 0.994, a model's prevalences from standard estimators stayed biased \parencite[pp.~7--8]{linder2026multicalibrated}. \textcite{egami2023dsl} combine model labels with a human-coded probability sample in a design-based estimator. Such an estimator adjusts shares towards the human coding without testing it.

\subsection{Research Gap and Contribution}\label{subsec:research_gap}

The closest studies each leave a gap. \textcite{ali2025hydrogenframing} code elements by hand, and \textcite{toetzke2022energydiscourse} plan to classify energy discourse in English translation. \textcite{lupo2024humanlevelcoding} validate only Swedish model coding, and other non-English annotation studies cover one language or pair \parencite{nemkova2025annotationskills,hohenwalde2025chatgptactors}. Testing only 7--12-billion-parameter open-weight models, \textcite[p.~188]{halterman2026codebookllms} find limited zero-shot adherence to three political-science codebooks. \textcite{linder2026multicalibrated} measure prevalence bias in four languages, but for one topic label. No reviewed study validates one language-model instrument at element level, per language and on aggregate prevalences.

Codebook prompting, per-language scoring and prevalence correction are each established. For one proprietary model on German and English paragraphs as written, this thesis combines them at element level, beside the coder's own retest. Two constraints shape the evaluation protocol: a single coder, and 120 paragraphs with rare labels, fewer than \textcite[p.~564]{song2020validations} advise.


\section{System Design}\label{sec:system_design}
RQ1 asks why each part of the pipeline is the way it is. Four decisions answer it, each guarding the validation against one threat. One frozen codebook supplies the definitions for the model's request and the coder's sheets, though the instructions each receives differ (Section~\ref{subsec:codebook}). A split registry keeps calibration and development paragraphs out of every evaluation figure. One stateless call codes each paragraph untranslated, so agreement is measured on source text. A pre-run gate and an append-only journal tie every prediction to its configuration. Sections~\ref{subsec:architecture},~\ref{subsec:codebook} and~\ref{subsec:extraction_call} cover the pipeline (Figure~\ref{fig:architecture}), the codebook (Table~\ref{tab:codebook_elements}) and the call.

\subsection{Architecture Overview}\label{subsec:architecture}

A split registry (Figure~\ref{fig:architecture}, centre) gives every corpus paragraph a permanent identity and a fixed role, such as calibration or evaluation.

% 2026-11-09, second slimming pass (G-52): caption tightened again; no key or fact removed.
% 2026-11-09: caption shortened; the evaluator is the solid green box, not a dashed frame (consensus F-09)
% Pipeline architecture (Section 4.1, moved from the appendix 2026-11-05); it implements the evaluation design of Figure 1a.
% Drawn by hand in TikZ in the thesis figure style (preamble): sans-serif text, neutral fills, no hatching; the
% stores are shaded light grey. Restyled 2026-11-05: labels no longer crossed by arrows or by the dashed frames, the
% paragraph arrows enter the runner and the coder from the side. Every box names a module that exists in tools/; the dashed boxes are the
% two halves that share the codebook as their single source. No measured quantity appears in this figure.
% Palette 2026-11-06: the model half frame and its label take figModel, the human half frame and its label figInk (both stay dashed and labelled, so they read in greyscale by position and label); the call journal, the model-side store, is filled figModelLight (this supersedes "shaded light grey" above for that store), the reference labels store keeps figGreyC; the evaluator is outlined figAdjusted at 1pt (it joins both halves; its weight and its position where both arrows meet carry the cue); the arrows from derive and from labels into it take figModel and figInk. New styles derive from the shared ones; no geometry, text or value changed.
\begin{figure}[htbp]
\centering
\begin{tikzpicture}[
  font=\figfontsmall,
  node distance=6mm and 9mm,
  box/.style={draw=figInk, line width=0.5pt, rectangle, rounded corners=1.5pt, minimum width=31mm, minimum height=9mm,
    align=center, text width=31mm, inner sep=2pt, fill=white,
    execute at begin node={\hyphenpenalty=10000\exhyphenpenalty=10000}},
  store/.style={box, fill=figGreyC, minimum height=11mm},
  src/.style={box, line width=1pt, fill=figGreyC, text width=32mm},
  half/.style={dashed, line width=0.5pt, rounded corners=2pt, inner sep=4mm},
  halflabel/.style={font=\figfontsmall\itshape},
  halfmodel/.style={half, draw=figModel},
  halfhuman/.style={half, draw=figInk},
  halflabelmodel/.style={halflabel, text=figModel},
  halflabelhuman/.style={halflabel, text=figInk},
  storemodel/.style={store, fill=figModelLight},
  evalbox/.style={box, draw=figAdjusted, line width=1pt},
  arr/.style={-{Latex[length=1.8mm]}, line width=0.6pt},
  arrmodel/.style={arr, draw=figModel},
  arrhuman/.style={arr, draw=figInk},
  lbl/.style={font=\figfontsmall, text=figGreyA, inner sep=1pt},
]
% --- single source
\node[src] (codebook) {\textbf{Frozen codebook}\\ one versioned file, hash-pinned};

% --- model half (left column)
\node[box, below left=16mm and 24mm of codebook] (render) {\textbf{Renderer}\\ system prompt and tool schema};
\node[box, below=of render] (gate) {\textbf{Pre-run gate}\\ contract, forced tool use, cache read-back};
% 2026-11-06 (review G3): the runner box names the three executions instead of the symbol k.
\node[box, below=of gate] (runner) {\textbf{Extraction runner}\\ one stateless call per paragraph, three executions};
\node[storemodel, below=of runner] (journal) {\textbf{Append-only call journal}\\ private, raw responses};
\node[box, below=of journal] (derive) {\textbf{Derivation}\\ terminal outcomes, predictions};

% --- human half (right column)
\node[box, below right=16mm and 24mm of codebook] (sheets) {\textbf{Sheet generator}\\ two-pass coding sheets};
\node[box, below=of sheets] (coder) {\textbf{Human coder}\\ pass 1: technologies; pass 2: elements, evaluation, target};
\node[box, below=of coder] (finalize) {\textbf{Finalisation}\\ strip text, validate against the codebook};
\node[store, below=of finalize] (labels) {\textbf{Reference labels}\\ identifiers and labels only};

% --- shared centre column: corpus, registry, sample, evaluator
\node[box, below=10mm of codebook] (corpus) {\textbf{Corpus}\\ articles, segmented into paragraphs};
\node[box, below=of corpus] (registry) {\textbf{Split registry}\\ identities, hash-chained assignment log, frozen frame};
\node[box, below=of registry] (sample) {\textbf{Evaluation sample}\\ stratified draw, design weights};
\node[evalbox, anchor=north] (evaluate) at ([yshift=-12mm]codebook.south |- derive.south) {\textbf{Evaluator}\\ agreement, stability, displacement};

% --- frames around the halves, labelled below so that no arrow crosses the label
\begin{scope}[on background layer]
  \node[halfmodel, fit=(render)(derive)] (mh) {};
  \node[halfhuman, fit=(sheets)(labels)] (hh) {};
\end{scope}
\node[halflabelmodel, below=1mm of mh.south] {model half};
\node[halflabelhuman, below=1mm of hh.south] {human half};

% --- edges
\draw[arr] (codebook.west) -| (render.north);
\draw[arr] (codebook.east) -| (sheets.north);
\draw[arr] (render) -- (gate);
\draw[arr] (gate) -- (runner);
\draw[arr] (runner) -- (journal);
\draw[arr] (journal) -- (derive);
\draw[arr] (sheets) -- (coder);
\draw[arr] (coder) -- (finalize);
\draw[arr] (finalize) -- (labels);
\draw[arr] (corpus) -- (registry);
\draw[arr] (registry) -- (sample);
% the paragraph labels sit above the outer segment beside the sample box, clear of the dashed half frames
\draw[arr] (sample.west) -- ++(-17mm,0) node[lbl, midway, above] {paragraphs} |- (runner.east);
\draw[arr] (sample.east) -- ++(17mm,0) node[lbl, midway, above] {paragraphs} |- (coder.west);
\draw[arr] (sample.south) -- (evaluate.north) node[lbl, pos=0.4, left=1mm, align=right] {design weights,\\ manifest digests};
\draw[arrmodel] (derive.east) -| ([xshift=-8mm]evaluate.north);
\draw[arrhuman] (labels.west) -| ([xshift=8mm]evaluate.north);
\end{tikzpicture}
\caption[Pipeline architecture]{\textbf{One frozen codebook supplies the model's prompt and the coder's sheets; model and coder labels are compared on a design-weighted sample.}}
\label{fig:architecture}
\end{figure}


One builder assembles every model request. A pre-run gate refuses any configuration whose codebook differs from its recorded digest, a cryptographic fingerprint, so no edited codebook reaches the model unnoticed. Every run records the digests of its prompt, tool configuration, codebook and sample manifest (Table~\ref{tab:electronic}). Raw responses enter an append-only journal before derivation, so a crash cannot lose a paid call.

The coder's labels are checked against the codebook and stored without paragraph text as the scoring reference. As a precaution, complete paragraph texts, raw responses and coding sheets stay outside the repository, because no basis for storing the hand-collected texts was clarified (Table~\ref{tab:topic_requirements}).

\subsection{The Codebook as Research Instrument}\label{subsec:codebook}

The codebook, a single versioned file, defines technologies, elements with inclusion and exclusion rules, and derived frames (Table~\ref{tab:codebook_elements}). It was frozen as version~1.0 after two calibration sessions.

The prompt, tool schema and coding-sheet entry fields are rendered from this file, and the prompt carries the full definitions motivated in Section~\ref{subsec:stream_llm} (Appendix~\ref{app:instrument}). Of the codebook's eleven exclusion rules, the renderer passes four to the model as written, rewords four and omits three, as decided on 2 September 2026 (Table~\ref{tab:rule_concordance}). Which documents the coder consulted was not logged.

% Generated by tools/pipeline/figures.py codebook from codebook-v1.0.yaml.
% Do not edit by hand.
\begin{table}[htbp]
\centering
\small
\begin{threeparttable}
\caption[Frame elements and frames of codebook v1.0]{Frozen codebook (v1.0): 17 frame elements and 9 derived frames. A frame fires when any of its elements fires; coder and model mark elements only.}
\label{tab:codebook_elements}
\begin{tabular}{@{}ll>{\raggedright\arraybackslash}p{80mm}l@{}}
\toprule
\textbf{Frame} & \textbf{Element} & \textbf{Label} & \textbf{Source} \\
\midrule
Costs & F1a & Cost or price & Topic description \\
 & F1b & Subsidy, levy or public funding &  \\
\addlinespace
Jobs & F2a & Employment & Topic description \\
 & F2b & Investment, industry or competitiveness (narrowed in v1.0) &  \\
 & F2c & Deployment growth, records and market share (added in v1.0) &  \\
\addlinespace
Climate & F3a & Emissions or climate effect & Topic description \\
 & F3b & Transition or phase-out pathway &  \\
\addlinespace
Security & F4a & Import dependence or independence & Topic description \\
 & F4b & Security of supply or geopolitics &  \\
\addlinespace
Reliability & F5a & Intermittency and dispatchability & Topic description \\
 & F5b & Grid and system integration &  \\
\addlinespace
Feasibility & F6a & Technical feasibility and maturity & Calibration \\
 & F6b & Technical performance and potential &  \\
\addlinespace
Siting & F7a & Land and area requirement & Calibration \\
 & F7b & Landscape, visual impact and local amenity &  \\
\addlinespace
Regulation & F8a & Regulation, permitting and legal process & Calibration \\
\addlinespace
Other & F9 & Other frame (requires a written note) & Residual \\
\bottomrule
\end{tabular}
\begin{tablenotes}\footnotesize
\item A named but unframed technology takes \emph{no frame for this technology}. Codes such as F1a are identifiers, not F1 scores.
\end{tablenotes}
\end{threeparttable}
\end{table}


The source is shared, but history and context are not. The prompt omits the calibration history and category frequencies to avoid installing a label prior, whereas the coder ran the calibration and cannot unlearn its outcomes. Both instruments restrict coding to the paragraph, yet the coder saw more context (Section~\ref{subsec:limitations_future}).

\subsection{The Extraction Call}\label{subsec:extraction_call}

One stateless, history-free call returns all technology records of a paragraph. It carries only the system prompt, tool schema and paragraph, and forces one tool whose schema prescribes the answer format. One sentence may contrast two technologies, and in \textcite[p.~9]{lupo2024humanlevelcoding}, joint coding performed like separate runs at almost a third of their cost. Adding labelled examples worked best in that study but would have revised the frozen instrument.

English instructions accompany paragraphs in their original language, as in \textcite[pp.~1184, 1186]{toernberg2025outperform}. The provider exposes no seed, so repeated byte-identical requests measure run-to-run variation (Section~\ref{subsec:metrics}).

Responses are validated in a fixed order: first the stop reason, so that a truncated answer counts as a failure, not an abstention. The schema follows, then relational rules a schema cannot express, such as one evaluation per frame present.


\section{Methodology}\label{sec:methodology}
A pipeline's labels are useful only if they agree with human coding on paragraphs it was not tuned on. One primary evaluation tests this on a stratified probability sample of 120 paragraphs, 60 per language, labelled by one coder and coded in three identical model executions. Sections~\ref{subsec:research_design} to~\ref{subsec:added_methods} cover the design, the corpus and draw, the human reference, the metrics, the model setup and the later analyses. Chapter~\ref{sec:system_design} answers RQ1, and Section~\ref{subsec:metrics} defines the estimands for RQ2 to RQ4: element agreement per language, change rates per dimension and language, and prevalence displacement qualified by class support.

\subsection{Research Design}\label{subsec:research_design}

The pipeline codes zero-shot: a language model receives the codebook but no complete labelled example \parencite[p.~104]{feuerriegel2025nlptextdata}. Three reasons favour this over a trained classifier. First, one model codes German and English as written, unlike most framing studies reviewed in Section~\ref{subsec:stream_computational}. Second, each row carries seventeen binary elements, many of them rare, whereas feature- and embedding-based classifiers typically need over a thousand labelled examples \parencite[p.~104]{feuerriegel2025nlptextdata}. Avoiding such training reserves the single coder's labels for validation \parencite[p.~259]{ziems2024css}. Third, coder and model share one codebook's definitions (Section~\ref{subsec:codebook}).

Dictionaries count word frequencies and cannot capture context \parencite[p.~103]{feuerriegel2025nlptextdata}, such as which evaluation belongs to which technology and frame. The case for the model therefore rests not on an established technology-naming advantage (Section~\ref{subsec:agreement}) but on joint relational extraction and evaluation. Zero-shot coding has known failure modes (Section~\ref{subsec:stream_llm}). The validation measures agreement and displacement but identifies none as a cause.

The primary evaluation followed an internal pre-analysis specification recorded in the repository before evaluation coding and prediction. \emph{Registered} denotes this record, not a public registry or a hypothesis pre-registration (Appendix~\ref{app:protocol}). The supervisor endorsed the frozen design and the hand-built corpus on 10 September 2026. Calibration shaped the codebook, and development informed prompt decisions fixed before the draw. The reference and the registered 40-paragraph re-code were committed before the model run (Section~\ref{subsec:annotation}).

The \emph{registered run} comprises three pipeline executions on the evaluation sample under an identical request. The \emph{first execution} is scored for agreement and displacement whatever its outcome, as the pipeline runs once per paragraph in use. A majority vote would be another, costlier instrument and is not computed. All three executions enter the stability measure. Neither agreement nor displacement establishes the human reference's validity.

\subsection{A Purposive, Hand-Built Corpus and a Probability Sample}\label{subsec:data_collection}

The corpus was hand-collected from online news editions (Table~\ref{tab:corpus_facts}). Factiva, the topic description's preferred source, was considered but not used, and its licence was later reported to bar automated queries. The alternative, the Global Database of Events, Language, and Tone (GDELT), may lack full text (Appendix~\ref{app:topic_requirements}).

Purposive selection targeted six cases expected to strain the pipeline (Table~\ref{tab:corpus_facts}). Three are documented difficulties of target-dependent sentiment analysis. Coverage may evaluate two targets in opposite directions \parencites[p.~1669]{hamborg2021newsmtsc}[p.~3594]{xing2020arts} or quote other speakers' evaluations \parencite[p.~610]{bastan2020authorsentiment}. It may also be balanced or purely factual, with evaluation at most implicit \parencite[p.~1664]{hamborg2021newsmtsc}. Article-level case tags overlap, and no result is reported by case. The corpus is not a representative sample of energy coverage but an instrument test bed.

Articles are split into paragraphs of at least 15 words, which drops fragments such as sub-heads. A recall-oriented keyword prefilter over the seven technologies marks \emph{candidates}. A candidate is \emph{eligible} unless it belongs to the calibration or development split, has a transcription gap or comes from the one article the prompt quotes (Figure~\ref{fig:sample_flow}). Every selection is drawn from this \emph{eligible frame}, a sampling frame in the survey sense, not a news frame (Chapter~\ref{sec:background}).

% 2026-11-09, second slimming pass (G-55): legend wording tightened; no count or key removed.
% 2026-11-09: caption tightened (consensus N-01).
% Figure (Chapter 5, Data Collection) -- from the collected corpus to the samples. Drawn by hand in TikZ in the thesis
% figure style (preamble): sans-serif text, no hatching, colours from the preamble palette. Every number was checked on 2026-11-05
% against tools/annotation/split-assignment.json (articles with and without prose, assignment log, counts), the
% candidate and gap flags of build_sheets.collect(), and the manifests in tools/annotation/labels/:
% A2_frame_manifest.json (618; 225 de / 393 en; 67 articles; 18 strata; exemplar article excluded),
% A2_sample_manifest.json (120; 60 / 60; drawn 2026-09-06; 40 re-code flags, 20 / 20), A3_sample_manifest.json
% (55; 27 de / 28 en; ten articles; created 2026-10-14) and A2_frame_completion_manifest.json (443; created
% 2026-10-19). The three selections are disjoint and their union is the frame (120 + 55 + 443 = 618).
% Lower part: the bar widths are proportional to paragraph counts, 150 mm for the 618 paragraphs of the frame
% (0.2427 mm per paragraph): 120 -> 29.13 mm, 40 -> 9.71 mm, 55 -> 13.35 mm, 443 -> 107.52 mm.
% Palette 2026-11-06: the lower bars carry the preamble palette, one meaning per colour. Evaluation sample: thick figInk
% outline, white fill (was figGreyB fill); re-code: figInk fill (was figGreyA); fresh census: figChanged outline 1pt,
% white fill (was figInk 0.5pt); frame completion: figModel outline 0.5pt, figModelLight fill (was figGreyA / figGreyC).
% The remaining-frame bar, the guides and the upper flow stay black and grey. Greyscale cues: outline weight (thick =
% evaluation sample), fill lightness (black = re-code, white = coder-labelled, light = model only) and the bar's level.
\begin{figure}[htbp]
\centering
\begin{tikzpicture}[
  x=1mm, y=1mm,
  font=\figfontsmall,
  sfstep/.style={draw=figInk, line width=0.5pt, rectangle, rounded corners=1.5pt, fill=white, align=center,
    text width=70mm, minimum height=6.5mm, inner sep=1.5pt,
    execute at begin node={\hyphenpenalty=10000\exhyphenpenalty=10000}},
  sfframe/.style={sfstep, line width=1pt},
  sfsplit/.style={sfstep, draw=figGreyA, fill=figGreyC, text=figInk, text width=68mm, align=left, inner sep=2pt},
  sfexit/.style={anchor=west, align=left, text width=66mm, text=figGreyA, inner sep=1pt},
  arr/.style={-{Latex[length=1.8mm]}, line width=0.6pt},
  xarr/.style={-{Latex[length=1.5mm]}, line width=0.5pt, figGreyA},
  sfseg/.style={line width=0.5pt, draw=figInk},
  sfnote/.style={anchor=north west, align=left, inner sep=0pt,
    execute at begin node={\hyphenpenalty=10000\exhyphenpenalty=10000}},
  sflead/.style={line width=0.4pt, figGreyA},
]
% ------------------------------------------------------------------ upper part: corpus to eligible frame
\node[sfstep] (collected) at (40,0) {\textbf{95} articles collected\\ from 25 outlets in six countries};
\node[sfstep] (corpus) at (40,-14) {\textbf{1{,}047} paragraphs in \textbf{74} articles with body prose\\ 361 German, 686 English};
\node[sfstep] (reserve) at (40,-29) {Reserve split: \textbf{983} paragraphs in 69 articles};
\node[sfstep] (cand) at (40,-40) {\textbf{641} prefilter candidates in the reserve};
\node[sfstep] (gapfree) at (40,-51) {\textbf{623} candidates without a transcription gap};
\node[sfframe] (frame) at (40,-63) {Eligible frame: \textbf{618} paragraphs in 67 articles\\ 18 strata; 225 German, 393 English};

\draw[arr] (collected) -- (corpus);
\draw[arr] (corpus) -- (reserve) node[midway, left=1.5mm, font=\figfontsmall\itshape, text=figGreyA] {split registry};
\draw[arr] (reserve) -- (cand);
\draw[arr] (cand) -- (gapfree);
\draw[arr] (gapfree) -- (frame);

% exits and the other registry splits, to the right of the main column
\draw[xarr] (40,-7) -- (80,-7) node[sfexit] {21 articles without body prose};
\draw[xarr] (40,-21.75) -- (80,-21.75) node[sfsplit, anchor=west] {Other registry splits, used before the draw:\\
  calibration: 32 paragraphs in sessions of 20 and 12\\ development: 32 paragraphs};
\draw[xarr] (40,-34.5) -- (80,-34.5) node[sfexit] {342 not marked by the keyword prefilter};
\draw[xarr] (40,-45.5) -- (80,-45.5) node[sfexit] {18 with a transcription gap};
\draw[xarr] (40,-56.5) -- (80,-56.5) node[sfexit] {5 from the article quoted in the prompt};

% ------------------------------------------------------------------ lower part: the frame and its three selections
% level 1: the registered evaluation sample and the remaining frame
% guide lines from the frame box to the ends of the bar: the bar is the frame, partitioned
\draw[sflead, figGreyB] (frame.south west) -- (0,-74) (frame.south east) -- (150,-74);
\draw[sfseg, line width=1pt, fill=white] (0,-74) rectangle (29.13,-82);
\node[align=center, inner sep=0pt] at (14.565,-78) {Evaluation\\ sample \textbf{120}};
\draw[sfseg, draw=figGreyA, fill=white] (29.13,-74) rectangle (150,-82);
\node[text=figGreyA] at (89.565,-78) {Remaining frame: \textbf{498} paragraphs, selections from it made later};
% level 2: the re-code nested in the evaluation sample; the fresh census and the completion from the remaining frame
\draw[sflead, figGreyB] (0,-82) -- (0,-84) (29.13,-82) -- (29.13,-84) (150,-82) -- (150,-84);
\draw[sfseg, fill=figInk] (0,-84) rectangle (9.71,-92);
\node[text=white] at (4.855,-88) {\textbf{40}};
\node[anchor=west, inner sep=0pt, text=figInk] at (10.9,-88) {re-code};
\draw[sfseg, draw=figChanged, line width=1pt, fill=white] (29.13,-84) rectangle (42.48,-92);
\node at (35.805,-88) {\textbf{55}};
\draw[sfseg, draw=figModel, line width=0.5pt, fill=figModelLight] (42.48,-84) rectangle (150,-92);
\node at (96.24,-88) {Frame completion: \textbf{443} paragraphs};

% annotations: one column per selection, joined to its bar by a leader
\node[sfnote, text width=44mm] (a2) at (0,-96) {\textbf{Evaluation sample} (registered):\\ 60 German, 60 English;
  stratified probability draw on 6~Sep~2026, coded and scored.\\ \textbf{Re-code}: 40 of the 120 (20 per language),
  flagged at the draw and coded again, blind, before the model run. The other 80 formed the supplementary re-code after the results.};
\draw[sflead] (35.805,-92) -- (35.805,-94) -| (54,-96);
\node[sfnote, text width=44mm] (a3) at (54,-96) {\textbf{Fresh census} (added): all 55 paragraphs, 27 German and
  28 English, of ten articles with no documented prior use in the study: none of their paragraphs had been used in calibration, development or the evaluation sample;
  selected on 14~Oct~2026; coded; unweighted.};
\draw[sflead] (110,-92) -- (110,-96);
\node[sfnote, text width=40mm] (fc) at (110,-96) {\textbf{Frame completion} (added): all remaining frame paragraphs (443),
  listed on 19~Oct~2026; model predictions only, never coded; completes the frame for the prevalence adjustment.};
\end{tikzpicture}
\caption[From the corpus to the samples]{\textbf{The evaluation sample was drawn first from the frozen sampling frame of 618 paragraphs; the fresh census and frame completion followed from the remaining 498.} Bars: disjoint selections, count-proportional widths. Black: coder-labelled. Vermillion: later fresh census, same coder. Blue: model predictions only.}
\label{fig:sample_flow}
\end{figure}


The \emph{evaluation sample} was drawn at random, once, from the frozen eligible frame in 18 strata, one per outlet within a language (Table~\ref{tab:corpus}). One coder's resources for two passes per paragraph set its size at 120. Drawing 60 per language allows per-language figures. A proportional draw would have favoured English. Within a language, allocation is proportional to stratum size, capped at 24 per outlet. Paragraph $i$ of stratum $h$ has a known selection probability, as design-based estimation requires \parencite{egami2023dsl}. Its inverse is the design weight \parencite[p.~669]{horvitz1952generalization}
\begin{equation}
\pi_h = \frac{n_h}{N_h}, \qquad w_i = \frac{1}{\pi_{h(i)}},
\label{eq:weights}
\end{equation}
where $N_h$ is the stratum size in the eligible frame and $n_h$ the allocated draw. The sample covers 56 articles.

% Generated by tools/annotation/descriptives.py from the corpus manifest, the split
% registry, the frozen A2 frame manifest and the registered A2 sample manifest. Do not edit by hand.
% Summary by language since 2026-11-09 (slimming debate consensus F-03); outlet rows: tab:corpus_outlets.
\begin{table}[htbp]
\centering
\footnotesize
\setlength{\tabcolsep}{3pt}
\begin{threeparttable}
\caption[Corpus and evaluation allocation by language]{Corpus, eligible frame and evaluation allocation by language (by outlet: Table~\ref{tab:corpus_outlets}).}
\label{tab:corpus}
\begin{tabular}{@{}lrrrrrrr@{}}
\toprule
\textbf{Language} & \textbf{Outlets} & \textbf{Articles} & \textbf{Paragraphs} & \textbf{Candidates} & \textbf{Eligible} & \textbf{Drawn} & \textbf{Weight range} \\
\midrule
German & 13 & 40 & 361 & 265 & 225 & 60 & 3.43--5.00 \\
English & 6 & 34 & 686 & 431 & 393 & 60 & 5.50--7.67 \\
Total & 19 & 74 & 1{,}047 & 696 & 618 & 120 & \\
\bottomrule
\end{tabular}
\begin{tablenotes}\footnotesize
\item Only articles with body prose, and their outlets, are counted. Candidates pass the keyword prefilter; eligible candidates lie in the reserve and survive the exclusions in Figure~\ref{fig:sample_flow}. Clean Energy Wire forms no stratum.
\end{tablenotes}
\end{threeparttable}
\end{table}


\subsection{A Single Coder as the Human Reference}\label{subsec:annotation}

The author, a native German speaker who reads English at academic level, produced all human labels alone (Table~\ref{tab:sessions}). Under the paragraph-only rule (Section~\ref{subsec:codebook}), the first pass marks technologies and the second marks each technology's \emph{frame elements}, evaluations and \emph{comparison target}. These labels are the \emph{first coding}, the registered reference. Only the coder records speaker attribution (Table~\ref{tab:checks_4p}) and a confidence rating. Neither enters a score. The coder never saw the keyword prefilter, except in a discarded first-pass attempt whose interface suggested keyword-derived technologies.

The general category covers renewables or the energy transition as a bundle. The reference reads it per statement, so a bundle term beside named technologies adds a general row. The prompt's wording also admits a residual reading, general only where no single technology carries a statement. Its figures differ (Table~\ref{tab:disclosures}). A later check without general rows leaves the element score near the registered one (Table~\ref{tab:checks_4o}).

The registered blind re-code (wave 1), a random 20 paragraphs per language, preceded the model run (Figure~\ref{fig:sample_flow}). It measures not agreement between coders but the \emph{temporal stability of the coder} \parencite[p.~215]{krippendorff2004contentanalysis}. Its washout was shortened from 14 planned to at least five realised days (Table~\ref{tab:deviations}).

The \emph{coder's~$\alpha$} is unweighted nominal Krippendorff's $\alpha = 1 - D_o/D_e$, with observed and expected disagreement $D_o$ and $D_e$ \parencite[pp.~1, 3]{krippendorff2011computing}. It is computed on shared units, for elements on the technology rows both codings created (Tables~\ref{tab:conventions} and~\ref{tab:reliability}). A later full-grid analysis adds disagreement on row creation (Table~\ref{tab:alpha_full_grid}). The 0.667 and 0.800 benchmarks guide readings of temporal stability, not of validity \parencite[p.~241]{krippendorff2004contentanalysis}. Pooled-cell $\alpha$ cannot establish each element's stability, and the re-code's 50 shared technology rows say little about rare elements.

The \emph{supplementary re-code} (wave 2) covered the other 80 paragraphs at least 33 days after the first coding. It was declared before coding but was outside the registered plan. The first labels were hidden, but the coder knew the results, including the model's lower element count. In 25 paragraphs, a non-blind re-reading had shown one row's first-coding and model elements (Section~\ref{subsec:missed_results}). Both waves form the \emph{second coding}. Analyses using the supplementary re-code, alone or pooled, are exploratory and do not replace registered results.

\subsection{Evaluation Metrics}\label{subsec:metrics}

The metrics run from single labels to aggregate shares.

\paragraph{Precision, Recall and F1.} All agreement figures use one universe, the union of the (paragraph~$\times$~technology) rows of the reference and the first execution. Unmatched predicted rows are inventions whose marked elements are false positives. Unmatched reference rows are omissions whose marked elements are false negatives. A failed first execution yields only omissions, so headline figures count failures as errors. With $\mathrm{TP}_i$, $\mathrm{FP}_i$ and $\mathrm{FN}_i$ the tuple counts of paragraph $i$ (Figure~\ref{fig:anatomy}b), the design-weighted precision $P_w$, recall $R_w$ and F1 score $F_1^{w}$ are
\begin{equation}
P_w = \frac{\sum_i w_i\,\mathrm{TP}_i}{\sum_i w_i\,(\mathrm{TP}_i + \mathrm{FP}_i)}, \qquad
R_w = \frac{\sum_i w_i\,\mathrm{TP}_i}{\sum_i w_i\,(\mathrm{TP}_i + \mathrm{FN}_i)}, \qquad
F_1^{w} = \frac{2\,P_w R_w}{P_w + R_w}.
\label{eq:microf1}
\end{equation}
Recall below precision, a \emph{recall deficit}, means that the model marks fewer element tuples than the reference, in weighted count. Technology detection is scored likewise, as micro-F1 over (paragraph~$\times$~technology) rows.

\paragraph{Matched and End-to-End Scores.} Further scores separate value agreement from finding the frame or row. \emph{Matched-frame accuracy} scores evaluations on shared (paragraph, technology, frame) units, and \emph{matched-row accuracy} scores comparison targets on shared rows. \emph{End-to-end F1} compares (paragraph, technology, frame, evaluation) or (paragraph, technology, target) tuples as sets, so missed or invented frames and rows also count as errors.

\paragraph{Cluster Bootstrap.} Paragraphs of one article are not independent, so a cluster bootstrap recomputes each statistic on 2{,}000 resamples of whole articles \parencite[p.~370]{field2007bootstrapping}. Its 2.5th and 97.5th percentiles form a descriptive 95\% interval showing how far a figure moves when articles are redrawn. The interval ignores the outlet strata and claims no finite-population coverage (Table~\ref{tab:conventions}). A per-class cell prints its value if the reference supports the class with at least $s = 5$ units, and its support otherwise. This sealed threshold is a display convention, not evidence of statistical adequacy.

\paragraph{Run-to-Run Stability.} Stability is measured per dimension and language on executions ending with technologies present or off topic (Appendix~\ref{app:evaluator_conventions}). With $U_i$ the units of paragraph $i$, $E_i$ its successful executions and $v_{u,e}$ the value of unit $u$ in execution $e$,
\begin{equation}
\delta_i = \frac{1}{|U_i|\binom{|E_i|}{2}} \sum_{u \in U_i} \sum_{e < e'} \mathbb{1}\bigl[v_{u,e} \neq v_{u,e'}\bigr], \qquad
\Delta = \frac{\sum_i w_i\, \delta_i}{\sum_i w_i}.
\label{eq:stability}
\end{equation}
$\Delta$ is the weighted share of unit values, such as a row's element set, that change between two successful executions, and $\mathbb{1}[\cdot]$ is 1 if its condition holds and 0 otherwise.

\paragraph{Prevalence Displacement.} Displacement compares class shares at paragraph level (bottom of Figure~\ref{fig:anatomy}b). For a class $c$, such as a technology, an element or a (frame, evaluation) pair, $y_i = 1$ if any reference row of paragraph $i$ carries $c$, and $\hat{y}_i$ likewise for the first execution. Over the sample paragraphs, pooled or within one language, the displacement of class $c$ in percentage points of the eligible frame is
\begin{equation}
D_c = 100 \cdot \frac{\sum_i w_i\,(\hat{y}_i - y_i)}{\sum_i w_i}.
\label{eq:displacement}
\end{equation}
A negative $D_c$, a model prevalence below the reference, is a \emph{downward displacement}. As false positives and negatives can cancel, a small displacement does not establish accurate individual labels. The registered analysis applies no bias correction.

\subsection{Experimental Setup}\label{subsec:experimental_setup}

Claude Sonnet 4.6, the model the pipeline was built against, was kept without prior comparison, and the registered selection among four candidates lapsed (Table~\ref{tab:deviations}). Calls used Amazon Bedrock's Converse application programming interface (API) in \texttt{eu-central-1} under the European Union (EU) cross-region inference profile (Table~\ref{tab:model_config}). Decoding used temperature zero (Section~\ref{subsec:stream_llm}), which does not guarantee identical output. Truncated, malformed or rule-violating responses are retried once with the identical request, then recorded as parse failures. No label is edited afterwards, and no paragraph is resubmitted at the coder's discretion.

\subsection{Analyses Added after the Registered Result}\label{subsec:added_methods}\label{subsec:procedure}

The \emph{later analyses}, added after the registered result was known, ask whether its pattern recurs, how far it depends on the reference and how its prevalences can be adjusted (Table~\ref{tab:added_analyses}). They are descriptive, leave the registered figures unchanged and mostly reuse the inspected evaluation sample. Most were specified before their runs. Table~\ref{tab:deviations} dates the exceptions, including two configuration controls accepted afterwards. The \emph{fresh census} codes the 55 eligible-frame paragraphs of ten articles with no documented prior use in the study (Figure~\ref{fig:sample_flow}), but its reference was coded after the registered result was known.

Three of the later analyses check the instrument itself (Table~\ref{tab:instrument_checks}). The \emph{perturbation battery} rewords single codebook definitions, and the \emph{mandatory form} asks for a yes or no on every element per technology, still in one call. The \emph{retry candidate}, otherwise the registered instrument, feeds the validator's messages back in its one retry.

The \emph{design-based adjustment} turns the displacement of Section~\ref{subsec:metrics} into a correction (Figure~\ref{fig:dsl_scheme}). The design-based supervised learning estimator \parencite{egami2023dsl} adds the sample's design-weighted mean residual, reference minus model, which equals $-D_c$ in percentage points, to the model's share over the 618 eligible-frame paragraphs. The model's predictions, from three runs of the unchanged instrument (Figure~\ref{fig:sample_flow}), are assumed but not verified to be generated independently of sampling (Appendix~\ref{app:dsl_details}). The result estimates paragraph-level class prevalences under the reference's coding convention, not individual labels or element shares per technology.

% 2026-11-16 (author comment C35): new hand-drawn schematic of the design-based adjustment of Section 5.6, the method
% behind Figure fig:dsl. Thesis figure style (preamble): sans-serif, palette by meaning, no measured quantity drawn to
% scale. Colours: model-only selection figModelLight fill with figModel outline; model share figModel outline; coder
% sample and residual figInk; the adjusted share figAdjusted at 1pt (model and coder combined); the later-analysis tag
% figChanged. Greyscale cues: fill (light = model only), outline weight (thick = adjusted), and a direct label on
% every box. Worked example, investment, industry or competitiveness (F2b), as annotated in Figure fig:dsl:
%   first coding:  tab:dsl row F2b: model 19.4, registered displacement -18.3 (residual = -D_c, as the design
%                  weights sum to the 618 frame paragraphs, Eq. eq:weights), adjusted 37.7 [31.1, 44.3];
%   second coding: tab:checks_4p panel (f), investment, second-coding column: -13.6; tab:checks_4o panel (6),
%                  adjusted, second coding: 33.1 [27.1, 39.0]. 19.4 + 13.6 = 33.0 against the printed 33.1 is
%                  rounding of unrounded inputs (19.417 in images/fig_dsl.tex), hence "rounded" in the source note.
\begin{figure}[b]
\centering
\begin{tikzpicture}[x=1mm, y=1mm, font=\figfontsmall,
  box/.style={draw=figInk, line width=0.5pt, rectangle, rounded corners=1.5pt, align=center, text width=34mm,
    inner sep=2pt, minimum height=14mm, fill=white,
    execute at begin node={\hyphenpenalty=10000\exhyphenpenalty=10000}},
  boxmodel/.style={box, draw=figModel},
  storemodel/.style={boxmodel, fill=figModelLight},
  adjbox/.style={box, draw=figAdjusted, line width=1pt, text width=36mm, minimum height=24mm},
  sumnode/.style={draw=figInk, circle, line width=0.6pt, inner sep=0.8pt, font=\figfont\bfseries},
  arr/.style={-{Latex[length=1.8mm]}, line width=0.6pt},
  arrmodel/.style={arr, draw=figModel},
  arrhuman/.style={arr, draw=figInk},
  arradj/.style={arr, draw=figAdjusted, line width=1pt},
  lbl/.style={font=\figfontsmall, text=figGreyA, inner sep=1pt, align=center},
]
% inputs (left column), the two estimates (centre), their sum (right)
\node[storemodel] (frame) at (17,15) {\textbf{Eligible frame}\\ 618 paragraphs,\\ model labels only};
\node[box] (sample) at (17,-15) {\textbf{Probability sample}\\ 120 paragraphs, coder's\\ reference and model labels};
\node[boxmodel, text width=38mm] (share) at (72,15) {\textbf{Model share}\\ over all 618 paragraphs\\ 19.4\,\%};
\node[box, text width=38mm] (resid) at (72,-15) {\textbf{Mean residual}, design-weighted,\\ reference minus model ($-D_c$)\\ first coding: $+18.3$ pp\\ second coding: $+13.6$ pp};
\node[sumnode] (sum) at (101,0) {$+$};
\node[adjbox, anchor=west] (adj) at (108,0) {\textbf{Adjusted share}\\[1pt] first coding:\\ 37.7\,\% [31.1, 44.3]\\[1pt] second coding:\\ 33.1\,\% [27.1, 39.0]};
\node[font=\figfontsmall\itshape, text=figChanged, anchor=south east, inner sep=1pt] at ([yshift=1mm]adj.north east) {later analysis};
% edges
\draw[arr, draw=figGreyA] (frame.south) -- (sample.north) node[lbl, midway, right=1mm, align=left] {stratified\\ random draw};
\draw[arrmodel] (frame.east) -- (share.west);
\draw[arrhuman] (sample.east) -- (resid.west);
\draw[arrmodel] (share.east) -| (sum.north);
\draw[arrhuman] (resid.east) -| (sum.south);
\draw[arradj] (sum.east) -- (adj.west);
% source note (figure text, not counted)
\node[lbl, anchor=north west, text width=150mm, align=left] at (0,-27) {Investment, industry or competitiveness, in \% of the 618 eligible-frame paragraphs, values rounded. pp: percentage points. First coding: Table~\ref{tab:dsl}. Exploratory second-coding results: Tables~\ref{tab:checks_4o} and~\ref{tab:checks_4p}. $D_c$: Equation~\eqref{eq:displacement}. Brackets: approximate 95\% design-based intervals. The second coding combines the registered 40-paragraph re-code, made before the model run, and the supplementary 80-paragraph re-code, made after the results.};
\end{tikzpicture}
\caption[Design-based adjustment]{\textbf{The adjustment adds the sample's mean residual to the model's share, here for investment.}}
\label{fig:dsl_scheme}
\end{figure}



\addtocontents{toc}{\protect\needspace{11\baselineskip}}
\section{Results}\label{sec:results}
Against the first coding, agreement is high on technologies and shared-frame evaluations, but the model marks fewer frame elements and most reported element shares are lower. Later reference checks show that this recall deficit's size depends on which of the coder's two codings is the reference. Registered results answer RQ2 in Sections~\ref{subsec:agreement} to~\ref{subsubsec:categorical}, RQ3 in~\ref{subsec:stability_results} and RQ4 in~\ref{subsec:displacement_results}, beside these checks. Sections~\ref{subsec:dsl_results} and~\ref{subsec:recurrence} cover later analyses (Table~\ref{tab:added_analyses}): prevalence adjustment (Figure~\ref{fig:dsl}) and further articles and models (Figure~\ref{fig:model_comparison}).

Unless stated otherwise, intervals and per-class reporting follow Section~\ref{subsec:metrics}, and estimates are design-weighted to the 618-paragraph eligible frame. The \emph{first coding} is the registered evaluation reference; registered temporal stability uses the 40 paragraphs re-coded before the model run. Model comparisons with that re-code are later analyses. The \emph{second coding} combines those 40 with the supplementary 80, re-coded after the results. Analyses using the latter are exploratory.

\subsection{Technologies Agree Closely, Elements Less}\label{subsec:agreement}

Technology detection reaches a micro-F1 of 0.935 [0.901, 0.963], whereas the registered headline, element micro-F1, is 0.643 [0.581, 0.705] (Table~\ref{tab:primary}). Every eligible-frame paragraph passed the keyword prefilter, so its screening is not evaluated. Excluding the 15 paragraphs from calibration or development articles changes no headline figure by more than 0.01 (Appendix~\ref{app:detailed_results}).

% Generated by tools/pipeline/figures.py tables from a2_sonnet_20261013/metrics.json (sha256 9d69e490475a885f7216ceb5c49d295b9d7e55662b7092ca6faa1636e46c57f4). Do not edit by hand.
\begin{table}[htbp]
\centering
\footnotesize
\setlength{\tabcolsep}{4pt}
\begin{threeparttable}
\caption[Agreement with the first coding by dimension and language]{The model's first-execution agreement with the first coding, by dimension and language. Design-weighted point estimates; brackets: 2.5--97.5 percentile intervals from 2{,}000 article resamples (56 pooled; 28 per language); $n$: paragraphs.}
\label{tab:primary}
\begin{tabular}{@{}>{\raggedright\arraybackslash}p{24mm}>{\raggedright\arraybackslash}p{28mm}>{\raggedright\arraybackslash}p{31mm}>{\raggedright\arraybackslash}p{31mm}>{\raggedright\arraybackslash}p{31mm}@{}}
\toprule
\textbf{Dimension} & \textbf{Score} & \textbf{Pooled ($n$\,=\,120)} & \textbf{German ($n$\,=\,60)} & \textbf{English ($n$\,=\,60)} \\
\midrule
Technology detection & micro-F1 & 0.935 [0.901, 0.963] & 0.884 [0.781, 0.952] & 0.960 [0.933, 0.984] \\
Frame elements & micro-F1 & 0.643 [0.581, 0.705] & 0.621 [0.532, 0.710] & 0.653 [0.569, 0.731] \\
Evaluation & matched-frame accuracy & 0.874 [0.812, 0.924] & 0.841 [0.757, 0.916] & 0.890 [0.807, 0.953] \\
 & end-to-end F1 & 0.617 [0.555, 0.679] & 0.597 [0.498, 0.694] & 0.626 [0.542, 0.706] \\
Comparison target & matched-row accuracy & 0.794 [0.697, 0.879] & 0.881 [0.792, 0.951] & 0.755 [0.616, 0.872] \\
 & end-to-end F1 & 0.742 [0.647, 0.831] & 0.779 [0.661, 0.868] & 0.725 [0.590, 0.844] \\
\bottomrule
\end{tabular}
\begin{tablenotes}\footnotesize
\item Scores: Section~\ref{subsec:metrics}. Matched units per column: 208/94/114 frames, 158/69/89 rows.
\end{tablenotes}
\end{threeparttable}
\end{table}


A later comparison scores the prefilter's keywords against the first coding. They lead the model in point estimate on the five technologies with their own pattern but trail it across all seven (Table~\ref{tab:checks_4p}). The coder had seen their suggestions in a discarded attempt.

\label{subsec:run_record}\label{subsec:sensitivity}%
The two parts of the eligible frame come from different outlets, so per-language differences are not effects of language alone. Technology micro-F1 is 0.884 in German and 0.960 in English, and element micro-F1 is 0.621 and 0.653 (Table~\ref{tab:primary}).

All five parse failures in the 120 first executions fell on German paragraphs, as did 11 of the 13 in all 360 executions (Table~\ref{tab:attrition}). Later checks find lower German technology F1 end to end, but no clear element difference (Table~\ref{tab:checks_4o}). The five failed paragraphs carry ten of the fifteen missed German technology rows. Among valid first executions, a selected subset, the technology contrast is inconclusive (Table~\ref{tab:checks_4p}).

A later comparison of translated and original German input is inconclusive for element micro-F1 (Appendix~\ref{app:language}).

\subsection{Fewer Elements Than the First Coding}\label{subsec:directional}\label{subsubsec:per_element}

Agreement varies more than twofold across the fifteen reported frame elements (Figure~\ref{fig:elements}a, Table~\ref{tab:element_prf}). Seven reach an F1 of 0.7 or more, led by intermittency at 0.830. Four fall below 0.5: land requirement, landscape, technical performance and regulation, the last two with a registered re-code $\alpha$ of 0.462 and 0.598.

% Generated by tools/pipeline/figures.py from a2_sonnet_20261013/metrics.json (sha256 9d69e490475a885f7216ceb5c49d295b9d7e55662b7092ca6faa1636e46c57f4). Do not edit by hand.
% Palette 2026-11-06: panel b marks model blue (figBelow = figModel) below and orange (figAbove) above the reference, grey on it; panel a takes model blue from the preamble styles pooled, precisionmark and recallmark. Shapes, fills and positions unchanged.
\begin{figure}[htbp]
\centering
\begin{tikzpicture}
\begin{axis}[thesisfig, name=pa,
  width=50mm, height=100mm, scale only axis,
  xmin=0, xmax=1, xtick={0,0.2,0.4,0.6,0.8,1}, xticklabels={0,0.2,0.4,0.6,0.8,1}, xlabel={Agreement with the reference},
  ymin=0.4, ymax=15.6, y dir=reverse, rowaxis,
  ytick={1,2,3,4,5,6,7,8,9,10,11,12,13,14,15},
  yticklabels={{Intermittency (F5a)},{Cost or price (F1a)},{Subsidies (F1b)},{Grid integration (F5b)},{Security of supply (F4b)},{Technical feasibility (F6a)\,\dag},{Import dependence (F4a)},{Investment, industry (F2b)},{Deployment growth (F2c)\,\dag},{Emissions, climate (F3a)},{Transition pathway (F3b)\,\dag},{Landscape (F7b)},{Regulation (F8a)\,\dag},{Technical performance (F6b)\,\dag},{Land requirement (F7a)}},
  legend columns=3,
]
\draw[figGreyB, line width=0.8pt] (axis cs:0.7890,1) -- (axis cs:0.8750,1);
\draw[figGreyB, line width=0.8pt] (axis cs:0.6430,2) -- (axis cs:0.8830,2);
\draw[figGreyB, line width=0.8pt] (axis cs:0.6740,3) -- (axis cs:0.8170,3);
\draw[figGreyB, line width=0.8pt] (axis cs:0.6890,4) -- (axis cs:0.7960,4);
\draw[figGreyB, line width=0.8pt] (axis cs:0.6990,5) -- (axis cs:0.7610,5);
\draw[figGreyB, line width=0.8pt] (axis cs:0.6230,6) -- (axis cs:0.8600,6);
\draw[figGreyB, line width=0.8pt] (axis cs:0.5540,7) -- (axis cs:1.0000,7);
\draw[figGreyB, line width=0.8pt] (axis cs:0.4950,8) -- (axis cs:0.9760,8);
\draw[figGreyB, line width=0.8pt] (axis cs:0.6130,9) -- (axis cs:0.6170,9);
\draw[figGreyB, line width=0.8pt] (axis cs:0.4810,10) -- (axis cs:0.6640,10);
\draw[figGreyB, line width=0.8pt] (axis cs:0.4860,11) -- (axis cs:0.5990,11);
\draw[figGreyB, line width=0.8pt] (axis cs:0.3200,12) -- (axis cs:1.0000,12);
\draw[figGreyB, line width=0.8pt] (axis cs:0.3200,13) -- (axis cs:0.7880,13);
\draw[figGreyB, line width=0.8pt] (axis cs:0.3910,14) -- (axis cs:0.4850,14);
\draw[figGreyB, line width=0.8pt] (axis cs:0.2190,15) -- (axis cs:0.6990,15);
\addplot[precisionmark] coordinates {(0.7890,1) (0.8830,2) (0.8170,3) (0.7960,4) (0.6990,5) (0.8600,6) (1.0000,7) (0.9760,8) (0.6170,9) (0.6640,10) (0.5990,11) (1.0000,12) (0.7880,13) (0.3910,14) (0.6990,15)};
\addlegendentry{Precision}
\addplot[recallmark] coordinates {(0.8750,1) (0.6430,2) (0.6740,3) (0.6890,4) (0.7610,5) (0.6230,6) (0.5540,7) (0.4950,8) (0.6130,9) (0.4810,10) (0.4860,11) (0.3200,12) (0.3200,13) (0.4850,14) (0.2190,15)};
\addlegendentry{Recall}
\addplot[pooled] coordinates {(0.8300,1) (0.7441,2) (0.7387,3) (0.7386,4) (0.7286,5) (0.7229,6) (0.7129,7) (0.6573,8) (0.6151,9) (0.5581,10) (0.5363,11) (0.4844,12) (0.4551,13) (0.4329,14) (0.3336,15)};
\addlegendentry{F1 (registered)}
\end{axis}
\node[panellabel] at ([yshift=6mm]pa.north -| current bounding box.west) {a};
\begin{axis}[thesisfig, name=pb, at={(pa.south east)}, anchor=south west, xshift=12mm,
  width=38mm, height=100mm, scale only axis,
  xmin=-30, xmax=20, xtick={-30,-20,-10,0,10,20},
  xlabel={Displacement (percentage points)},
  ymin=0.4, ymax=15.6, y dir=reverse, rowaxis, ytick=\empty,
]
\draw[refline] (axis cs:0,0.4) -- (axis cs:0,15.6);
\addplot[only marks, mark=*, draw=figAbove, mark options={draw=figAbove, fill=white, line width=0.8pt}, error bars/.cd, x dir=both, x explicit, error bar style={figAbove}] coordinates {(1.607,1) -= (3.487,0) += (3.644,0) (3.425,11) -= (7.563,0) += (8.282,0) (2.670,14) -= (6.203,0) += (6.696,0)};
\addplot[only marks, mark=*, draw=figBelow, mark options={draw=figBelow, fill=white, line width=0.8pt}, error bars/.cd, x dir=both, x explicit, error bar style={figBelow}] coordinates {(-5.690,2) -= (8.857,0) += (8.175,0) (-0.771,3) -= (3.804,0) += (3.907,0) (-2.508,4) -= (3.082,0) += (2.508,0) (-2.427,7) -= (3.206,0) += (2.427,0) (-1.998,9) -= (6.530,0) += (6.298,0) (-2.739,10) -= (4.327,0) += (3.991,0)};
\addplot[only marks, mark=*, draw=figBelow, mark options={draw=figBelow, fill=figBelow, line width=0.8pt}, error bars/.cd, x dir=both, x explicit, error bar style={figBelow}] coordinates {(-4.924,6) -= (4.369,0) += (3.877,0) (-18.303,8) -= (6.121,0) += (6.070,0) (-6.230,12) -= (5.701,0) += (4.627,0) (-8.202,13) -= (7.080,0) += (6.277,0) (-5.596,15) -= (5.446,0) += (4.812,0)};
\addplot[only marks, mark=*, draw=figGreyA, mark options={draw=figGreyA, fill=white, line width=0.8pt}, error bars/.cd, x dir=both, x explicit, error bar style={figGreyA}] coordinates {(0.000,5) -= (4.444,0) += (4.271,0)};
\node[figtext, anchor=east, xshift=17mm] at (axis cs:20,1) {\textcolor{figGreyA}{suppressed}};
\node[figtext, anchor=east, xshift=17mm] at (axis cs:20,2) {0.678};
\node[figtext, anchor=east, xshift=17mm] at (axis cs:20,3) {0.899};
\node[figtext, anchor=east, xshift=17mm] at (axis cs:20,4) {\textcolor{figGreyA}{suppressed}};
\node[figtext, anchor=east, xshift=17mm] at (axis cs:20,5) {\textcolor{figGreyA}{suppressed}};
\node[figtext, anchor=east, xshift=17mm] at (axis cs:20,6) {0.625};
\node[figtext, anchor=east, xshift=17mm] at (axis cs:20,7) {\textcolor{figGreyA}{suppressed}};
\node[figtext, anchor=east, xshift=17mm] at (axis cs:20,8) {0.788};
\node[figtext, anchor=east, xshift=17mm] at (axis cs:20,9) {0.369};
\node[figtext, anchor=east, xshift=17mm] at (axis cs:20,10) {\textcolor{figGreyA}{suppressed}};
\node[figtext, anchor=east, xshift=17mm] at (axis cs:20,11) {0.562};
\node[figtext, anchor=east, xshift=17mm] at (axis cs:20,12) {\textcolor{figGreyA}{suppressed}};
\node[figtext, anchor=east, xshift=17mm] at (axis cs:20,13) {0.598};
\node[figtext, anchor=east, xshift=17mm] at (axis cs:20,14) {0.462};
\node[figtext, anchor=east, xshift=17mm] at (axis cs:20,15) {\textcolor{figGreyA}{suppressed}};
\node[figtext, anchor=south east, xshift=17mm] at (axis cs:20,0.4) {Re-code $\alpha$};
\end{axis}
\node[panellabel] at ([xshift=-3mm, yshift=6mm]pb.north west) {b};
\end{tikzpicture}
\caption[Agreement and prevalence displacement per frame element]{\textbf{Against the first coding, most frame elements have lower recall than precision and lower paragraph shares.} Design-weighted, 120 paragraphs. Elements with at least 5 reference rows, sorted by F1. \textbf{a},~F1 and, added later, precision and recall over (paragraph, technology, element) tuples. \textbf{b},~Paragraph-share displacement (model minus reference), with 2.5--97.5 percentile intervals. Filled: interval excludes zero. Blue: below, orange: above, grey: on the reference (dashed). Paragraph shares count an element on any row, so the panels can differ. Transition pathway: precision 0.599, recall 0.486, yet displacement $+3.4$ [$-4.1$, $+11.7$] points. Right: unweighted $\alpha$ of the registered re-code (40 paragraphs), suppressed below 5 first-coding positives. \dag: $\alpha < 0.667$.}
\label{fig:elements}
\end{figure}


\label{subsubsec:error_decomposition}%
A later split of registered predictions gives pooled element precision 0.76 and recall 0.56 against the first coding, a recall deficit (Table~\ref{tab:checks_4o}). Precision exceeds recall for 12 of the 15 reported elements. For derived frames, detection precision is at least 0.677 for every reported frame, whereas recall falls to 0.301 for siting (Table~\ref{tab:frame_prf}).

\label{subsec:reliability_results}%
The reference's temporal stability limits interpretation. In the registered 40-paragraph re-code, unweighted $\alpha$ ranges from 0.369 to 0.899 across eight reportable elements (Figure~\ref{fig:elements}). Later article-bootstrap intervals for all eight (Table~\ref{tab:alpha_full_grid}) reach below 0.667, the lowest level at which \textcite[p.~242]{krippendorff2004contentanalysis}, with stated hesitation, would accept reliability data. Stability over time is the weakest form of reliability (p.~215). Pooled-cell $\alpha$ is 0.641 [0.555, 0.723] on shared rows and 0.619 [0.530, 0.697] on the full grid. Each element result is therefore read beside its own $\alpha$ where one is reportable (Figure~\ref{fig:elements}).

On the 40 paragraphs re-coded before the model run, neither model score is shown equivalent to the coder's own retest. Scored later in the model's metric, the two codings agree at an element F1 of 0.656 (Table~\ref{tab:checks_4p}). Against the first coding, the model reaches 0.571, below the retest, with no interval for the difference. Against the second coding, it reaches 0.620, which differs from the retest by $-0.035$ [$-0.129$, $+0.063$] (Table~\ref{tab:recode_waves}).

On these 40 paragraphs, the recall deficit narrows against the second coding but persists in point estimate. Model recall there is 0.584, above that against the first coding by $+0.088$ [$+0.002$, $+0.178$], a contrast declared before the supplementary re-code (Table~\ref{tab:recode_comparison}). Precision is 0.662, so their gap is $+0.078$, against $+0.177$ with the first coding. Its interval and paired change, both defined after the supplementary re-code was inspected, are in Table~\ref{tab:recode_waves}.

Exploratory checks on all 120 paragraphs show that the second coding drops most first-coding tuples the model misses but few that it finds. It no longer holds 102 of the 168 misses and 36 of the 208 hits (unweighted; Table~\ref{tab:checks_4p}). Against the full second coding, precision and recall nearly balance in point estimate, at 0.695 and 0.666, as they do without the 25 re-read paragraphs (Table~\ref{tab:recode_waves}). This balance mixes elements with opposite gaps (Table~\ref{tab:recode_elements}) and does not establish unbiased paragraph shares. These checks measure one coder's agreement over time, not an intercoder ceiling or evidence of validity. Neither coding is established as correct, and a dropped tuple is not thereby a reference error.

\label{subsec:missed_results}%
A later re-reading of 45 randomly drawn misses within matched rows points in both directions: 16 clear misses, 15 possible first-coding over-marks and 10 borderline. Of the other 4, two rest on cues outside the paragraph (Table~\ref{tab:missed_elements}). Investment accounts for 13 of the 45, only 3 of them clear misses. These are one coder's non-blind judgements of their own reference, not error rates.

\subsection{Evaluations Agree on Shared Frames, Comparison Targets Less}\label{subsubsec:categorical}

Against the first coding, matched-frame accuracy is 0.874 (Table~\ref{tab:primary}). A later comparator that always chooses the most frequent value, positive, reaches 0.630 on the same 208 units and weights (Table~\ref{tab:checks_4p}). Positive and negative evaluations agree best, with per-value F1 of 0.916 and 0.933. Neutral and mixed, with 11 and 9 matched units, reach only 0.558 and 0.571 (Table~\ref{tab:per_value}).

End-to-end F1, also counting unmatched frames and rows, falls to 0.617. The model missed 105 reference frames within matched rows, and 20 reference technology rows with 27 frame--evaluation units.

In a later comparison on the 40 paragraphs re-coded before the model run, the coder's two codings agree on evaluations at 0.876. The model reaches 0.802 against the first coding on these paragraphs. The matched sets differ, so these accuracies are not ranked.

The comparison target agrees less well. Its matched-row accuracy is 0.794 (Table~\ref{tab:primary}), but a later constant guess of ``none'' already reaches 0.674. On German rows, the guess scores 0.887 against the model's 0.881 (Appendix~\ref{app:later_checks}). The comparison target's end-to-end macro-F1, the unweighted mean of per-class F1, is 0.379. Whole rows match exactly at a rate of 0.215 over 184 rows (Table~\ref{tab:primary_full}).

\subsection{Outputs Vary, Scores Stay Similar}\label{subsec:stability_results}

Across the three executions of the identical request, element micro-F1 ranges from 0.632 to 0.648. Averaged over pairs of successful executions, however, a paragraph's output differs in at least one label at a rate of 0.150 [0.091, 0.213] over 117 paragraphs (Table~\ref{tab:stability}). A paragraph's technology set changes at 0.037, and a technology row's element set at 0.088 [0.047, 0.133] over 115 paragraphs.

\subsection{Most Element Prevalences Are Lower}\label{subsec:displacement_results}

% 2026-11-11: rewritten for the fourth review's plan (M16); earlier versions and number sources in git (state d5a2bfb).
% 2026-11-24: the instrument-check block moved to tab_instrument_checks (appendix).
% OUTLINE: The overview retains adjustment, recurrence, translation and exploratory reference sensitivity; instrument checks have a separate appendix overview, and no later analysis replaces a registered figure.
\begin{table}[!tbp]
\centering
\scriptsize
\setlength{\tabcolsep}{3pt}
\renewcommand{\arraystretch}{1.1}
\begin{threeparttable}
\caption[Analyses added after the registered result]{Descriptive analyses added after the registered result. None replaces a registered figure. Instrument checks: Table~\ref{tab:instrument_checks}. Reference and context checks: Appendix~\ref{app:later_checks}.}
\label{tab:added_analyses}
\begin{tabular}{@{}>{\raggedright\arraybackslash}p{36mm}>{\raggedright\arraybackslash}p{50mm}>{\raggedright\arraybackslash}p{65mm}@{}}
\toprule
\textbf{Analysis (appendix tables)} & \textbf{Question; material} & \textbf{Result} \\
\midrule
Prevalence adjustment (Table~\ref{tab:dsl}) & Eligible-frame class shares? 618 model-labelled paragraphs, 120 human-coded & Investment 37.7\% (model only 19.4\%) under the first coding's convention \\
\addlinespace
Fresh census (Table~\ref{tab:fresh_displacement}) & Recurrence on articles without prior use? 55 paragraphs of ten articles; reference coded after the registered result & Element F1 0.563, precision above recall; eight of ten reported elements displaced downward \\
\addlinespace
Other models (Tables~\ref{tab:model_config} and~\ref{tab:pr_gaps}) & Is the gap model-specific? Both samples, mostly one execution each & In point estimate, precision above recall against both first codings in 11 of 12 configurations, all but Claude Opus~5 \\
\addlinespace
Translated input (Appendix~\ref{app:language}) & Does translation change agreement? 60 German evaluation paragraphs & Inconclusive for agreement \\
\addlinespace
Supplementary re-code (exploratory; Tables~\ref{tab:alpha_waves} and~\ref{tab:recode_waves}) & Reference sensitivity? 80 paragraphs re-coded after the results & Second coding, both waves (120 paragraphs): precision 0.695, recall 0.666; 290 element tuples against 376 in the first coding \\
\bottomrule
\end{tabular}
\begin{tablenotes}\footnotesize
\item Evaluation sample design-weighted, census unweighted. F1: element micro-F1. Claude~5 and 5.5: provider's default sampling.
\end{tablenotes}
\end{threeparttable}
\end{table}


Against the first coding, paragraph-level displacement is downward for eleven of the fifteen reported elements, upward for three and zero for one (Figure~\ref{fig:elements}b; all classes in Table~\ref{tab:displacement}). Investment, industry or competitiveness is displaced most, by $-18.3$ [$-24.4$, $-12.2$] points. The first coding marks it on 44 of 120 paragraphs and the model on 24. Five downward article intervals exclude zero, but no upward one does. Under later design-based intervals, four remain, as technical feasibility drops out (Section~\ref{subsec:dsl_results}; Table~\ref{tab:checks_4p}). The intervals are pointwise and unadjusted for the fifteen elements.

Pooled technology prevalences differ by at most 5 points. Within German, solar differs by $-10.1$, an undecomposed end-to-end figure counting failed executions as unmarked.

Among frame--evaluation pairs, positive evaluations in the jobs frame are displaced most, by $-13.5$ [$-19.6$, $-7.6$] points. As employment has two reference paragraphs, investment and deployment growth carry this frame.

An exploratory check against the second coding gives nine lower and six higher element shares (Table~\ref{tab:checks_4p}). Land requirement, among the nine, is below the reporting threshold there. Two article intervals lie above zero, for transition pathway and technical performance (registered re-code $\alpha$ 0.562 and 0.462), and one below, for investment at $-13.6$ [$-20.4$, $-7.6$]. Only investment has an article interval below zero under both readings.

Sample size and reference stability limit reporting. The reporting threshold suppresses displacements for employment, the residual other frame and 15 of 31 frame--evaluation pairs marked by either side. Reported element intervals are 5.6 to 17.0 points wide, so even the $-5.7$ points of cost or price are indistinguishable from zero. Scored later on the 40 paragraphs re-coded before the model run, the two codings agree on derived-frame presence at an F1 of 0.736 [0.642, 0.819] (Table~\ref{tab:checks_4p}). Such agreement over time does not validate the frames or their shares.

\subsection{Adjusted Prevalences Shift with the Reference Coding}\label{subsec:dsl_results}\label{subsec:after_registered}\label{subsec:instrument_checks}

The exploratory adjustment of Section~\ref{subsec:added_methods} targets a chosen coding convention rather than correcting reference error. The model marks investment on 19.4\% of eligible-frame paragraphs, but its sample share is 18.3 points below the first coding's. The adjustment therefore estimates investment's eligible-frame prevalence at 37.7\% [31.1, 44.3] under the first coding (Figure~\ref{fig:dsl}) and, with the supplementary re-code included, 33.1\% [27.1, 39.0] under the second. Technical feasibility (registered re-code $\alpha$ 0.625) rises from the model's 16.5\% to 21.4\% under the first coding but falls to 12.8\% under the second (Tables~\ref{tab:dsl} and~\ref{tab:checks_4o}).

% Generated by tools/pipeline/figures.py dsl from dsl_prevalence.json (sha256 e3990a480c22851f83744c21892dd4fcd80f0f2b1ead760eeb6be2da4179bb63). Do not edit by hand.
% Analysis added after the registered result (specification §4d); exploratory.
% Palette 2026-11-06: adjusted estimate green (figAdjusted, model and coder combined), model-only share model blue (figModel), sample-only estimate black (figInk); shapes, fills and row offsets unchanged, the legend keys the series (caption shortened 2026-11-09, K-18).
\begin{figure}[htbp]
\centering
\begin{tikzpicture}
\begin{axis}[thesisfig,
  width=72mm, height=78mm, scale only axis,
  xmin=-2, xmax=50, xtick={0,10,20,30,40,50}, xlabel={Prevalence in the eligible frame (\% of paragraphs)},
  ymin=0.4, ymax=15.6, y dir=reverse, rowaxis,
  ytick={1,2,3,4,5,6,7,8,9,10,11,12,13,14,15},
  yticklabels={{Investment, industry (F2b)},{Cost or price (F1a)},{Technical feasibility (F6a)},{Deployment growth (F2c)\,\ddag},{Regulation (F8a)},{Transition pathway (F3b)\,\ddag},{Subsidies (F1b)},{Grid integration (F5b)},{Landscape (F7b)},{Emissions, climate (F3a)},{Land requirement (F7a)},{Intermittency (F5a)},{Import dependence (F4a)},{Security of supply (F4b)},{Technical performance (F6b)\,\ddag}},
  legend columns=1, legend style={at={(1,0)}, anchor=south west, xshift=3mm, yshift=13mm},
]
\fill[figGreyC] (axis cs:-2,0.5) rectangle (axis cs:50,1.5);
\draw[guideline] (axis cs:0,0.4) -- (axis cs:0,15.6);
\addplot[only marks, mark=x, mark size=2.2pt, draw=figModel, mark options={draw=figModel, line width=0.8pt}] coordinates {(19.417,0.73) (20.550,1.73) (16.505,2.73) (19.256,3.73) (10.032,4.73) (21.036,5.73) (14.401,6.73) (9.385,7.73) (3.722,8.73) (6.796,9.73) (1.780,10.73) (8.576,11.73) (3.398,12.73) (5.178,13.73) (7.605,14.73)};
\addlegendentry{Model only}
\addplot[only marks, mark=*, draw=figAdjusted, fill=figAdjusted, mark options={draw=figAdjusted, fill=figAdjusted}, error bars/.cd, x dir=both, x explicit, error bar style={figAdjusted}] coordinates {(37.720,1.00) -= (6.619,0) += (6.619,0) (26.241,2.00) -= (6.480,0) += (6.480,0) (21.429,3.00) -= (5.135,0) += (5.135,0) (21.254,4.00) -= (7.179,0) += (7.179,0) (18.235,5.00) -= (5.630,0) += (5.630,0) (17.611,6.00) -= (7.007,0) += (7.007,0) (15.172,7.00) -= (3.772,0) += (3.772,0) (11.893,8.00) -= (2.602,0) += (2.602,0) (9.951,9.00) -= (4.360,0) += (4.360,0) (9.535,10.00) -= (3.790,0) += (3.790,0) (7.376,11.00) -= (4.372,0) += (4.372,0) (6.969,12.00) -= (3.157,0) += (3.157,0) (5.825,13.00) -= (2.527,0) += (2.527,0) (5.178,14.00) -= (3.963,0) += (3.963,0) (4.935,15.00) -= (5.218,0) += (5.218,0)};
\addlegendentry{Adjusted, 95\% interval}
\addplot[only marks, mark=*, draw=figInk, mark options={draw=figInk, fill=white, line width=0.8pt}, error bars/.cd, x dir=both, x explicit, error bar style={figInk}, error mark options={rotate=90, mark size=1.6pt, line width=0.8pt, draw=figInk}] coordinates {(40.056,1.27) -= (8.459,0) += (8.459,0) (25.301,2.27) -= (7.565,0) += (7.565,0) (17.869,3.27) -= (6.657,0) += (6.657,0) (23.122,4.27) -= (6.704,0) += (6.704,0) (16.915,5.27) -= (5.958,0) += (5.958,0) (21.146,6.27) -= (6.770,0) += (6.770,0) (15.347,7.27) -= (6.024,0) += (6.024,0) (10.183,8.27) -= (4.018,0) += (4.018,0) (8.670,9.27) -= (4.851,0) += (4.851,0) (9.201,10.27) -= (4.776,0) += (4.776,0) (8.023,11.27) -= (4.586,0) += (4.586,0) (9.331,12.27) -= (5.131,0) += (5.131,0) (7.349,13.27) -= (4.048,0) += (4.048,0) (7.727,14.27) -= (4.916,0) += (4.916,0) (7.673,15.27) -= (3.721,0) += (3.721,0)};
\addlegendentry{Sample only, $\pm$\,1.96 SE}
\node[figtext, anchor=north west, align=left, xshift=3mm, inner sep=0pt] at (axis cs:50,0.55) {\textbf{Investment} (worked example)\\adjusted 37.7\% [31.1, 44.3]\\sample only 40.1\% [31.6, 48.5]\\model only 19.4\%};
\end{axis}
\end{tikzpicture}
\caption[Adjusting element prevalences with the coded sample]{\textbf{The adjustment raises 11 of 15 element prevalences above model-only shares, investment most, and narrows 12 intervals relative to the sample alone.} Shares in \% of 618 eligible-frame paragraphs, under the first coding's convention. Adjusted: design-based supervised learning (Section~\ref{subsec:added_methods}), 95\% intervals (normal approximation, no coverage guarantee). Sample only: $\pm$\,1.96 design-based standard errors (SE). Model only: failed executions (43 paragraphs) unmarked. \ddag: adjusted interval wider than sample-only, as their model errors vary more than the coder's labels. Fixed model predictions exclude run-to-run variation from both intervals.}
\label{fig:dsl}
\end{figure}


\subsection{Most Configurations Share the Recall Deficit}\label{subsec:recurrence}

% Generated by tools/pipeline/compare_models.py (sha256 a491c5f986c48a28) from a2_sonnet_20261013, mc_opus46_a2_20261019, mc_sonnet45_a2_20261019, mc_haiku45_a2_20261019, mc_novapro_a2_20261019, mc_nemotron3s_a2_20261019r3, mc_qwen3_235b_a2_20261019, mc_mercury3_instant_a2_20261030, mc_mercury3_low_a2_20261030, mc_sonnet5_a2_20261101, mc_opus5_a2_20261101, mc_opus55_a2_20261108r2, mc_sonnet_req5_a2_20261101, mc_opus_req5_a2_20261101, mc_opus5_req55_a2_20261108r2, a3_sonnet_20261019, mc_opus46_a3_20261019, mc_sonnet45_a3_20261019, mc_haiku45_a3_20261019, mc_novapro_a3_20261019, mc_nemotron3s_a3_20261019, mc_qwen3_235b_a3_20261019, mc_mercury3_instant_a3_20261030r3, mc_mercury3_low_a3_20261030r2, mc_sonnet5_a3_20261101, mc_opus5_a3_20261101, mc_opus55_a3_20261108r2, mc_sonnet_req5_a3_20261101, mc_opus_req5_a3_20261101, mc_opus5_req55_a3_20261108r2; robustness model comparison (specification §4f; Mercury 3 §4g; Claude 5 §4k; Opus 5.5 §4m), not the registered result. Do not edit by hand.
% Palette 2026-11-06: the registered model, its joining line and the sample legend images model blue (figModel); Opus 5 (pentagon) and, since 2026-11-08, Opus 5.5 (triangle) vermillion (figHighlight = figChanged), the other configurations grey (figContext); the caption says blue diamonds.
\begin{figure}[htbp]
\centering
\begin{tikzpicture}
\begin{axis}[thesisfig,
  width=80mm, height=80mm, scale only axis,
  xmin=0.2, xmax=1, ymin=0.2, ymax=1, xtick={0.2,0.4,0.6,0.8,1}, ytick={0.2,0.4,0.6,0.8,1},
  xticklabels={0.2,0.4,0.6,0.8,1}, yticklabels={0.2,0.4,0.6,0.8,1},
  xlabel={Element recall}, ylabel={Element precision},
  legend columns=3,
]
\draw[refline] (axis cs:0.2,0.2) -- (axis cs:1,1);
\node[figtext, text=figGreyA, rotate=45, anchor=south] at (axis cs:0.30,0.30) {precision = recall};
\addplot[figGreyB, densely dotted, line width=0.6pt, domain=0.2505:1, samples=60, forget plot] {0.4*x/(2*x-0.4)};
\node[figtext, text=figGreyA, anchor=south west, inner sep=1pt] at (axis cs:0.251,1) {F1 = 0.4};
\addplot[figGreyB, densely dotted, line width=0.6pt, domain=0.4291:1, samples=60, forget plot] {0.6*x/(2*x-0.6)};
\node[figtext, text=figGreyA, anchor=south west, inner sep=1pt] at (axis cs:0.429,1) {F1 = 0.6};
\addplot[figGreyB, densely dotted, line width=0.6pt, domain=0.6672:1, samples=60, forget plot] {0.8*x/(2*x-0.8)};
\node[figtext, text=figGreyA, anchor=south west, inner sep=1pt] at (axis cs:0.667,1) {F1 = 0.8};
\draw[figContext, line width=0.4pt, shorten <=2pt, shorten >=2pt] (axis cs:0.5860,0.7941) -- (axis cs:0.6286,0.7395);
\draw[figContext, line width=0.4pt, shorten <=2pt, shorten >=2pt] (axis cs:0.4659,0.7625) -- (axis cs:0.5071,0.8068);
\draw[figContext, line width=0.4pt, shorten <=2pt, shorten >=2pt] (axis cs:0.2891,0.7208) to[bend right=35] (axis cs:0.3929,0.8088);
\draw[figContext, line width=0.4pt, shorten <=2pt, shorten >=2pt] (axis cs:0.2711,0.4942) -- (axis cs:0.3714,0.4727);
\draw[figContext, line width=0.4pt, shorten <=2pt, shorten >=2pt] (axis cs:0.3143,0.6577) -- (axis cs:0.3286,0.5897);
\draw[figContext, line width=0.4pt, shorten <=2pt, shorten >=2pt] (axis cs:0.4052,0.7061) -- (axis cs:0.4500,0.6632);
\draw[figContext, line width=0.4pt, shorten <=2pt, shorten >=2pt] (axis cs:0.3079,0.7043) -- (axis cs:0.3143,0.7458);
\draw[figContext, line width=0.4pt, shorten <=2pt, shorten >=2pt] (axis cs:0.4548,0.8003) -- (axis cs:0.5571,0.8864);
\draw[figContext, line width=0.4pt, shorten <=2pt, shorten >=2pt] (axis cs:0.5295,0.8110) -- (axis cs:0.5571,0.7800);
\draw[figModel, line width=0.6pt, shorten <=2pt, shorten >=2pt] (axis cs:0.5572,0.7591) -- (axis cs:0.4929,0.6571);
\draw[figHighlight, line width=0.6pt, shorten <=2pt, shorten >=2pt] (axis cs:0.7314,0.7479) -- (axis cs:0.7357,0.6398);
\draw[figHighlight, line width=0.6pt, shorten <=2pt, shorten >=2pt] (axis cs:0.6603,0.8061) -- (axis cs:0.7571,0.7910);
\addplot[context, forget plot] coordinates {(0.5860,0.7941)};
\addplot[context, open, forget plot] coordinates {(0.6286,0.7395)};
\addplot[context, forget plot] coordinates {(0.4659,0.7625)};
\addplot[context, open, forget plot] coordinates {(0.5071,0.8068)};
\addplot[context, forget plot] coordinates {(0.2891,0.7208)};
\addplot[context, open, forget plot] coordinates {(0.3929,0.8088)};
\addplot[context, forget plot] coordinates {(0.2711,0.4942)};
\addplot[context, open, forget plot] coordinates {(0.3714,0.4727)};
\addplot[context, forget plot] coordinates {(0.3143,0.6577)};
\addplot[context, open, forget plot] coordinates {(0.3286,0.5897)};
\addplot[context, forget plot] coordinates {(0.4052,0.7061)};
\addplot[context, open, forget plot] coordinates {(0.4500,0.6632)};
\addplot[context, forget plot] coordinates {(0.3079,0.7043)};
\addplot[context, open, forget plot] coordinates {(0.3143,0.7458)};
\addplot[context, forget plot] coordinates {(0.4548,0.8003)};
\addplot[context, open, forget plot] coordinates {(0.5571,0.8864)};
\addplot[context, forget plot] coordinates {(0.5295,0.8110)};
\addplot[context, open, forget plot] coordinates {(0.5571,0.7800)};
\addplot[only marks, mark=diamond*, mark size=3.6pt, draw=white, fill=white, forget plot] coordinates {(0.5572,0.7591)};
\addplot[registered, forget plot] coordinates {(0.5572,0.7591)};
\addplot[only marks, mark=diamond*, mark size=3.6pt, draw=white, fill=white, forget plot] coordinates {(0.4929,0.6571)};
\addplot[registered, open, forget plot] coordinates {(0.4929,0.6571)};
\addplot[highlight, forget plot] coordinates {(0.7314,0.7479)};
\addplot[highlight, open, forget plot] coordinates {(0.7357,0.6398)};
\addplot[highlightb, forget plot] coordinates {(0.6603,0.8061)};
\addplot[highlightb, open, forget plot] coordinates {(0.7571,0.7910)};
\node[figtext, below=3pt, align=center, inner sep=1pt, fill=white] at (axis cs:0.4929,0.6571) {Sonnet 4.6\\(registered)};
\node[figtext, right=3pt, inner sep=1pt, fill=white, text=figHighlight, font=\figfontsmall\bfseries] at (axis cs:0.7314,0.7479) {Opus 5};
\node[figtext, above=3pt, inner sep=1pt, fill=white, text=figHighlight, font=\figfontsmall\bfseries] at (axis cs:0.6603,0.8061) {Opus 5.5};
\node[figtext, text=figGreyA, below=2pt, inner sep=1pt, fill=white] at (axis cs:0.6286,0.7395) {Opus 4.6};
\node[figtext, text=figGreyA, above=2pt, inner sep=1pt, fill=white] at (axis cs:0.5571,0.8864) {Mercury 3 low};
\draw[figGreyA, line width=0.3pt] (axis cs:0.3143,0.8388) -- (axis cs:0.3143,0.7578);
\node[figtext, text=figGreyA, anchor=south, align=center, inner sep=1pt, fill=white] at (axis cs:0.3593,0.8388) {Mercury 3\\instant};
\draw[figGreyA, line width=0.3pt] (axis cs:0.2711,0.5472) -- (axis cs:0.2711,0.5062);
\node[figtext, text=figGreyA, anchor=south, align=center, inner sep=1pt, fill=white] at (axis cs:0.2711,0.5472) {Nova Pro};
\addlegendimage{registered}\addlegendentry{Sonnet 4.6, registered}
\addlegendimage{highlight}\addlegendentry{Opus 5}
\addlegendimage{highlightb}\addlegendentry{Opus 5.5}
\addlegendimage{context}\addlegendentry{Other configurations (9)}
\addlegendimage{context}\addlegendentry{Evaluation sample}
\addlegendimage{context, open}\addlegendentry{Fresh census}
\end{axis}
\end{tikzpicture}
\caption[Element precision against recall by model]{\textbf{In point estimate, eleven of twelve configurations mark fewer frame elements than the first coding on both samples: precision exceeds recall.} First executions; evaluation sample design-weighted, fresh census unweighted. Lines join each configuration's samples. Several configurations differed from the registered decoding settings (Table~\ref{tab:model_config}).}
\label{fig:model_comparison}
\end{figure}


\label{subsec:fresh_results}%
\enlargethispage{\baselineskip}%
Against the fresh census's first coding (Appendix~\ref{app:fresh}), the unchanged pipeline reaches a lower unweighted element micro-F1 of 0.563 [0.511, 0.622]. Its gap, precision minus recall, is $+0.164$ [$+0.071$, $+0.245$] end to end. A later check on the 47 paragraphs with a valid first execution, a selected subset, gives $+0.057$ [$-0.006$, $+0.108$] (Table~\ref{tab:checks_4o}). The census reference was coded after the registered result, its stability is unmeasured, and the intervals rest on ten articles.

\label{subsec:model_comparison}%
Against both first codings, eleven of twelve configurations have positive gaps in point estimate, all but Claude Opus~5 (Figure~\ref{fig:model_comparison}). Nine have intervals above zero on both samples (Table~\ref{tab:pr_gaps}). All twelve used the same texts, prompt and tool schema, so these are not independent replications; none was selected (Appendix~\ref{app:model_comparison}). A later valid-output check gives the same eleven positive gaps (Table~\ref{tab:checks_4o}). Exploratory rescoring against the second coding leaves fewer positive gaps; the registered model's gap interval includes zero (Table~\ref{tab:checks_4p}).

Against the first codings, Opus~5's gap is near zero on the evaluation sample and negative on the census. Three further executions reproduce this pattern (Appendix~\ref{app:model_comparison}). Under Opus~5's request settings, however, three other configurations kept the deficit.


\section{Discussion}\label{sec:discussion}
Against the first coding, the model marks fewer elements, and most element shares are lower. In a later check against the coder's re-code of 40 paragraphs, made before the model run, the recall deficit persists in point estimate but is smaller (Section~\ref{subsec:directional}). Opposite errors can partly cancel in aggregate figures, whereas one-directional errors shift them \parencite[p.~2091]{edwards2026rogangladen}. Aggregate shares therefore need a defended human reference. Sections~\ref{subsec:interpretation} and~\ref{subsec:mechanisms} interpret the pattern and weigh possible sources, none isolated by the design. Sections~\ref{subsec:theoretical_implications} to~\ref{subsec:limitations_future} cover implications, limitations and next steps.

\subsection{The Recall Deficit Depends on the Reference Reading}\label{subsec:interpretation}\label{subsec:comparison_prior}

In a later split against the first coding, element precision exceeds recall (Figure~\ref{fig:overview}b), whereas recall exceeds precision in 20 of 27 codebook-prompted task replications by \textcite{pangakis2025humansintheloop}. The higher recall against the re-code made before the model run shows this asymmetry to be reference-sensitive, not an established measure of model-side omission.

The gap recurs in most configurations, so the shared setup remains a candidate explanation. Yet a later fresh-census check finds the gap small among valid outputs, a selected subset (Section~\ref{subsec:recurrence}). Against the evaluation sample's first coding, a later Opus~5 run is near balance in point estimate, so the setup does not force a deficit.

\subsection{Possible Sources of Element Disagreement}\label{subsec:mechanisms}

Against the first coding, most first-execution misses occur on technology rows where the model marked another element (later check, Tables~\ref{tab:fn_causes} and~\ref{tab:checks_4p}). The union of all three registered executions raises recall only from 0.557 to 0.596, so the disagreements persist across executions. Persistence does not establish which side is wrong.

On the reference side, the second coding holds fewer element tuples than the first (Table~\ref{tab:added_analyses}). This may reflect threshold drift, a known risk of coding over time \parencite[p.~565]{song2020validations}, learning, memory or, after the results, knowledge of them, which the design cannot separate.

Two mechanisms from the literature, both conjectures here, could produce a deficit of this sign. A model asked for several labels may stop listing early, as autoregressive models tended to concentrate probability on one label per step \parencite[p.~2472]{ma2025multilabel}. Label bias, a preference for some answers regardless of the input, may also favour absence \parencite{reif2024labelbias}. Both studies examine open-weight models' output probabilities on other tasks, but Sonnet~4.6 returns only text. The mandatory form's recall gain, clear only on the fresh census and confounded with a doubled output cap, fits the first conjecture without isolating it (Table~\ref{tab:form_variant}). Under the first conjecture, Mercury~3, a diffusion model that does not generate strictly left to right, should show a narrower gap than autoregressive models, as declared before its runs \parencite{inception2026mercury3,inception2025mercury}. Its gap was larger than the registered model's, though model, temperature and provider also differed (Table~\ref{tab:pr_gaps}). A third candidate, ambiguous definitions, is not singled out: in the perturbation battery, the model's labels move with definition wording but also with frame order, element labels and a negative-control sentence (Appendix~\ref{app:battery}).

The coder's non-blind re-reading of misses points to both model and reference (Section~\ref{subsec:missed_results}). Its judgements are not error rates. Of the 13 sampled investment misses, 5 were judged possibly over-marked. Four of these report deployment growth or decline, a separate element (registered re-code $\alpha$ 0.369) that the model marked each time. Investment's comparatively high stability on shared rows thus does not show that its coding follows the intended boundary (Table~\ref{tab:alpha_waves}).

\subsection{Theoretical Implications}\label{subsec:theoretical_implications}\label{subsec:implications}

The instrument measures presence, not salience, and most elements can fire on a topical reference alone. The residual frame appears on two evaluation paragraphs in the first coding and on six of 55 in the census reference (Tables~\ref{tab:displacement} and~\ref{tab:fresh_displacement}). Neither count shows that the inventory is complete.

Measuring frames through elements \parencite{ali2025hydrogenframing,liu2023strategyframing} exposed differences a single label would hide. Within the feasibility frame, F1 is 0.723 for technical feasibility and 0.433 for technical performance, yet their later intervals overlap (Table~\ref{tab:checks_4o}). Technical performance is unstable in both re-code waves, so the contrast is tentative (Table~\ref{tab:alpha_waves}).

Labels that discriminate well can still yield biased shares when counted \parencite[pp.~7--8]{linder2026multicalibrated}, so label and prevalence agreement remain distinct. The later design-based adjustment moves shares towards the convention of its human-coded sample, whose accuracy its authors name as a condition \parencite[Sec.~6]{egami2023dsl}, and does not correct reference error (Section~\ref{subsec:dsl_results}). Under the second coding, made mostly after the results, it gives an exploratory sensitivity, not a better estimate. Label agreement, reference stability and derived shares thus need joint validation.

\subsection{Managerial Implications}\label{subsec:managerial_implications}

Media-monitoring teams can report the technologies detected in prefiltered paragraphs, and the evaluations of frames a coder confirms, once both are validated per language on their own material. Whole technology records rarely match exactly, so exporting them without review is not supported (Section~\ref{subsubsec:categorical}).

Element, frame and evaluation shares should not be reported without a human-coded probability sample. A blind re-code completed before model results are inspected assesses the coder's temporal stability. Re-reading sampled disagreements can indicate whether model or reference may err. Once the reference is secured, shares need adjustment towards a defended convention where warranted. A second coder can help defend the convention, but no step establishes validity by itself. This workflow transfers to other monitoring projects, which must define their own frames and elements and revalidate.

Each source, language, country, technology or period to be compared needs its own human-coded sample. When analysts compare frame shares across outlets, countries and periods, misses that vary by group can exaggerate, hide or reverse the difference human coding would show \parencite[p.~2092]{edwards2026rogangladen}. Against the first coding, German technology detection was lower end to end, mostly on paragraphs whose first execution failed. Among valid outputs, the contrast is inconclusive (Section~\ref{subsec:agreement}). All five such failures were German, but the fresh census had four per language (Table~\ref{tab:checks_4o}). With element displacement intervals 5.6 to 17.0 points wide, the present sample is too small per cell (Table~\ref{tab:topic_requirements}).

Model calls scale to larger corpora, but sourcing and validation do not yet. The registered run's 360 executions took 12.5 minutes and cost an estimated USD 2.18 in model calls. Coding and validation labour was not measured. Larger applications need a licensed or API-based source (Table~\ref{tab:topic_requirements}). They should budget retries, report attrition per language, fix configurations, journal every call and revalidate new model versions (Section~\ref{subsec:recurrence}).

\subsection{Limitations and Future Research}\label{subsec:limitations_future}\label{subsec:limitations}

Both human readings come from one coder, the author, and neither is independently validated. Calibration covered 32 paragraphs on one day. The final split of deployment growth from investment was not trial-coded before the freeze. The registered re-code came after as few as five days instead of 14, so memory may inflate its temporal stability and sustain the deficit against it. Imperfect human labels can also depress estimated classification performance, as simulations show \parencite[p.~560]{song2020validations}.

The supplementary re-code followed the registered results and the re-reading of misses, so it may have moved towards the model. Omitting the 25 re-read paragraphs does not remove knowledge of the results and leaves a non-random subset (Appendix~\ref{app:later_checks}). Model-agreement and prevalence intervals condition on the chosen labels and do not propagate uncertainty about alternative human readings. A discarded first attempt with keyword suggestions may have shaped technology ticks (Section~\ref{subsec:annotation}). The general category's figures also depend on its reading.

Model and coder differed in instructions (Section~\ref{subsec:codebook}) and in context. The coder saw headline, outlet, language and often other article paragraphs. The model saw only the paragraph, although the codebook's unit definition names headline and metadata. Context could reveal elements despite the shared paragraph-only rule and so inflate the apparent deficit. Adding headline, outlet and language in one later execution per sample did not consistently raise recall or narrow the gap (Table~\ref{tab:context_variant}) but left the coder's knowledge of other paragraphs untested. With one coder, mostly single runs of further models and variation of only the model's context, the deficit cannot be apportioned among model, codebook, reference and context.

The reference and context limits point to the next test. Two independent coders would apply a narrower, untested codebook to the model's paragraph-only input on fresh paragraphs. An element would require a technology-anchored claim (not a topic mention), a stated boundary to its neighbours and a rule on whether one cue may support two elements. The general category would need one explicit reading. Procedural reporting would need a settled boundary to regulation framing. Candidate rules are in Table~\ref{tab:battery_sharpen}. A model-assisted review of drafts may find gaps early (Table~\ref{tab:development_review}).

The sample of 120 paragraphs from 56 articles is small and clustered: the element micro-F1 interval spans about $\pm 0.06$ and per-language intervals up to $\pm 0.09$. Population estimates target the eligible frame of a purposive, hand-collected corpus, defined by a keyword prefilter with unmeasured misses. The later analyses (Table~\ref{tab:added_analyses}) and appendix checks mostly reuse inspected samples and report intervals unadjusted for multiplicity, so none is confirmatory. The evidence covers one topic, one codebook and two languages from different outlets, not other frames or domains.

The data rule restricts article text to the registered run's EU route. Exceptions were authorised for several later steps. Mercury~3 received evaluation, census, calibration and development paragraphs via its developer's interface, hosted at an unknown location. After the EU quota ran out, 62\% of battery executions used a global route, possibly outside the EU. Assistant agents outside the EU route read paragraphs to derive the battery's candidate rules. A language-model assistant consulted during the fresh census received census text but recommended no label, and none was changed (Appendix~\ref{app:reference_disclosures}). Some assistant data routes cannot be reconstructed from the archived material (Appendix~\ref{app:guide_llm}). The records document no rights basis for storing or externally processing the hand-collected texts, so they are withheld as a precaution (Table~\ref{tab:topic_requirements}).

Proprietary Sonnet~4.6 exposes no seed. Its provider may change or end access to this version \parencite{feuerriegel2026guidellm}. The coding interface's build is unarchived.

As the model comparisons reused inspected samples, production selection needs a fresh selection sample, a rule fixed beforehand and a separate, untouched evaluation sample.

The topic description's broader project asks how energy coverage changes around the 2024 United States presidential election. It needs a human-coded probability sample per country, language and period, a sample of keyword-rejected paragraphs to measure the filter's misses, and source, model version and data route fixed for the study period.


\section{Conclusion}\label{sec:conclusion}
This thesis asked how a multilingual LLM pipeline can be designed and validated for extracting technology frames from international renewable energy news. The question matters because firms, investors and agencies monitoring coverage can use frame shares only from a validated instrument. The answer rests on validation against one coder on a 120-paragraph probability sample, design-weighted to the 618-paragraph eligible frame of a keyword-prefiltered, purposive corpus.

On design (RQ1), one frozen codebook supplies common definitions for the model's prompt and output schema and the coder's sheets. Its exclusion rules reach them differently, so identical application is not established. Zero-shot coding reserves the coder's labels for validation. One stateless call codes a paragraph's technologies jointly, as one sentence may contrast two, and untranslated, as translation can change labels. A registry, a raw-response journal and a design-weighted evaluator make every prediction traceable to its configuration.

On agreement (RQ2), technology detection (micro-F1 0.935) and evaluations of shared frames (matched-frame accuracy 0.874) agree closely with the first coding. End-to-end evaluation F1 is 0.617. Comparison targets reach 0.794 matched-row accuracy, against 0.674 for a later constant guess of ``none''. In later contrasts, German technology detection is lower when failed executions count as misses but not clearly so among valid first executions, a selected subset. The languages also come from different outlets.

Frame elements agree less, with a micro-F1 of 0.643 against the first coding and no clear language difference. A later split gives precision 0.76 against recall 0.56, so the model marks fewer elements. In later checks on the 40 paragraphs re-coded before the model run, the model's element F1 is 0.571 against the first coding and 0.620 against the second. There, the recall deficit narrows against the second coding but persists in point estimate. The two codings agree at 0.656, one coder's agreement over time, not a ceiling or validity. Neither coding is established as correct, so the central finding is a recall deficit against one coder, its size depending on the reference reading.

On stability (RQ3), identical repeated executions give similar scores but not identical outputs. Element micro-F1 ranges from 0.632 to 0.648. Yet some label of a paragraph changes between two successful executions at a design-weighted rate of 0.150.

On aggregate prevalences (RQ4), 11 of 15 reported element shares lie below the first coding's in point estimate, five with pointwise intervals excluding zero. Against the second coding, made mostly after the results, an exploratory sensitivity gives nine lower and six higher, two of the latter with intervals above zero. Only investment lies below both codings with intervals excluding zero, by 18.3 and 13.6 percentage points. A later design-based adjustment raises investment's share under either coding's convention but does not correct reference error. Sample size sets the limit: classes with fewer than five supporting reference paragraphs stay suppressed, and intervals 5.6 to 17.0 points wide make small displacements indistinguishable from zero.

The contribution is a tested combination, not a new model or method. It links label agreement, output variation, the coder's stability and share displacement, and, in later analyses, the reference dependence of adjusted shares. Once validated per language on the material at hand, such a pipeline can report technologies and the evaluations of confirmed frames. Its element and frame shares need their own check against a human-coded probability sample, with a reference secured before adjustment.

With one coder who had more context than the model, the deficit's sources cannot be apportioned. A later Opus~5 run shows that the setup does not force a deficit. Neither a second coder nor re-reading disagreements establishes validity alone. The next test is whether a narrower codebook, applied by two independent coders to the model's paragraph-only input, raises agreement between the coders and with the model on fresh paragraphs.

Other monitoring projects that define their own frames and revalidate can reuse the validation workflow.


\newpage
\pagenumbering{Roman}
\setcounter{page}{\value{frontmatterpages}} % continues the front matter's Roman numbering
\begingroup
\raggedright
\setlength{\emergencystretch}{3em} % long URLs and compound words: no line into the margin
\printbibliography[title={References}, heading=bibintoc]
\endgroup

\appendix
\section{Appendix}\label{app:main}

\makeatletter
\newcommand{\appfloatbarrier}{\par\begingroup\let\@elt\relax\edef\@tempa{\@botlist\@deferlist\@dbldeferlist}%
  \ifx\@tempa\@empty\endgroup\else\endgroup\clearpage\fi}
\newcommand{\appcaption}[2]{\let\app@origcaption\caption
  \def\caption{\let\caption\app@origcaption\@ifnextchar[{\app@capopt{#1}{#2}}{\app@capopt{#1}{#2}[]}}}
\long\def\app@capopt#1#2[#3]#4{\if\relax\detokenize{#2}\relax\caption[#1]{#4}\else\caption[#1]{#2}\fi}
\long\def\app@grabto#1\app@eof{\long\gdef\app@content{#1}}
\newcommand{\appgrab}[1]{\begingroup\everyeof{\app@eof}\expandafter\app@grabto\@@input #1 \endgroup}
\newcommand{\appcontains}[1]{\long\def\app@test##1#1##2\app@end{\if\relax\detokenize{##2}\relax\in@false\else\in@true\fi}%
  \expandafter\app@test\app@content#1\app@end}
\newcommand{\appLTcaption}[2]{\let\app@origLTc@ption\LT@c@ption
  \long\def\LT@c@ption##1[##2]##3{\app@origLTc@ption##1[#1]{#2}}}
\long\def\app@ifcontains#1#2{\begingroup\edef\x{\endgroup\noexpand\in@{\detokenize{#2}}{\detokenize{#1}}}\x
  \ifin@\expandafter\@firstoftwo\else\expandafter\@secondoftwo\fi}
\newtoks\app@keep
\long\def\app@rowloop#1\\{\app@ifcontains{#1}{app@rowend}{}{%
  \app@ifdrop{#1}{}{\app@keep\expandafter{\the\app@keep#1\\}}\app@rowloop}}
\long\def\app@ifdrop#1{\begingroup\edef\x{\endgroup\noexpand\in@{\app@droppattern}{\detokenize{#1}}}\x
  \ifin@\expandafter\@firstoftwo\else\expandafter\@secondoftwo\fi}
\long\def\app@ltsplit#1\endlastfoot#2\end#3\app@stop{\app@keep{}\app@rowloop#2\app@rowend\\%
  \expandafter\app@ltout\expandafter{\the\app@keep}{#1}{#3}}
\long\def\app@ltout#1#2#3{#2\endlastfoot#1\end#3}
\newcommand{\appdroprows}[1]{\edef\app@droppattern{\detokenize{#1}}\expandafter\app@ltsplit\app@content\app@stop}
\ExplSyntaxOn
\cs_new_protected:Npn \appreplace #1#2 { \regex_replace_all:nnN {#1} {#2} \app@content }
\ExplSyntaxOff
\newcommand{\appthousands}{\appreplace{over\ 2000\ }{over\ 2\cB\{,\cE\}000\ }}
\newcommand{\appinput}[1]{\appgrab{#1}\appthousands\app@content}
\newcommand{\appkeepheads}{\appreplace{\cE\}\cE\}\ *\c{\\}}{\cE\}\cE\}\c{\\}*}\appreplace{\c{\\}\ *\c{midrule}}{\c{\\}*\c{noalign}\cB\{\c{nobreak}\c{vskip}2pt\c{hrule}\ height\ 0.5pt\c{nobreak}\c{vskip}2pt\c{nobreak}\cE\}}}
\newcommand{\appout}{\app@content}
\newcommand{\appneedspace}[1]{\par\begingroup\setlength{\dimen@}{#1}\vskip\z@\@plus\dimen@\penalty-100
  \vskip\z@\@plus-\dimen@\vskip\dimen@\penalty9999\vskip-\dimen@\vskip\z@skip\endgroup}
\begingroup
\renewenvironment{table}[1][]{\par\addvspace{\intextsep}\noindent\begin{minipage}{\linewidth}\def\@captype{table}%
  \let\footnotesize\scriptsize\let\small\footnotesize}%
  {\end{minipage}\par\addvspace{\intextsep}}
\makeatother

\noindent Several generated tables below keep the short labels of the analysis code. Table~\ref{tab:label_key} gives the term of the main text for each.

\begin{table}[htbp]\centering\footnotesize
\caption[Key to the labels of the analysis code]{Key to the short labels of the analysis code and their terms in the main text.}\label{tab:label_key}
\begin{tabular}{@{}>{\raggedright\arraybackslash}p{36mm}>{\raggedright\arraybackslash}p{112mm}@{}}
\toprule
\textbf{Label in a table} & \textbf{Term in the main text} \\
\midrule
A0, A0b; A1 & calibration sessions; development split \\
A2; A2R; A2R2 & evaluation sample (120 paragraphs); the coder's registered blind re-code of 40 of them (wave 1); the supplementary re-code of the other 80 (wave 2) \\
A3; fresh sample & fresh census (55 paragraphs of ten articles) \\
score~A & matched-frame accuracy (evaluation) or matched-row accuracy (comparison target) \\
score~B & end-to-end F1 \\
DSL estimate & adjusted prevalence (design-based supervised learning, Equation~\eqref{eq:dsl_estimator}) \\
\bottomrule
\end{tabular}
\end{table}

\subsection{Studies That Shaped the Design}\label{app:literature}

\begingroup\scriptsize
\setlength{\tabcolsep}{4pt}
\begin{longtable}{@{}>{\raggedright\arraybackslash}p{26mm}>{\raggedright\arraybackslash}p{38mm}>{\raggedright\arraybackslash}p{48mm}>{\raggedright\arraybackslash}p{38mm}@{}}
\caption[Studies that shaped the design]{Studies that shaped the design, by stream: the material and method, the finding this thesis relies on, and what was taken from it.}\label{tab:literature} \\
\toprule
\textbf{Study} & \textbf{Material and method} & \textbf{Finding relied on} & \textbf{Taken into this thesis} \\
\midrule
\endfirsthead
\toprule
\textbf{Study} & \textbf{Material and method} & \textbf{Finding relied on} & \textbf{Taken into this thesis} \\
\midrule
\endhead
\bottomrule
\endfoot
\multicolumn{4}{@{}l}{\emph{Media framing of energy technologies}} \\
\textcite{stephens2009windmedia} & 678 wind-power articles, three United States newspapers, 1990--2007; manual framing & aesthetic-risk frame in 44, 14 and 6\% of the papers' articles; the political frame dropped for failing reliability; 129 of 678 articles (19\%) link wind to climate & frames differ by place; evaluation coded separately from the frame; reliability reported per element, not used to drop one \\
\textcite{rochyadireetz2019framing} & 1{,}772 articles, eleven democracies; Entman's frame elements; cluster analysis & positive frame dominant in every country, Germany least positive; reliability described only procedurally & element-based coding; storage and transmission as subjects in their own right \\
\textcite{ali2025hydrogenframing} & 2{,}292 hydrogen articles, four countries; 61 frame components; factor analysis & agreement reported only as one average over all elements; 25.3\% of variance explained; rare elements excluded & a fixed element-to-frame mapping instead of factor analysis \\
\addlinespace
\multicolumn{4}{@{}l}{\emph{Computational analysis of energy and climate discourse}} \\
\textcite{toetzke2022energydiscourse} & about five million articles, fourteen outlets, seven countries; planned translation to English; paragraph-level classifiers & design only (workshop slides); 27 context labels from the literature and an expert study & the paragraph as coding unit; keyword matching to technologies; its labels related to the frames (Table~\ref{tab:frame_crosswalk}) \\
\textcite{guo2023openframing} & 1{,}300 English and 350 German headlines; one classifier per language & precision 0.83 and 0.73; frame distributions compared across languages without a bilingual test & one instrument validated per language on a bilingual sample \\
\addlinespace
\multicolumn{4}{@{}l}{\emph{Large language models as annotators}} \\
\textcite{gilardi2023chatgpt} & 6{,}183 texts; zero-shot ChatGPT against crowd workers and trained annotators & on average about 25 points more accurate than the crowd; raw agreement only; human coding of the 14 policy frames in news articles unsuccessful & the codebook as the prompt; agreement reported beyond raw percentages \\
\textcite{alizadeh2025opensource} & the same group's tweets and news articles; GPT-3.5, LLaMA and FLAN fine-tuned on 50--1{,}500 texts, GPT-4 zero-shot as comparison & fine-tuned open-weight models match or surpass zero-shot GPT-3.5 and GPT-4, FLAN and LLaMA-1 coming close to zero-shot GPT-3.5 at about 250 texts; few-shot prompting inconsistent; recommend at least two runs with reported agreement and decoding at zero or low temperature & the labels spent on measurement rather than fine-tuning; repeated executions with an agreement statistic \\
\textcite{ziems2024css} & 25 English benchmarks, thirteen models & rarely beat fine-tuned RoBERTa baselines; in four analysed tasks, many errors from preferring the most generic label; results averaged over paraphrased prompts & one fixed prompt; per-class reporting \\
\textcite{reif2024labelbias} & 279 classification tasks, ten open-weight models; bias measured directly & a preference for certain labels regardless of the input, substantial after standard calibration and few-shot fine-tuning; anti-correlated with performance across demonstration sets & label bias as one possible, untested mechanism of the recall deficit; the aggregate corrected, not the instrument \\
\textcite{lupo2024humanlevelcoding} & Swedish policy documents; three hierarchical codebook tasks & full definitions with fifteen examples best; joint coding of interdependent labels at almost a third of the cost; the model's codes close to those of an average human coder & full definitions rendered into the prompt; one joint call per paragraph \\
\textcite{toernberg2025outperform} & social-media posts, eleven countries; one English prompt & highest accuracy in every country in point estimate, not significantly so in all; lower accuracy in multi-party systems, not fully explained by language & one untranslated instrument; no claim about a language effect \\
\textcite{pangakis2025humansintheloop} & 27 tasks from eleven datasets; GPT-4 & median F1 0.707, a third of the tasks below 0.5 on precision or recall; a 250-item development subset and 1{,}000 held-out items; temperature without notable effect in the seven tasks tested & a development split and one evaluation sample; stability measured, not tuned \\
\textcite{atreja2025whatsinaprompt} & four tasks including news frames; sixteen prompts, four models & definitions raise news-frame accuracy (GPT-3.5, GPT-4o, PaLM 2); asking for an explanation has a mixed effect; news frames 0.75 accuracy, 0.71 macro-F1 for GPT-4o & definitions rendered; the explanation field tested as a variant, not assumed \\
\textcite{hohenwalde2025chatgptactors} & German news, actor categories; GPT-3.5 and GPT-4 models at temperature 0 & the human codebook alone insufficient, decomposition helped; 100 development and 200 test articles; team $\alpha = .63$ (.77 after merging the science categories) & the German benchmark for a single reference; the split precedent \\
\textcite{baumann2025llmhacking} & 37 tasks, eighteen models, 2{,}361 hypotheses & wrong conclusions in about 31\% of hypotheses for state-of-the-art models and about half for small ones; a frame task among the riskiest for some models; for controlling false positives, 100 human labels beat 100{,}000 model labels & a configuration fixed in advance, one reported run, a human reference \\
\textcite{enomoto2025fairtranslationese} & instructions with translationese removed & which instruction language works better varies with the task and the language of the labels & no claim from the German--English difference \\
\addlinespace
\multicolumn{4}{@{}l}{\emph{Validation and aggregates}} \\
\textcite{song2020validations} & review of validation practice; simulation & less reliable human labels distort estimated performance, mostly downward; more than 1{,}300 items (1\% of a simulated 130{,}000-article corpus) advised at $\alpha \geq 0.7$ & the reference's temporal stability measured and reported; the sample size named as a limit \\
\textcite{egami2023dsl} & design-based supervised learning estimator & predicted-label estimates valid only under restrictive error assumptions; correction from a human-coded probability sample & a probability sample with design weights; the correction as a declared exploratory analysis only \\
\textcite{angelopoulos2023ppi} & prediction-powered inference; seven applications in six scientific domains & valid intervals from predictions plus a small gold set, for any predictor; no gain when predictions are poor or the unlabelled set is small relative to the gold set & the design-based interval of the exploratory estimate; the expected gain stated per class \\
\textcite{linder2026multicalibrated} & one language model against expert-coded topic labels; four countries, four languages & per-language AUC 0.983--0.994 of the probability scores, yet biased prevalences from standard estimators; counted binary labels (AUC 0.960) biased by 1.6 to 4.8 points; about 13{,}400 calibration labels & aggregate displacement as a registered quantity \\
\textcite{pontiki2014semeval,pontiki2016semeval} & SemEval aspect-based sentiment tasks & set-valued scoring: precision over returned, recall over gold annotations; polarity scored on gold aspects & matched-frame accuracy and end-to-end F1; the value \emph{mixed} (from the 2014 task) \\
\textcite{card2020littlepower} & power analysis for NLP evaluation & on a 147-item test set with a 94.5\% baseline, the smallest gain detectable at 80\% power is about 5.3 points & no small-difference claims from 32 or 55 paragraphs \\
\textcite{halterman2026codebookllms} & three political-science codebooks; label-free behavioural tests & label-set compliance and category-order sensitivity tested before accuracy & instrument checks named for a later version \\
\addlinespace
\multicolumn{4}{@{}p{156mm}@{}}{\scriptsize Read as preprints \parencite{gilardi2023chatgpt,jumle2024findingframes,lupo2024humanlevelcoding,nemkova2025annotationskills,xing2024crosslingualinconsistency,hasan2026englishrewrites,liang2025multilingualbehaviour,linder2026multicalibrated,baumann2025llmhacking,reif2024labelbias}; \textcite{toetzke2022energydiscourse} was consulted as workshop slides; its design, its planned steps and its corpus figures are taken from them.} \\
\end{longtable}
\endgroup

\subsection{Corpus and Registry}\label{app:corpus_details}

\begingroup\scriptsize
\setlength{\tabcolsep}{4pt}
\begin{longtable}{@{}>{\raggedright\arraybackslash}p{30mm}>{\raggedright\arraybackslash}p{124mm}@{}}
\caption[Corpus collection, exclusions and registry]{Corpus collection, exclusions and registry. The frame and sample counts are in Table~\ref{tab:corpus} and Figure~\ref{fig:sample_flow}.}\label{tab:corpus_facts} \\
\toprule
\textbf{Item} & \textbf{Record} \\
\midrule
\endfirsthead
\toprule
\textbf{Item} & \textbf{Record} \\
\midrule
\endhead
\bottomrule
\endfoot
Collection & Without a database, assisted by Claude Code (overview of tools used, part~B), 3 August to 2 September 2026. Batches 1--4: 35 articles retrieved on 3 and 4 August through web search and the author's browser sessions (paywalled outlets through the author's subscriptions, The Economist through the university's licence). Batch 5: 60 browser prints supplied by the author and transcribed on 1 and 2 September, with optical character recognition for raster prints; each print is kept beside its transcript. \\
Database route & Factiva explored, not used (Table~\ref{tab:topic_requirements}). The licence position relayed on 1 October 2026 (ETH Zurich library: automated queries under the academic Factiva licence are not permissible) closes that route. \\
Selection cases & The collection brief asked for at least one article per case that breaks an extraction pipeline; articles per case (an article may serve several): evaluation voiced by a quoted source 35; mixed or neutral evaluation 29; several technologies with conflicting evaluations 23; procedural reporting without a frame 20; nuclear energy against renewables 5; energy security without a stated mechanism 3. \\
Manifest & Outlet, country, language, publication date, headline, canonical URL and retrieval date for each of the 95 articles; publication dates from 28 February 2023 to 2 September 2026, 57 articles from 2026. \\
Segmentation & Body prose split into paragraphs of at least 15 words: 1,047 paragraphs (686 English, 361 German) from the 74 articles with body prose. \\
Exclusions & 21 articles without body prose (19 extracts, 2 metadata only). Article A011, one paragraph of which is quoted in the prompt, is excluded from the frame. Reserve candidates with a marked transcription gap (18) are excluded from the frame; all 26 paragraphs with a gap stay in the registry, flagged. \\
Prefilter & A keyword prefilter over the seven technologies, tuned for recall, marks 696 paragraphs (66\%) as candidates; it decides frame membership only and was hidden from the coder, with one exception (Table~\ref{tab:disclosures}). \\
Registry & Each paragraph has an immutable identity (article, paragraph, hash of the segmented text). Three closed corpus versions: 198 identities (29 August), 679 and 170 (2 September); each binds the hashes of its input files, generator code and interpreter, so later drift is detected. Splits are read from an append-only, hash-chained assignment log, never recomputed. A containment check with eight-word shingles at a threshold of 0.15 guards against re-entry under a new identity; its quarantine is empty. \\
Splits & Calibration 20 (A0) and 12 (A0b), development 32 (A1), reserve 983. The development session coded 35 paragraphs: the 32 of its split and three that the registry assigns to A0b. \\
Freeze and draw & Freezing the frame appended a terminal registry event whose reference is the frame manifest's digest; nothing can enter the reserve afterwards. Allocation proportional to stratum size within language, largest remainder, cap of 24 per outlet (binding only for The New York Times, 24 of 184); within a stratum, paragraphs sorted by text digest and drawn with seed 20260905. The registered re-code drew 20 paragraphs per language with seed 20260906. \\
Sample dates & Paragraphs from 2026, 2025, 2024 and 2023: 70, 36, 8 and 6 (articles 35, 13, 5 and 3 of 56); time is not a stratum, and no per-period figure is reported. \\
\end{longtable}
\endgroup

% Generated by tools/annotation/descriptives.py from the corpus manifest, the split
% registry, the frozen A2 frame manifest and the registered A2 sample manifest. Do not edit by hand.
% Appendix A.2 since 2026-11-09 (slimming debate consensus F-03): the main text carries the summary tab:corpus.
\begin{table}[htbp]
\centering
\footnotesize
\setlength{\tabcolsep}{3pt}
\begin{threeparttable}
\caption[Corpus, eligible frame and evaluation-sample allocation by outlet]{Corpus, eligible frame and evaluation-sample allocation by outlet; each outlet with eligible paragraphs is one stratum within its language, with stratum size $N_h$ in the eligible frame, draw $n_h$ and design weight $w_h = N_h/n_h$.}
\label{tab:corpus_outlets}
\begin{tabular}{@{}llrrr@{\hspace{12pt}}rrr@{}}
\toprule
 & & \multicolumn{3}{c}{\textbf{Corpus}} & \multicolumn{3}{c}{\textbf{Stratum}} \\
\cmidrule(lr){3-5}\cmidrule(l){6-8}
\textbf{Outlet} & \textbf{Country} & \textbf{Articles} & \textbf{Paragraphs} & \textbf{Candidates} & $\boldsymbol{N_h}$ & $\boldsymbol{n_h}$ & $\boldsymbol{w_h}$ \\
\midrule
\multicolumn{8}{@{}l}{\textit{German-language outlets}} \\
taz & Germany & 8 & 80 & 64 & 45 & 12 & 3.75 \\
DIE ZEIT & Germany & 7 & 52 & 44 & 24 & 7 & 3.43 \\
Süddeutsche Zeitung & Germany & 4 & 40 & 23 & 23 & 6 & 3.83 \\
BILD & Germany & 6 & 38 & 28 & 28 & 8 & 3.50 \\
Neue Zürcher Zeitung & Switzerland & 2 & 36 & 22 & 22 & 6 & 3.67 \\
tagesschau.de & Germany & 3 & 31 & 24 & 24 & 6 & 4.00 \\
Handelsblatt & Germany & 2 & 26 & 12 & 12 & 3 & 4.00 \\
DER SPIEGEL & Germany & 3 & 21 & 18 & 18 & 5 & 3.60 \\
Hamburger Abendblatt & Germany & 1 & 11 & 9 & 9 & 2 & 4.50 \\
BR24 & Germany & 1 & 8 & 7 & 7 & 2 & 3.50 \\
Berliner Morgenpost & Germany & 1 & 7 & 4 & 4 & 1 & 4.00 \\
STERN.de & Germany & 1 & 7 & 6 & 5 & 1 & 5.00 \\
Salzburger Nachrichten (SN.at) & Austria & 1 & 4 & 4 & 4 & 1 & 4.00 \\
\textit{All German} & & 40 & 361 & 265 & 225 & 60 & \\
\addlinespace
\multicolumn{8}{@{}l}{\textit{English-language outlets}} \\
The New York Times & United States & 13 & 311 & 184 & 184 & 24 & 7.67 \\
The Economist & United Kingdom & 12 & 162 & 125 & 108 & 18 & 6.00 \\
The Guardian & United Kingdom & 4 & 95 & 68 & 68 & 12 & 5.67 \\
Al Jazeera & Qatar & 3 & 94 & 37 & 22 & 4 & 5.50 \\
The Washington Post & United States & 1 & 18 & 11 & 11 & 2 & 5.50 \\
Clean Energy Wire & Germany & 1 & 6 & 6 & 0 & -- & -- \\
\textit{All English} & & 34 & 686 & 431 & 393 & 60 & \\
\midrule
Total & 6 countries & 74 & 1{,}047 & 696 & 618 & 120 & \\
\bottomrule
\end{tabular}
\begin{tablenotes}\footnotesize
\item Articles: articles with body prose (articles without it are not counted). Candidates: paragraphs the keyword prefilter marks as naming at least one studied technology. $N_h$: candidates in the reserve that carry no transcription gap and do not come from the article quoted in the prompt. Clean Energy Wire forms no stratum: its one article is the one quoted in the prompt. Allocation within language by the largest-remainder method, cap 24 per stratum (binding only where $n_h = 24$).
\end{tablenotes}
\end{threeparttable}
\end{table}


\subsection{Requirements of the Topic Description}\label{app:topic_requirements}

Table~\ref{tab:topic_requirements} traces each requirement of the topic description to its status in this thesis, with the reason and the place where the thesis documents it.

\begingroup\scriptsize
\setlength{\tabcolsep}{4pt}
\begin{longtable}{@{}>{\raggedright\arraybackslash}p{32mm}>{\raggedright\arraybackslash}p{14mm}>{\raggedright\arraybackslash}p{106mm}@{}}
\caption[Requirements of the topic description]{Requirements of the topic description: delivered, changed or not done, with the reason and where the thesis documents it.}\label{tab:topic_requirements} \\
\toprule
\textbf{Requirement} & \textbf{Status} & \textbf{Reason; where} \\
\midrule
\endfirsthead
\toprule
\textbf{Requirement} & \textbf{Status} & \textbf{Reason; where} \\
\midrule
\endhead
\bottomrule
\endfoot
\multicolumn{3}{@{}l}{\emph{Data source and retrieval}} \\
Multilingual retrieval; Factiva preferred, otherwise e.g.\ GDELT & changed & Hand collection by keyword search instead of a database route: 95 articles, 3 August to 2 September 2026 (Table~\ref{tab:corpus_facts}). Kept on the supervisor's advice of 10 September 2026. Factiva was explored but could not be used: the licence excluded automated queries (bulk retrieval, below). GDELT was not used: it may provide only links, headlines or excerpts, as the topic description notes, whereas the coding unit is the body paragraph. \\
Coverage: countries, outlets, languages, years & changed & Set by the hand collection, not by a database: outlet, country, language and publication date are recorded for each article (Appendix~\ref{app:data_dictionary}). \\
Reproducible bulk retrieval & not done & Clarified on 1 October 2026, after the frame freeze: automated queries under the academic licence are not permissible. A comparison of six aggregator databases reports a 100-article download limit for Factiva and Nexis Uni and extra cost for APIs where offered \parencite{gilbert2024newsaggregators}. Instead, the registry's closed corpus versions make the hand-collected corpus reproducible, not its retrieval. \\
Full-text storage & not clarified & Not clarified, for a database or for the texts collected by hand: the study records document no basis for storing them. As a precaution they are kept outside the repository, which withholds article and paragraph text (Appendix~\ref{app:electronic}). \\
Processing by an external LLM & not clarified & No basis for processing the hand-collected texts with an external provider is documented in the study records; the author's routing approvals do not establish copyright or contractual permission. Under the author's own data rule, article text stays on the EU route of Amazon Bedrock, which the registered run used (Table~\ref{tab:model_config}). Later steps outside that route, each authorised by the author, are listed in Section~\ref{subsec:limitations_future}, Tables~\ref{tab:disclosures} and~\ref{tab:deviations} and Item D.1 of the checklist (Appendix~\ref{app:guide_llm}). \\
Redistribution restrictions & not clarified & Complete articles and paragraphs are not shared; the repository holds identifiers, labels and metadata. The frozen codebook files and the prompt carry short quotations from calibration paragraphs (Appendix~\ref{app:instrument}); whether these may be redistributed was not clarified. \\
\multicolumn{3}{@{}l}{\emph{Corpus preparation}} \\
Cleaning & delivered & Body prose kept, with optical character recognition for raster prints; 21 articles without body prose excluded; 26 paragraphs with a marked transcription gap flagged in the registry (Table~\ref{tab:corpus_facts}). \\
Deduplication & delivered & A containment check with eight-word shingles at a threshold of 0.15 guards against a duplicate re-entering under a new identity; its quarantine is empty (Table~\ref{tab:corpus_facts}). \\
Metadata & delivered & Outlet, country, language, publication date, headline, canonical URL and retrieval date for each article (Appendix~\ref{app:data_dictionary}). \\
Segmentation: paragraphs, possibly with headline and limited context & changed & Paragraphs of body prose with at least 15 words: 1,047. The coder saw headline, outlet and language, the model the paragraph alone; giving the model these lines was tested later (Appendix~\ref{app:later_checks}). \\
\multicolumn{3}{@{}l}{\emph{Extraction}} \\
Relevant passages & delivered & A keyword prefilter over the seven technologies marks 696 candidate paragraphs for the frame, and the model can return a paragraph as off topic. Screening was not evaluated, because every sampled paragraph passed the prefilter. \\
Technologies & changed & Seven categories: those of the topic description plus storage (Table~\ref{tab:literature}); the hydrogen production pathway is retained in its own field (Appendix~\ref{app:data_dictionary}). \\
Frames & changed & Nine frames instead of five: the five of the topic description, three added after the first calibration session (technical feasibility, siting, regulation) and a residual, built from 17 elements by a fixed mapping (Table~\ref{tab:codebook_elements}). \\
Evaluation: positive, negative, neutral, mixed or uncertain & changed & Positive, negative, neutral and mixed (Table~\ref{tab:literature}); ``uncertain'' is replaced by the coder's confidence rating, recorded on every row but not analysed. \\
Multilingual LLM extraction; prompt and output structure & delivered & One untranslated instrument for both languages: the frozen codebook rendered as the system prompt, the paragraph as the only user message and one forced tool call with a pinned schema (Appendix~\ref{app:instrument}). \\
Handling of uncertainty & changed & No uncertainty value and no abstention for doubt: off topic and a technology without a frame are negative classifications. Refusals and failed outputs are recorded as such, three executions measure instability, and the coder's confidence rating is recorded as above. \\
\multicolumn{3}{@{}l}{\emph{Validation}} \\
Manually annotated sample of suitable size & delivered & A probability sample of 120 paragraphs from 56 articles coded by one coder (the first coding) and re-coded with the first labels hidden, 40 before the model run (the registered re-code) and 80 after the results were known (together the second coding), and a fresh census of 55 paragraphs (Figure~\ref{fig:sample_flow}); the single coder and the sample size are limits. \\
Across languages and technologies & delivered & Agreement reported per language as well as pooled, and F1 per technology (Table~\ref{tab:per_value}). \\
Across countries and periods & not done & Outlets are nested in countries, and time is not a stratum; the 120 paragraphs give too few per country or period for a figure (8 and 6 from 2024 and 2023; Table~\ref{tab:corpus_facts}). A sample stratified by country or period would allow it. \\
\multicolumn{3}{@{}l}{\emph{Time frame, data structure and scope}} \\
Time frame from 1 January 2023 & changed & Publication dates from 28 February 2023 to 2 September 2026, as found by the purposive hand collection; later months can be added as new corpus versions of the registry (Appendix~\ref{app:data_dictionary}). \\
Structure for added countries and months; aggregation by country, outlet, technology, frame and month & delivered & The records carry every key (Appendix~\ref{app:data_dictionary}). No aggregate by these keys is reported, because the corpus is a purposive collection; the adjusted prevalences describe the frozen frame only (Appendix~\ref{app:dsl_details}). \\
Prototype that can be scaled & delivered & One registered run of 360 executions took 12.5 minutes at an estimated USD 2.18 (Appendix~\ref{app:detailed_results}); scaling is discussed in Section~\ref{subsec:managerial_implications}. \\
Effect of the 2024 United States election & not done & Out of scope, as the topic description states. \\
\end{longtable}
\endgroup

Table~\ref{tab:frame_crosswalk} gives the origin of each frame. The five frames of the topic description were related to the two inventories of Section~\ref{subsec:framing} before calibration. The other three were added after the first calibration session, in which 9 of 20 paragraphs fitted no element; no category-level derivation from the inventories was recorded for them.

\begin{table}[htbp]
\centering\scriptsize
\setlength{\tabcolsep}{4pt}
\caption[Origin of the nine frames]{Origin of the nine frames: the topic description's wording, the context labels of \textcite{toetzke2022energydiscourse} and the dimension of the Policy Frames Codebook \parencite{boydstun2020policyframes} to which each of its five frames was related, and the frames added later.}\label{tab:frame_crosswalk}
\begin{tabular}{@{}>{\raggedright\arraybackslash}p{18mm}>{\raggedright\arraybackslash}p{38mm}>{\raggedright\arraybackslash}p{52mm}>{\raggedright\arraybackslash}p{38mm}@{}}
\toprule
\textbf{Frame} & \textbf{Origin} & \textbf{Context labels of Toetzke et al.} & \textbf{Policy Frames Codebook} \\
\midrule
Costs & topic description: costs and affordability & technology cost and price & Economics \\
Jobs & topic description: jobs, investment, and the economy & employment; private investment; public finance and investment & Economics; Capacity and Resources \\
Climate & topic description: climate and emissions & greenhouse gas emissions; climate change & no counterpart \\
Security & topic description: energy security and independence & energy security; trade & Security and Defense \\
Reliability & topic description: reliability and system integration & technology efficiency; technology characteristics; infrastructure & Capacity and Resources \\
Feasibility, Siting, Regulation & added after the first calibration session (codebook v0.2) & none recorded; the three recur in manual studies (Section~\ref{subsec:stream_framing}) & none recorded \\
Other & residual, with a required note & --- & its residual category \\
\bottomrule
\end{tabular}
\end{table}

\subsection{Output Records and Aggregation Keys}\label{app:data_dictionary}

Table~\ref{tab:data_dictionary} lists the output files and their fields. Joining \texttt{paragraph\_id} and its \texttt{article\_id} with the corpus manifest gives every record its outlet, country, language and publication date. Frames are derived from the element indicators by the fixed element-to-frame mapping, and the month from the publication date.

\begin{table}[htbp]
\centering\scriptsize
\setlength{\tabcolsep}{4pt}
\caption[Output records and aggregation keys]{Output records and their fields: the prediction and terminal files of every run and the corpus manifest, with the unit of one row in parentheses.}\label{tab:data_dictionary}
\begin{tabular}{@{}>{\raggedright\arraybackslash}p{38mm}>{\raggedright\arraybackslash}p{114mm}@{}}
\toprule
\textbf{File (one row per)} & \textbf{Fields} \\
\midrule
Predictions, \modelid{<run_id>_predictions.csv} (technology record of an execution) & \texttt{run\_id}; \texttt{run\_index}; \texttt{paragraph\_id}; \texttt{technology}; \texttt{wind\_siting}; \texttt{hydrogen\_pathway}; \texttt{F1a} \ldots{} \texttt{F9} (the 17 element indicators); \texttt{no\_frame\_for\_this\_tech}; \texttt{eval\_costs} \ldots{} \texttt{eval\_other} (one evaluation value per frame); \texttt{comparison\_target} \\
Terminal outcomes, \modelid{<run_id>_terminal.csv} (execution of a paragraph) & \texttt{run\_id}; \texttt{run\_index}; \texttt{paragraph\_id}; \texttt{lang}; \texttt{outcome}; \texttt{n\_attempts}; \texttt{codes}; \texttt{resolved\_from\_attempt}; \texttt{episode}; \texttt{gate} \\
Corpus manifest, \modelid{manifest.csv} (article) & outlet, country, language, publication date, headline, canonical URL, retrieval date \\
\bottomrule
\end{tabular}
\end{table}

A new month is appended through the registry: its articles enter as a new closed corpus version, with immutable paragraph identities and an append-only, hash-chained assignment log, so earlier records stay unchanged. The frozen frame and samples do not change; prevalences for new months therefore need their own frame and a fresh human-coded sample.

\subsection{Disclosures about the Reference}\label{app:reference_disclosures}

\begingroup\appcaption{Annotation sessions}{}% Generated by tools/annotation/descriptives.py from labels/sessions.csv.
% Do not edit by hand.
\begin{table}[htbp]
\setlength{\tabcolsep}{3pt}\footnotesize
\centering
\small
\begin{threeparttable}
\caption[Annotation sessions]{Annotation sessions. The calibration and development sessions preceded the evaluation sample and none of their labels enters the reported evaluation. The evaluation session (A2) is the reference and A2R its registered blind re-code. A2R2, the second re-code wave of the other 80 paragraphs, followed the registered results. The fresh census's sessions are in Appendix~\ref{app:fresh}.}
\label{tab:sessions}
\begin{tabular}{@{}llllrl@{}}
\toprule
\textbf{Session} & \textbf{Pass} & \textbf{Codebook} & \textbf{Coded} & \textbf{Rows} & \textbf{Unit} \\
\midrule
Calibration I (A0) & 1 & v0.1 & 2026-08-29 & 20 & paragraphs \\
Calibration I (A0) & 2 & v0.3 & 2026-08-29 & 27 & paragraph--technology rows \\
Calibration II (A0b) & 1 & v0.2 & 2026-08-29 & 12 & paragraphs \\
Calibration II (A0b) & 2 & v0.3 & 2026-08-29 & 37 & paragraph--technology rows \\
Development (A1) & 1 & v1.0 & 2026-09-01 & 35 & paragraphs \\
Development (A1) & 2 & v1.0 & 2026-09-01 & 53 & paragraph--technology rows \\
Evaluation (A2) & 1 & v1.0 & 2026-09-28 to 09-30 & 120 & paragraphs \\
Evaluation (A2) & 2 & v1.0 & 2026-09-30 to 10-08 & 178 & paragraph--technology rows \\
Re-code (A2R) & 1 & v1.0 & 2026-10-13 & 40 & paragraphs \\
Re-code (A2R) & 2 & v1.0 & 2026-10-13 & 54 & paragraph--technology rows \\
Second re-code (A2R2) & 1 & v1.0 & 2026-11-10 & 80 & paragraphs \\
Second re-code (A2R2) & 2 & v1.0 & 2026-11-10 & 113 & paragraph--technology rows \\
\bottomrule
\end{tabular}
\begin{tablenotes}\footnotesize
\item The first pass records technologies present; the second expands to one row per paragraph and technology. Both calibration sessions were recoded in the second pass under codebook v0.3 before the freeze as v1.0.
\end{tablenotes}
\end{threeparttable}
\end{table}
\endgroup

\begingroup\scriptsize
\setlength{\tabcolsep}{4pt}
\begin{longtable}{@{}>{\raggedright\arraybackslash}p{36mm}>{\raggedright\arraybackslash}p{64mm}>{\raggedright\arraybackslash}p{52mm}@{}}
\caption[Disclosures about the reference]{Disclosures about the reference: what happened, how it was handled, and what it means for the reported figures. The labels are recorded with the label files.}\label{tab:disclosures} \\
\toprule
\textbf{What happened} & \textbf{How it was handled} & \textbf{Consequence for the figures} \\
\midrule
\endfirsthead
\toprule
\textbf{What happened} & \textbf{How it was handled} & \textbf{Consequence for the figures} \\
\midrule
\endhead
\bottomrule
\endfoot
Discarded first attempt (week before 25 September 2026): the build of the coding interface then in use displayed keyword-derived technology suggestions, the prefilter's lexicon, beside the paragraph & Attempt discarded; no label used or carried over; its export opened by the coder only. Interface rebuilt without any suggestion, highlighting, default or pre-selection; first pass coded anew from blank sheets & The exposure on the paragraphs seen cannot be undone; it could raise technology agreement; treated as a limitation \\
Shortened washout: registered 14 days, so the re-code was admissible from 22 October 2026 & Author's decision of 12 October to recover schedule; re-coded on 13 October, five days after the last first-coding date; the evaluator applied a minimum of four days and records both values & The coder's $\alpha$ may be inflated by recall; limited evidence of temporal stability, not a bound on a 14-day result \\
One exported first-pass row carried a screening value on a paragraph naming no studied technology, where the value is undefined & Cleared before finalisation; correction recorded with the label file & None \\
Six paragraphs screened as framed while every one of their second-pass rows is unframed & The second pass governs frames and evaluations; the screening flag enters no metric; cases not re-read & None \\
The general category admits two readings: per statement (a bundle term beside named technologies yields an extra row) or as a residual & The first coding follows the per-statement reading (40 of 120 paragraphs, 22 of them beside an individual technology); the prompt carries the same definition; codebook unchanged & Under the residual reading the general category's F1 (0.827) and displacement ($-4.7$ points) would shift \\
Fresh census (14 October): after setting first-pass labels, the author put a consistency question about the general category to a language-model assistant outside the pipeline & The assistant read the working sheet and the finalised evaluation labels, reported counts and recommended no label; no label was changed & None on the labels; census text reached the assistant's provider outside the Bedrock EU route \\
Fresh census, second pass (19 October): the export left the comparison target blank on ten rows, all marked as carrying no frame & Set to none by the coding agent on the author's written instruction that the blanks mean none; checked cell by cell against the export, which is kept & None on technologies, elements or evaluations; the census comparison-target figures use the completed cells \\
Re-reading of 45 sampled missed elements (returned 9 November): the coder saw the model's output beside each missed element & The coder judged every item alone, in a separate sheet with fixed categories (\S4l (6)). Afterwards, on the author's instruction, the coding agent set two categories that contradicted the coder's own note to borderline; no reference label changed & None on the labels or any registered figure; the counts are a non-blind judgement (Appendix~\ref{app:later_checks}) \\
A development-set sheet exported by the coding interface with a model's labels in the human slot & Quarantined outside the annotation flow; never entered a label file & None \\
The coding assistant handled exported sheets that carry paragraph text (census second pass, supplementary re-code) and, in the checks of November, local files with paragraph text & What exactly was read, and over which route, was not recorded for the label edits; the assistant's provider lies outside the Bedrock EU route & None on the labels; further exposure of paragraph text to that provider is likely and is not itemised \\
The coding interface is not archived in any build & Its written specification (no suggestion, highlighting, default or pre-selection) is kept with the label files & Its behaviour and data flow cannot be reconstructed independently \\
Supplementary re-code (wave 2, 10 November): coded after the registered results and the re-reading of misses, 25 of whose paragraphs it contains; its declaration (\S4n) and labels entered the repository in one commit (\texttt{6fe59fe}), so their order rests on the specification's record & Declared beforehand (\S4n); the 25 kept by the author's decision and reported with and without; on eight unframed rows the blank comparison target set to none by the agent, as in the first wave; no reference label changed & None on the registered figures; analyses using this re-code are exploratory; it may lean towards the model's reading \\
\end{longtable}
\endgroup

\subsection{Evaluator Conventions}\label{app:evaluator_conventions}

Table~\ref{tab:conventions} lists the conventions the evaluator applies. Macro averages are taken over the classes $C^{+}$ with positive reference support in the population scored, pooled or one language; with design-weighted per-class scores $F_{1,c}^{w}$, the macro-F1 is
\begin{equation}
F_1^{\mathrm{macro}} = \frac{1}{|C^{+}|} \sum_{c \in C^{+}} F_{1,c}^{w}.
\label{eq:macrof1}
\end{equation}
For the registered model comparisons, intervals recompute the statistic $T$ on article resamples drawn with replacement. Article multiplicities $m_a^{(b)}$ multiply the paragraph design weights $w_i$, so that
\begin{equation}
\begin{aligned}
T^{(b)} &= T\bigl(\{\,m_{a(i)}^{(b)}\, w_i\,\}\bigr), \qquad B = 2{,}000, \\
\mathrm{CI}_{95} &= \bigl[\,q_{0.025}(T^{(1)}, \dots, T^{(B)}),\; q_{0.975}(T^{(1)}, \dots, T^{(B)})\,\bigr],
\end{aligned}
\label{eq:resample}
\end{equation}
where $a(i)$ is the article of paragraph $i$ and $q_p$ the $p$-quantile of the $B$ resampled values.

\begingroup\scriptsize
\setlength{\tabcolsep}{4pt}
\begin{longtable}{@{}>{\raggedright\arraybackslash}p{27mm}>{\raggedright\arraybackslash}p{80mm}>{\raggedright\arraybackslash}p{45mm}@{}}
\caption[Evaluator conventions]{Evaluator conventions: the situation, the rule the evaluator applies, and its effect on the reported figures.}\label{tab:conventions} \\
\toprule
\textbf{Situation} & \textbf{Rule} & \textbf{Effect} \\
\midrule
\endfirsthead
\toprule
\textbf{Situation} & \textbf{Rule} & \textbf{Effect} \\
\midrule
\endhead
\bottomrule
\endfoot
Weighting & Model agreement and displacement use paragraph design weights $w_i = 1/\pi_i$; unweighted figures are also retained. The coder's $\alpha$, raw agreement and reported support counts are unweighted. On the 40-paragraph re-code, the common one-third selection probability within each language cancels in the design-weighted ratios & Full-sample weighted estimates target the 618-paragraph eligible frame, not the 1{,}047-paragraph corpus \\
Technologies and elements & Micro-F1 over (paragraph $\times$ technology) units and (paragraph $\times$ technology $\times$ element) tuples on the outer join over every sample paragraph & The headline quantities \\
Matched units (matched-frame and matched-row accuracy; score~A) & Evaluation: accuracy and per-value F1 over (paragraph, technology, frame) units both sides fired; comparison target: over (paragraph, technology) rows both sides named; excluded rows and frames are counted & Conditional on detection; post hoc, because a stateless call cannot be handed the reference frames \\
End to end (end-to-end F1; score~B) & (paragraph, technology, frame, value) and (paragraph, technology, target) tuples compared as sets: a true positive only if identical, a missing row or frame a false negative, an invented one a false positive & A missed frame lowers end-to-end F1 but not matched-frame accuracy \\
Exact match & The whole relation row agrees in every field; point estimate & Secondary \\
Macro averages & Equation~\eqref{eq:macrof1}; class sets frozen once from the final reference labels in a class manifest per dimension, language and table; undefined over fewer than two supported classes. A per-language macro uses that language's own set: for the elements, 16 classes each, 15 of them shared, so the two macros are not on one class set & Unstable over rare classes; the defined-class counts over the resamples are electronic (Appendix~\ref{app:electronic}) \\
Macro inside a resample & A class with zero support and zero predictions drops out (available-class macro-F1, with the count of defined classes); zero support with predictions gives F1 = 0 and stays; false positives always enter the micro counts & The macro interval is over available classes \\
Intervals & Equation~\eqref{eq:resample}: articles resampled within language for a per-language figure, over all sampled articles for a pooled one, never within outlet & Descriptive; no finite-population coverage guarantee is claimed \\
Reportability & Sealed threshold $s = 5$, fixed blind on 7 September 2026 (commit \texttt{7c61b3d}). Support unit: agreement tables, the unweighted reference row (an evaluation value: the frame-bearing reference unit; matched accuracy: the matched unit); displacement, the paragraph carrying the class; reliability, the aligned unit the first coding marks positive (categorical items: the unit inside the applicability mask) & A cell below $s$ prints $n<s$ with its support; the class stays in the macro set; single-article support is marked \\
Seal & The evaluator refuses a threshold that differs from the seal, a seal rewritten after its first commit, or one not committed before the first coded reference row, and records the seal's digest & The threshold could not follow the labels; the digest was also sent to the supervisor on 15 September \\
Undefined values & Per-class F1 with zero support and zero predictions; $\alpha$ when all ratings agree without variation; an empty matched intersection, a macro set without a defined class, a resample without a unit at a level & Printed as undefined, never as a number \\
Key-set check & The (paragraph, technology) keys of the second pass must equal those implied by the first-pass ticks, and every sample paragraph needs a first-pass row & A mismatch refuses the evaluation \\
Execution stability & Pairs of successful executions of a paragraph; paragraph level: the sorted technology set; relation level: element set, evaluation map keyed by frame, comparison target and subtype for each technology in the union over the successful executions, an absent relation counting as a value & Paragraphs with fewer than two successful executions have no pair and are counted \\
Coder's temporal stability & Krippendorff's $\alpha$ on aligned units: the technology tick per paragraph; the element flag on rows both codings created; the evaluation where both fired the frame; the comparison target on all aligned rows; a technology ticked once only is a disagreement, not a missing rating & Non-aligned rows are counted; registered washout 14 days, applied minimum four, realised five (Table~\ref{tab:disclosures}) \\
Displacement & Model minus reference prevalence in percentage points over every sample paragraph, with reference and predicted paragraph counts; occurrences of an element under several technologies are not prevalences & The interval carries no uncertainty from the reference \\
Clean domain & The headline recomputed on the 105 paragraphs from articles no part of which was used in calibration or development & A sensitivity, not a contamination effect \\
\end{longtable}
\endgroup

\subsection{Registration Record and Deviations}\label{app:protocol}

\emph{Registered} means fixed in the internal pre-analysis specification, recorded in the repository before the evaluation sample was coded or predicted; it is not a public or a hypothesis pre-registration. The specification holds the questions, estimands and reporting rules; the metric layer is written out in the pipeline's completion plan, and the frozen frame and sample manifests document the draw. The specification was committed on 5, 6 and 7 September 2026 (commits \texttt{607a89e}, \texttt{5af3778}, \texttt{1870c02}, \texttt{400dbc8}), the completion plan on 6 and 7 September, all before the discarded first coding attempt and before the first label file of 30 September. The threshold seal was committed on 7 September (commit \texttt{7c61b3d}). A comparison of the committed files bounds later changes. Sections 1 and 2 of the completion plan, which hold the research questions, estimands and metric layer, are byte-identical at \texttt{400dbc8}, at the evaluation commit \texttt{460b344} of 13 October and on 15 November 2026. Sections 1 to 3 of the specification differ only in one file path. Its other changes before the run are dated in place, between 10 September and 12 October 2026. The run record's starting state is commit \texttt{07e1f8d}, which finalised the blind re-code three seconds before the run started, and its \texttt{git\_at\_start} field records a clean working tree. The specification is still headed as a draft, and its sign-off block was never filled in. Table~\ref{tab:deviations} lists each departure from the specification with its handling.

\begingroup\scriptsize
\setlength{\tabcolsep}{4pt}
\begin{longtable}{@{}>{\raggedright\arraybackslash}p{16mm}>{\raggedright\arraybackslash}p{54mm}>{\raggedright\arraybackslash}p{42mm}>{\raggedright\arraybackslash}p{38mm}@{}}
\caption[Amendments and deviations]{Amendments of and deviations from the pre-analysis specification, with their handling and consequence for the reading. Amendments are struck in place with a date and a reason.}\label{tab:deviations} \\
\toprule
\textbf{Date} & \textbf{Departure} & \textbf{Handling} & \textbf{Consequence} \\
\midrule
\endfirsthead
\toprule
\textbf{Date} & \textbf{Departure} & \textbf{Handling} & \textbf{Consequence} \\
\midrule
\endhead
\bottomrule
\endfoot
6 and 15 Sep & The production model was to be chosen among four candidates by a development analysis. Claude Sonnet 4.6 was kept by the author's choice on 6 September, and the selection deadline passed on 15 September & Selection not run; the registered model is Claude Sonnet 4.6 & None on the registered run \\
24 Sep & Conditional Opus 4.6 comparison closed by the slippage rule & Reopened on 19 October only as a robustness analysis (\S4f) & Enters no registered figure \\
before 25 Sep & First-pass attempt in a build with technology suggestions & Discarded; interface rebuilt (Table~\ref{tab:disclosures}) & Exposure is a limitation \\
12 Oct & Washout shortened from 14 planned to five realised days & Minimum of four days written into the metrics & Coder's $\alpha$ possibly inflated \\
14 Oct & Development iterations after the registered tables had been inspected, as dated departures from the adaptation firewall & Sealed before their runs; development split only & No revised instrument authorised (Table~\ref{tab:development_review}) \\
19 Oct & Adjusted prevalences, closed on 5 September, reopened at the supervisor's request & Declared exploratory (\S4d) & Not a registered figure \\
19 Oct & Nemotron 3 Super: a completed run on the evaluation sample was extended by an erroneous resume & Replaced by a fresh run; the replaced run scores 0.431 against 0.425 & Descriptive row only \\
30 Oct & Firewall amended to allow the mandatory-form runs on the inspected samples (\S4j) & Dated amendment before any call & Descriptive comparison \\
30 Oct & Battery: the EU route's daily token quota was exhausted; 62\% of executions ran on the global profile, which may process requests outside the EU & Within-route computation fixed after a smoke test had printed interim contrasts, before any variant result was read & Planted-gap contrast on the edge between the two computations (Table~\ref{tab:battery_primary}) \\
30 Oct & Mercury 3 runs and the codebook review's second arm sent paragraphs to the developer's API outside Bedrock EU & Transfer cleared by the author as compliant with data protection & Rows differ from the registered request in more than the model \\
1 Nov & Configuration controls for Sonnet 4.6 and Opus 4.6 added to the author's request & Ratified by the author after the runs & Controls read as added \\
5 Nov & Language comparison run with three executions instead of the one declared on 19 October; extraction began before the translation check was sealed & Amendment before any translation; no agreement figure before the seal & Exploratory; inconclusive for agreement (Appendix~\ref{app:language}) \\
8 Nov & Opus 5.5 rejects the forced tool choice and disabled reasoning; a control, a high-effort run and an Opus 5 repeat added to the author's request & Smallest accepted request declared before any run (\S4m); additions declared as the agent's & Rows differ from the registered request in tool choice and reasoning \\
10 Nov & The paired change of the precision--recall gap and the per-element directions were added after the second-wave results had been inspected (\S4n (3)), then extended to the registered wave & Declared before their computation & Descriptive; marked as later in Section~\ref{subsec:reliability_results} \\
\end{longtable}
\endgroup

\subsection{Detailed Results of the Registered Run}\label{app:detailed_results}

The tables below back the registered figures of run \texttt{a2\_sonnet\_20261013} that Chapter~\ref{sec:results} summarises. The run made 379 attempts for its 360 executions on 13 October 2026, took 12.5 minutes and cost an estimated USD 2.18; its request is the first row of Table~\ref{tab:model_config}. In execution one, 168 of the 376 first-coding element tuples were missed and 68 tuples were invented (unweighted counts). Excluding the 15 paragraphs from articles used in calibration or development changes no pooled headline score by more than 0.01; on the remaining 105 paragraphs from 50 articles, the element micro-F1 is 0.640 [0.575, 0.706]. This clean-domain table, the run record with tokens, gate and environment, and the evaluator's complete tables with the unweighted figures, the comparison-target classes and the classes defined across resamples, are electronic (Appendix~\ref{app:electronic}).

% Generated by tools/pipeline/figures.py tables from a2_sonnet_20261013/metrics.json (sha256 9d69e490475a885f7216ceb5c49d295b9d7e55662b7092ca6faa1636e46c57f4). Do not edit by hand.
\begin{table}[htbp]
\centering
\scriptsize
\setlength{\tabcolsep}{4pt}
\begin{threeparttable}
\caption[Agreement with the first coding by dimension and language, complete]{Complete headline agreement: the micro scores of Table~\ref{tab:primary} with the macro-F1 and the whole-row exact match, by dimension and language. Design-weighted point estimates; brackets: 2.5--97.5 percentile intervals from 2{,}000 article resamples (56 pooled; 28 per language); $n$: paragraphs.}
\label{tab:primary_full}
\begin{tabular}{@{}>{\raggedright\arraybackslash}p{24mm}>{\raggedright\arraybackslash}p{28mm}>{\raggedright\arraybackslash}p{31mm}>{\raggedright\arraybackslash}p{31mm}>{\raggedright\arraybackslash}p{31mm}@{}}
\toprule
\textbf{Dimension} & \textbf{Score} & \textbf{Pooled ($n$\,=\,120)} & \textbf{German ($n$\,=\,60)} & \textbf{English ($n$\,=\,60)} \\
\midrule
Technology detection & micro-F1 & 0.935 [0.901, 0.963] & 0.884 [0.781, 0.952] & 0.960 [0.933, 0.984] \\
Frame elements & micro-F1 & 0.643 [0.581, 0.705] & 0.621 [0.532, 0.710] & 0.653 [0.569, 0.731] \\
Evaluation & matched-frame accuracy & 0.874 [0.812, 0.924] & 0.841 [0.757, 0.916] & 0.890 [0.807, 0.953] \\
 & end-to-end F1 & 0.617 [0.555, 0.679] & 0.597 [0.498, 0.694] & 0.626 [0.542, 0.706] \\
Comparison target & matched-row accuracy & 0.794 [0.697, 0.879] & 0.881 [0.792, 0.951] & 0.755 [0.616, 0.872] \\
 & end-to-end F1 & 0.742 [0.647, 0.831] & 0.779 [0.661, 0.868] & 0.725 [0.590, 0.844] \\
\addlinespace
Frame elements & macro-F1 & 0.615 [0.501, 0.657] (17) & 0.600 [0.481, 0.662] (16) & 0.618 [0.469, 0.665] (16) \\
Evaluation & matched-frame macro-F1 & 0.744 [0.643, 0.830] (4) & 0.733 [0.526, 0.860] (4) & 0.746 [0.579, 0.866] (4) \\
 & end-to-end macro-F1 & 0.523 [0.444, 0.582] (4) & 0.513 [0.358, 0.639] (4) & 0.516 [0.411, 0.591] (4) \\
Comparison target & matched-row macro-F1 & 0.446 [0.352, 0.601] (7) & 0.876 [0.453, 0.986] (3) & 0.475 [0.333, 0.667] (7) \\
 & end-to-end macro-F1 & 0.379 [0.295, 0.507] (8) & 0.490 [0.271, 0.905] (5) & 0.467 [0.323, 0.659] (7) \\
\addlinespace
Whole relation row & exact match & 0.215 (184 rows) & 0.228 (87 rows) & 0.208 (97 rows) \\
\bottomrule
\end{tabular}
\begin{tablenotes}\footnotesize
\item Matched-frame accuracy: the share of (paragraph, technology, frame) units fired by both the first coding and the model on which the evaluation agrees (208 pooled, 94 German, 114 English units). Matched-row accuracy: the share of (paragraph, technology) rows named by both on which the comparison target agrees (158 pooled, 69 German, 89 English rows). End-to-end F1 is the micro-F1 over all (paragraph, technology, frame, evaluation) or (paragraph, technology, comparison target) tuples, so a frame or row that one side missed counts as an error. Micro-F1 pools the tuples of all classes.
\item Macro-F1 is the point estimate over the classes with positive reference support (count after the interval); its interval is the available-class macro over resamples. Per-language macros use each language's own class set and are not direct language contrasts. The exact-match rate carries no interval. Unweighted figures and the per-population defined-class distributions are in the evaluator's output file \texttt{table1\_primary.csv} (electronic appendix).
\end{tablenotes}
\end{threeparttable}
\end{table}


\begingroup\let\footnotesize\scriptsize% Generated by tools/pipeline/figures.py tables from a2_sonnet_20261013/metrics.json (sha256 9d69e490475a885f7216ceb5c49d295b9d7e55662b7092ca6faa1636e46c57f4). Do not edit by hand.
% OUTLINE: The registered re-code reports one coder's temporal stability on aligned units after the realised minimum washout.
\begingroup\scriptsize
\setlength{\tabcolsep}{4pt}
\setlength{\LTpost}{3pt}
\begin{longtable}{@{}>{\raggedright\arraybackslash}p{19mm}lrrr>{\raggedright\arraybackslash}p{15mm}r>{\raggedright\arraybackslash}p{24mm}@{}}
\caption[Temporal stability of the reference]{Intra-coder temporal stability on the registered 40-paragraph re-code (Run \texttt{a2\_sonnet\_20261013}, reportable). Nominal Krippendorff's $\alpha$ and raw agreement compare the first coding with the blind re-code. The table reports aligned units, both codings' supports and article support. The realised minimum interval was 5 days, rather than the registered 14; this deviation is disclosed.}\label{tab:reliability} \\
\toprule
\textbf{Level} & \textbf{Item} & \textbf{$\alpha$} & \textbf{Raw} & \textbf{$n$} & \textbf{Support 1st / 2nd} & \textbf{Articles} & \textbf{Mask} \\
\midrule
\endfirsthead
\toprule
\textbf{Level} & \textbf{Item} & \textbf{$\alpha$} & \textbf{Raw} & \textbf{$n$} & \textbf{Support 1st / 2nd} & \textbf{Articles} & \textbf{Mask} \\
\midrule
\endhead
\bottomrule
\endfoot
\bottomrule
\endlastfoot
Technology & general & 0.779 & 0.925 & 40 & 8 / 9 & 7 &  \\
Technology & solar & 1.000 & 1.000 & 40 & 11 / 11 & 9 &  \\
Technology & wind & 1.000 & 1.000 & 40 & 9 / 9 & 7 &  \\
Technology & hydrogen & 1.000 & 1.000 & 40 & 10 / 10 & 5 &  \\
Technology & other\_renewable & $n<5$ (1) & $n<5$ (1) & 40 & 1 / 1 & 1 &  \\
Technology & storage & $n<5$ (3) & $n<5$ (3) & 40 & 3 / 5 & 3 &  \\
Technology & nuclear & 1.000 & 1.000 & 40 & 9 / 9 & 6 &  \\
\addlinespace
Element & F1a & 0.678 & 0.900 & 50 & 10 / 9 & 7 & one coding only: 1 + 4 \\
Element & F1b & 0.899 & 0.980 & 50 & 6 / 5 & 5 &  \\
Element & F2a & $n<5$ (0) & $n<5$ (0) & 50 & 0 / 1 & 0 &  \\
Element & F2b & 0.788 & 0.900 & 50 & 20 / 17 & 14 &  \\
Element & F2c & 0.369 & 0.820 & 50 & 8 / 9 & 7 &  \\
Element & F3a & $n<5$ (4) & $n<5$ (4) & 50 & 4 / 4 & 4 &  \\
Element & F3b & 0.562 & 0.900 & 50 & 8 / 5 & 7 &  \\
Element & F4a & $n<5$ (3) & $n<5$ (3) & 50 & 3 / 2 & 3 &  \\
Element & F4b & $n<5$ (2) & $n<5$ (2) & 50 & 2 / 3 & 1 &  \\
Element & F5a & $n<5$ (4) & $n<5$ (4) & 50 & 4 / 6 & 3 &  \\
Element & F5b & $n<5$ (3) & $n<5$ (3) & 50 & 3 / 2 & 3 &  \\
Element & F6a & 0.625 & 0.920 & 50 & 8 / 4 & 6 &  \\
Element & F6b & 0.462 & 0.920 & 50 & 5 / 3 & 2 &  \\
Element & F7a & $n<5$ (2) & $n<5$ (2) & 50 & 2 / 3 & 1 &  \\
Element & F7b & $n<5$ (2) & $n<5$ (2) & 50 & 2 / 1 & 2 &  \\
Element & F8a & 0.598 & 0.880 & 50 & 12 / 6 & 9 &  \\
Element & F9 & $n<5$ (0) & $n<5$ (0) & 50 & 0 / 0 & 0 &  \\
\addlinespace
Evaluation & costs & 1.000 & 1.000 & 12 & 12 & 10 & 12 of 17 \\
Evaluation & jobs & 0.787 & 0.895 & 19 & 19 & 13 & 19 of 29 \\
Evaluation & climate & 0.492 & 0.750 & 8 & 8 & 8 & 8 of 12 \\
Evaluation & security & $n<5$ (4) & $n<5$ (4) & 4 & 4 & 4 & 4 of 5 \\
Evaluation & reliability & $n<5$ (4) & $n<5$ (4) & 4 & 4 & 3 & 4 of 11 \\
Evaluation & feasibility & 0.755 & 0.857 & 7 & 7 & 4 & 7 of 13 \\
Evaluation & siting & $n<5$ (3) & $n<5$ (3) & 3 & 3 & 2 & 3 of 5 \\
Evaluation & regulation & 0.756 & 0.833 & 6 & 6 & 5 & 6 of 12 \\
Evaluation & other & $n<5$ (0) & $n<5$ (0) & 0 & 0 & 0 & 0 of 0 \\
\addlinespace
Comparison target & all aligned rows & 0.624 & 0.880 & 50 & 50 & 27 &  \\
\end{longtable}
% OUTLINE: The note defines support suppression, undefined alpha and the applicable units for each label dimension.
\par\noindent{\footnotesize An item whose first-coding support is below 5 prints its state; $\alpha$ is undefined when all ratings across both codings are identical. For technology ticks, the unit is the paragraph. For elements, it is the aligned (paragraph $\times$ technology) row (50 rows). For evaluations, it is the aligned row on which both codings fired the frame; the mask column gives these rows of the rows on which either fired. For the comparison target, it is every aligned row.\par}
\par\bigskip
\endgroup
\endgroup

% Generated by tools/pipeline/error_breakdown.py (sha256 6c95b39e7a3b7aaa) from run a2_sonnet_20261013: metrics.json sha256 9d69e490475a885f7216ceb5c49d295b9d7e55662b7092ca6faa1636e46c57f4; declared in the pre-analysis specification §4c (2026-10-19), descriptive; the per-element F1 reproduces the registered per-class F1. Do not edit by hand.
\begin{table}[htbp]
\centering
\scriptsize
\setlength{\tabcolsep}{4pt}
\begin{threeparttable}
\caption[Per-element precision, recall and F1]{Per-element precision, recall and F1 of the first execution against the first coding (Run \texttt{a2\_sonnet\_20261013}, reportable; descriptive decomposition declared 2026-10-19): design-weighted, over the same tuples as the registered per-class F1, which the F1 column reproduces; reference support in rows in parentheses, cells with fewer than 5 rows print their state.}
\label{tab:element_prf}
\begin{tabular}{@{}l>{\raggedright\arraybackslash}p{40mm}>{\raggedright\arraybackslash}p{40mm}>{\raggedright\arraybackslash}p{40mm}@{}}
\toprule
\textbf{Element} & \textbf{Pooled P / R / F1 ($n$)} & \textbf{German} & \textbf{English} \\
\midrule
F1a & 0.883 / 0.643 / 0.744 (45) & 0.848 / 0.647 / 0.734 (18) & 0.897 / 0.642 / 0.748 (27) \\
F1b & 0.817 / 0.674 / 0.739 (30) & 0.814 / 0.697 / 0.751 (13) & 0.818 / 0.663 / 0.733 (17) \\
F2a & $n<5$ (3) & $n<5$ (3) & $n<5$ (0) \\
F2b & 0.976 / 0.495 / 0.657 (60) & 0.857 / 0.328 / 0.475 (19) & 1.000 / 0.543 / 0.704 (41) \\
F2c & 0.617 / 0.613 / 0.615 (35) & 0.555 / 0.639 / 0.594 (17) & 0.665 / 0.597 / 0.629 (18) \\
F3a & 0.664 / 0.481 / 0.558 (13) & 0.735 / 0.426 / 0.540 (7) & 0.628 / 0.520 / 0.569 (6) \\
F3b & 0.599 / 0.486 / 0.536 (38) & 0.589 / 0.618 / 0.604 (13) & 0.603 / 0.444 / 0.511 (25) \\
F4a & 1.000 / 0.554 / 0.713 (11) & 1.000 / 0.616 / 0.763 (5) & 1.000 / 0.523 / 0.687 (6) \\
F4b & 0.699 / 0.761 / 0.729 (13) & $n<5$ (3) & 0.657 / 0.724 / 0.689 (10) \\
F5a & 0.789 / 0.875 / 0.830 (20) & 0.576 / 0.767 / 0.658 (9) & 0.919 / 0.926 / 0.923 (11) \\
F5b & 0.796 / 0.689 / 0.739 (17) & 0.749 / 0.775 / 0.762 (13) & $n<5$ (4) \\
F6a & 0.860 / 0.623 / 0.723 (20) & 0.806 / 0.493 / 0.611 (8) & 0.876 / 0.671 / 0.760 (12) \\
F6b & 0.391 / 0.485 / 0.433 (15) & 0.833 / 0.460 / 0.593 (13) & $n<5$ (2) \\
F7a & 0.699 / 0.219 / 0.334 (14) & $n<5$ (2) & 1.000 / 0.243 / 0.392 (12) \\
F7b & 1.000 / 0.320 / 0.484 (12) & $n<5$ (4) & 1.000 / 0.342 / 0.510 (8) \\
F8a & 0.788 / 0.320 / 0.455 (28) & 0.859 / 0.361 / 0.508 (17) & 0.720 / 0.283 / 0.406 (11) \\
F9 & $n<5$ (2) & $n<5$ (0) & $n<5$ (2) \\
\bottomrule
\end{tabular}
\begin{tablenotes}\footnotesize
\item Precision is the share of predicted element tuples that the reference carries, recall the share of reference tuples the model predicted; a reference row without a predicted row counts every element as a false negative, a predicted row on a reference-off-topic paragraph as a false positive. No resampling interval (per-class cells carry none in the registered design).
\end{tablenotes}
\end{threeparttable}
\end{table}


\begin{table}[htbp]
\centering\scriptsize
\setlength{\tabcolsep}{4pt}
\begin{threeparttable}
\caption[Per-class F1 of technologies and evaluation values]{Per-class F1 of the technologies and the evaluation values, first execution against the first coding (Run \texttt{a2\_sonnet\_20261013}): design-weighted, reference support in parentheses; $n<5$: fewer than five supporting units. The per-element F1 is in Table~\ref{tab:element_prf}.}\label{tab:per_value}
\begin{tabular}{@{}>{\raggedright\arraybackslash}p{58mm}>{\raggedright\arraybackslash}p{28mm}>{\raggedright\arraybackslash}p{28mm}>{\raggedright\arraybackslash}p{28mm}@{}}
\toprule
\textbf{Class} & \textbf{Pooled} & \textbf{German} & \textbf{English} \\
\midrule
\multicolumn{4}{@{}l}{\emph{Technology (support: reference rows)}} \\
general & 0.827 (40) & 0.750 (17) & 0.862 (23) \\
solar & 0.930 (37) & 0.847 (23) & 1.000 (14) \\
wind & 0.988 (36) & 0.977 (24) & 1.000 (12) \\
hydrogen & 1.000 (20) & $n<5$ (3) & 1.000 (17) \\
other\_renewable & 0.820 (10) & 0.682 (6) & $n<5$ (4) \\
storage & 0.957 (13) & 1.000 (8) & 0.918 (5) \\
nuclear & 1.000 (22) & $n<5$ (3) & 1.000 (19) \\
\addlinespace
\multicolumn{4}{@{}l}{\emph{Evaluation value on matched frames (support: matched units)}} \\
positive & 0.916 (126) & 0.857 (48) & 0.934 (78) \\
negative & 0.933 (62) & 0.919 (37) & 0.946 (25) \\
neutral & 0.558 (11) & 0.588 (6) & 0.530 (5) \\
mixed & 0.571 (9) & $n<5$ (3) & 0.571 (6) \\
\addlinespace
\multicolumn{4}{@{}l}{\emph{Evaluation value end to end (support: frame-bearing reference units)}} \\
positive & 0.703 (189) & 0.648 (69) & 0.720 (120) \\
negative & 0.560 (118) & 0.612 (60) & 0.523 (58) \\
neutral & 0.370 (19) & 0.436 (10) & 0.320 (9) \\
mixed & 0.460 (14) & 0.356 (6) & 0.500 (8) \\
\bottomrule
\end{tabular}
\begin{tablenotes}\footnotesize
\item The matched evaluation units exclude 132 frame-bearing reference units the model did not predict, which with the 208 matched units make 340. Of the 132, 27 are carried by 20 missing technology rows and 105 are missing frames within matched rows (German 18 and 33, English 9 and 72). Per-class cells carry no resampling interval.
\end{tablenotes}
\end{threeparttable}
\end{table}

% Generated by tools/pipeline/figures.py tables from a2_sonnet_20261013/metrics.json (sha256 9d69e490475a885f7216ceb5c49d295b9d7e55662b7092ca6faa1636e46c57f4). Do not edit by hand.
\begin{table}[htbp]
\centering
\footnotesize
\setlength{\tabcolsep}{4pt}
\begin{threeparttable}
\caption[Terminal outcomes per execution]{Terminal outcome of every paragraph per execution (Run \texttt{a2\_sonnet\_20261013}, reportable). The first five outcome columns partition the population; the last column is the registered joint count of paragraphs on which both the first coding and the model name at least one technology, a subset of the first.}
\label{tab:attrition}
\begin{tabular}{@{}lr>{\raggedright\arraybackslash}p{20mm}>{\raggedright\arraybackslash}p{12mm}>{\raggedright\arraybackslash}p{12mm}>{\raggedright\arraybackslash}p{12mm}>{\raggedright\arraybackslash}p{10mm}>{\raggedright\arraybackslash}p{20mm}@{}}
\toprule
\textbf{Population} & \textbf{Exec.} & \textbf{Technologies present} & \textbf{Off topic} & \textbf{Parse failed} & \textbf{Refused} & \textbf{Error} & \textbf{Both sides present} \\
\midrule
Pooled & 1 & 112 & 3 & 5 & 0 & 0 & 111 \\
 & 2 & 113 & 2 & 5 & 0 & 0 & 112 \\
 & 3 & 114 & 3 & 3 & 0 & 0 & 113 \\
\addlinespace
German & 1 & 52 & 3 & 5 & 0 & 0 & 51 \\
 & 2 & 55 & 2 & 3 & 0 & 0 & 54 \\
 & 3 & 54 & 3 & 3 & 0 & 0 & 53 \\
\addlinespace
English & 1 & 60 & 0 & 0 & 0 & 0 & 60 \\
 & 2 & 58 & 0 & 2 & 0 & 0 & 58 \\
 & 3 & 60 & 0 & 0 & 0 & 0 & 60 \\
\bottomrule
\end{tabular}
\begin{tablenotes}\footnotesize
\item The first coding marks 1 paragraph as naming no studied technology (German). A parse failure is a response that remained schema- or contract-invalid after the one identical retry; it enters the agreement figures as omissions.
\end{tablenotes}
\end{threeparttable}
\end{table}


The change rates of Table~\ref{tab:stability} describe the instrument; they are not evidence about language.

\appinput{tables/tab_stability.tex}

% Generated by tools/pipeline/error_breakdown.py (sha256 6c95b39e7a3b7aaa) from run a2_sonnet_20261013: metrics.json sha256 9d69e490475a885f7216ceb5c49d295b9d7e55662b7092ca6faa1636e46c57f4; declared in the pre-analysis specification §4c (2026-10-19), descriptive; the per-element F1 reproduces the registered per-class F1. Do not edit by hand.
\begin{table}[htbp]
\centering
\scriptsize
\setlength{\tabcolsep}{4pt}
\begin{threeparttable}
\caption[Per-frame detection and evaluation, conditional and end to end]{Per-frame decomposition of the first execution against the first coding (Run \texttt{a2\_sonnet\_20261013}, reportable; descriptive decomposition declared 2026-10-19): precision, recall and F1 of frame detection at the (paragraph $\times$ technology $\times$ frame) level, the evaluation accuracy on the frames the model found (score A, with the matched units) and the end-to-end evaluation F1 (score B, with the reference support); design-weighted; a frame with fewer than 5 reference units in the population prints its state, and the score-A cell its own state when fewer than 5 units were matched.}
\label{tab:frame_prf}
\begin{tabular}{@{}lrrrrr@{}}
\toprule
\textbf{Frame} & \textbf{Detection P} & \textbf{R} & \textbf{F1 ($n$)} & \textbf{Evaluation, score A (units)} & \textbf{Evaluation, score B F1 ($n$)} \\
\midrule
\multicolumn{6}{@{}l}{\emph{Pooled}} \\
costs & 0.887 & 0.671 & 0.764 (70) & 0.879 (47) & 0.671 (70) \\
jobs & 0.934 & 0.618 & 0.744 (82) & 0.859 (53) & 0.639 (82) \\
climate & 0.677 & 0.530 & 0.595 (49) & 0.851 (26) & 0.506 (49) \\
security & 0.825 & 0.741 & 0.781 (20) & 1.000 (15) & 0.781 (20) \\
reliability & 0.904 & 0.852 & 0.877 (37) & 0.900 (31) & 0.789 (37) \\
feasibility & 0.760 & 0.636 & 0.692 (33) & 0.724 (20) & 0.501 (33) \\
siting & 0.814 & 0.301 & 0.440 (19) & 1.000 (6) & 0.440 (19) \\
regulation & 0.788 & 0.320 & 0.455 (28) & 0.912 (9) & 0.415 (28) \\
other & $n<5$ (2) & & & $n<5$ (1) & \\
\addlinespace
\multicolumn{6}{@{}l}{\emph{German}} \\
costs & 0.859 & 0.694 & 0.768 (27) & 0.943 (19) & 0.724 (27) \\
jobs & 0.834 & 0.746 & 0.787 (33) & 0.806 (25) & 0.634 (33) \\
climate & 0.643 & 0.557 & 0.597 (18) & 0.810 (10) & 0.484 (18) \\
security & 1.000 & 1.000 & 1.000 (5) & 1.000 (5) & 1.000 (5) \\
reliability & 0.806 & 0.812 & 0.809 (22) & 0.774 (18) & 0.626 (22) \\
feasibility & 0.822 & 0.526 & 0.641 (19) & 0.797 (10) & 0.511 (19) \\
siting & $n<5$ (4) & & & $n<5$ (1) & \\
regulation & 0.859 & 0.361 & 0.508 (17) & 0.835 (6) & 0.424 (17) \\
other & $n<5$ (0) & & & $n<5$ (0) & \\
\addlinespace
\multicolumn{6}{@{}l}{\emph{English}} \\
costs & 0.898 & 0.663 & 0.763 (43) & 0.854 (28) & 0.652 (43) \\
jobs & 1.000 & 0.565 & 0.722 (49) & 0.888 (28) & 0.641 (49) \\
climate & 0.691 & 0.521 & 0.594 (31) & 0.866 (16) & 0.514 (31) \\
security & 0.788 & 0.694 & 0.738 (15) & 1.000 (10) & 0.738 (15) \\
reliability & 1.000 & 0.886 & 0.939 (15) & 1.000 (13) & 0.939 (15) \\
feasibility & 0.730 & 0.717 & 0.724 (14) & 0.685 (10) & 0.496 (14) \\
siting & 1.000 & 0.311 & 0.474 (15) & 1.000 (5) & 0.474 (15) \\
regulation & 0.720 & 0.283 & 0.406 (11) & $n<5$ (3) & 0.406 (11) \\
other & $n<5$ (2) & & & $n<5$ (1) & \\
\bottomrule
\end{tabular}
\begin{tablenotes}\footnotesize
\item A frame is present on a side when at least one of its elements is; detection recall is therefore the share of reference frames the model found. Score A conditions on those found frames and asks whether the evaluation value agrees; score B counts a (paragraph, technology, frame, value) tuple as correct only when the frame was found and the value agrees, so the gap between the two columns is the agreement lost to frames the model did not find or invented. No resampling interval.
\end{tablenotes}
\end{threeparttable}
\end{table}


\appinput{tables/tab_displacement.tex}

\subsection{Adjusted Prevalences: Estimator and Standard Error}\label{app:dsl_details}

The design-based supervised learning estimator \parencite{egami2023dsl} is used here with the model's first-execution indicator as a fixed prediction, in place of the cross-fitted outcome model of the original estimator. Treating the prediction as fixed assumes that it was generated independently of whether a paragraph was sampled. It estimates the share of the $N = 618$ eligible-frame paragraphs on which a class $c$ is present as
\begin{equation}
\hat{\theta}_{\mathrm{DSL}}(c) = \frac{1}{N} \sum_{i=1}^{N} \Bigl[\, \hat{y}_i(c) + \frac{S_i}{\pi_i}\,\bigl(y_i(c) - \hat{y}_i(c)\bigr) \Bigr],
\label{eq:dsl_estimator}
\end{equation}
with $\hat{y}_i$ the model's first-execution indicator, 0 after a failed execution, $y_i$ the reference indicator, $S_i = 1$ for the 120 sampled paragraphs and $S_i = 0$ otherwise, and $\pi_i$ their registered inclusion probability. The second term is the bias-correction term of \textcite{egami2023dsl}, a design-weighted counterpart of the rectifier in prediction-powered inference \parencite[p.~669]{angelopoulos2023ppi}. Because the eligible frame is a census for the model's indicator, the only design variance under that assumption is that of this term. Under the registered stratified draw its standard error is
\begin{equation}
\widehat{\mathrm{SE}}\bigl(\hat{\theta}_{\mathrm{DSL}}\bigr) = \frac{1}{N} \sqrt{\sum_{g} \bigl(1 - f_g\bigr) \frac{n_g}{n_g - 1} \sum_{i \in g} \bigl(v_i - \bar{v}_g\bigr)^2 }, \qquad v_i = w_i \bigl(y_i(c) - \hat{y}_i(c)\bigr),
\label{eq:dsl_se}
\end{equation}
with $w_i = 1/\pi_i$, $\bar{v}_g$ the mean of $v_i$ in variance group $g$, $f_g = n_g/N_g$ for a group that is a single stratum and $f_g = 0$ for a pooled group. The eighteen strata form sixteen groups: the three single-draw strata are pooled with a neighbouring stratum of their language in allocation order (Berliner Morgenpost with DER SPIEGEL; STERN.de with Salzburger Nachrichten). A pooled group keeps the with-replacement form, which overstates its variance in expectation; pooling strata carries no such guarantee. The design-based interval is the estimate $\pm 1.96\,\widehat{\mathrm{SE}}$, so a class's standard error is the half-width of its design interval in Table~\ref{tab:dsl} divided by 1.96 (for investment, F2b, 3.4 points). It is derived for this draw, treats the predictor as fixed and has no guaranteed nominal coverage. The predictions come from three runs of the unchanged instrument, and the run followed the sample status. The registered run of 13 October 2026 covers the 120 sampled paragraphs and the fresh-census run of 19 October its 55. One execution on the same day covers the other 443 eligible-frame paragraphs (479 calls, USD 2.76). First executions without a valid record number 5 of 120 (4.2\%), 8 of 55 (14.5\%) and 30 of 443 (6.8\%), and count as zero indicators. The assumption of a prediction generated independently of sampling is therefore not verified: run date and paragraph composition differ together, and the residual is estimated on the registered run alone. If failures were more frequent outside the sample for reasons of the run, not of the paragraphs, the adjusted shares would lean downward. A class without any sampled coder--model difference has a zero-width interval, a plug-in artefact that says nothing about unsampled paragraphs. Table~\ref{tab:dsl} prints every class that reaches the reporting threshold. Of the 35 frame--evaluation pairs that either side marks anywhere in the eligible frame, nineteen fall below it; on the evaluation sample alone, either side marks 31 pairs, of which 15 fall below it (Section~\ref{subsec:displacement_results}). These pairs and two elements are in the electronic table.

Under the first coding, the coded sample alone puts investment at 40.1\%, close to the adjusted 37.7\% (Table~\ref{tab:dsl}). The model's labels lower the standard error from 4.3 to 3.4 percentage points. For hydrogen and nuclear, no sampled paragraph shows a coder--model difference, so their design-based intervals have zero width, an artefact of the variance estimate rather than precision.

\begingroup\let\footnotesize\scriptsize\renewcommand{\arraystretch}{0.9}
\appLTcaption{Adjusted prevalences of every reportable class}{Adjusted (DSL) prevalences of every reportable class in the frozen eligible frame (618 paragraphs), in \% of its paragraphs, beside the model-only share, the sample-only estimate of the first coding and the registered displacement.}%
\appgrab{tables/tab_dsl.tex}%
\appthousands%
\appreplace{is\ not\ reported;\ its\ support\ is\ in\ the\ last\ column\.}{is\ not\ printed;\ its\ estimate\ and\ support\ are\ in\ the\ electronic\ table\ (Appendix\c{nobreakspace}\c{ref}\cB\{app:electronic\cE\}).}%
\appdroprows{not reported}%
\endgroup

\subsection{Fresh Census}\label{app:fresh}

The census holds the 55 eligible-frame paragraphs of ten articles that no study record named when it was declared on 14 October 2026 (Section~\ref{subsec:added_methods}). The author coded it in two passes, on 14 and 19 October 2026, blind to the model's output on these paragraphs but after the registered result of 13 October was known; the census is therefore fresh in its articles, not in the coder's knowledge. The registered instrument ran on it on 19 October with three executions (run \texttt{a3\_sonnet\_20261019}, USD 1.09). Its figures are unweighted and describe the census itself; their intervals over ten articles are indicative, and the reference's reliability on the census is unmeasured. The registered instrument's agreement on the census is the v1.0 column of Table~\ref{tab:a3_candidate}; agreement by dimension and language is electronic.

\begingroup\appcaption{Retry candidate on the fresh census}{The registered instrument v1.0 against the frozen candidate v1.0r3, the registered request whose one retry returns the validator's messages to the model, on the fresh census (55 paragraphs, three executions, unweighted). Decided on 26 October 2026 after the v1.0 figure; descriptive only. Brackets: percentile intervals over 2{,}000 resamples of the ten articles; differences paired.}% Generated by tools/pipeline/compare_a3_candidate.py (sha256 1c8e779bed648fb8) from runs a3_sonnet_20261019 and a3_sonnet_r3_20261026 and their evaluator metrics; candidate comparison of specification §4b (author's decision 2026-10-26); descriptive. Do not edit by hand.
\begin{table}[htbp]
\centering
\scriptsize
\setlength{\tabcolsep}{4pt}
\begin{threeparttable}
\caption[Retry candidate on the fresh census]{Retry candidate on the fresh census (specification §4b): the registered instrument v1.0 and the frozen candidate v1.0r3, the registered request byte for byte under retry policy v3, in which the one retry after a schema- or contract-invalid response returns the model's own call with the validator's messages instead of resending the identical request. Both runs: 55 paragraphs, three executions, weight one. A descriptive comparison of complete configurations, decided by the author on 26 October 2026 after the v1.0 figure had been inspected; never a confirmed improvement. Brackets are the 2.5--97.5 percentile over 2{,}000 resamples of the ten census articles, the paired difference under common resamples. Runs \texttt{a3\_sonnet\_20261019} and \texttt{a3\_sonnet\_r3\_20261026}.}
\label{tab:a3_candidate}
\begin{tabular}{@{}>{\raggedright\arraybackslash}p{52mm}>{\raggedright\arraybackslash}p{44mm}>{\raggedright\arraybackslash}p{44mm}@{}}
\toprule
\textbf{Quantity} & \textbf{v1.0 (identical retry)} & \textbf{v1.0r3 (feedback retry)} \\
\midrule
\multicolumn{3}{@{}l}{\emph{Attrition over 165 executions}} \\
Invalid first responses & 24 & 20 \\
Recovered by the retry & 1 & 20 \\
Terminal parse failures & 23 (12 de, 11 en) & 0 (0 de, 0 en) \\
Paragraphs without a valid first execution & 8 & 0 \\
\addlinespace
\multicolumn{3}{@{}l}{\emph{Agreement with the reference, first execution (weight one; intervals indicative)}} \\
Element micro-F1, pooled & 0.563 [0.511, 0.622] & 0.662 [0.589, 0.737] \\
Element micro-F1, German / English & 0.556 / 0.574 & 0.658 / 0.667 \\
Element precision / recall & 0.657 / 0.493 & 0.707 / 0.621 \\
Technology micro-F1 & 0.865 [0.788, 0.927] & 0.951 [0.903, 0.986] \\
Evaluation, score A (units) / score B & 0.742 (66) / 0.447 & 0.704 (81) / 0.487 \\
Change rate over executions, any dimension & 0.099 & 0.145 \\
\addlinespace
\multicolumn{3}{@{}l}{\emph{Paired difference, candidate minus registered, common article resamples}} \\
Element micro-F1 & \multicolumn{2}{l}{+0.098 [0.033, 0.174]} \\
Technology micro-F1 & \multicolumn{2}{l}{+0.086 [0.027, 0.158]} \\
\addlinespace
Estimated cost, USD & 1.09 & 1.07 \\
\bottomrule
\end{tabular}
\begin{tablenotes}\footnotesize
\item Invalid first responses are executions whose first attempt returned a schema- or contract-invalid call; recovered means the retry returned a valid record; a terminal parse failure contributes omissions only to the agreement figures. The reference's reliability on the census is unmeasured; the per-element figures are in the evaluator's tables.
\end{tablenotes}
\end{threeparttable}
\end{table}
\endgroup

\begingroup\appcaption{Prevalence displacement on the fresh census}{}\appinput{tables/fresh_sample/tab_displacement.tex}\endgroup

\subsection{Model Comparison}\label{app:model_comparison}

The comparisons report each configuration's first execution on both samples; repeated executions are identified separately. No model was selected. Gaps include invalid outputs counted as omissions: Mercury~3 instant failed on 27 of 120 first executions on the evaluation sample, and its gap is $+0.245$ on the 93 valid ones against $+0.396$ end to end (Table~\ref{tab:checks_4o}). In the registered request, output is capped at 1{,}024 tokens, with no extended reasoning, and a one-hour cache point follows the constant prefix. Table~\ref{tab:model_config} lists how each request differed from the registered one, and Table~\ref{tab:pr_gaps} gives the precision--recall gap of every configuration with its interval. Figure~\ref{fig:model_comparison} plots each configuration's precision and recall; the complete results table with the configuration controls is electronic (Appendix~\ref{app:electronic}).

\begingroup\scriptsize
\setlength{\tabcolsep}{4pt}
\begin{longtable}{@{}>{\raggedright\arraybackslash}p{27mm}>{\raggedright\arraybackslash}p{41mm}>{\raggedright\arraybackslash}p{26mm}>{\raggedright\arraybackslash}p{40mm}>{\raggedright\arraybackslash}p{14mm}@{}}
\caption[Model-comparison configurations]{Configurations of the model comparison. Every request carries the pinned prompt and tool schema, and every request outside \S4m the forced tool choice; the output cap is 1,024 tokens unless stated. Route: Amazon Bedrock from eu-central-1, through an EU inference profile (\texttt{eu.} prefix) or in-region; date: the runs on both samples.}\label{tab:model_config} \\
\toprule
\textbf{Configuration} & \textbf{Model identifier} & \textbf{Route} & \textbf{Difference from the registered request} & \textbf{Date; \S} \\
\midrule
\endfirsthead
\toprule
\textbf{Configuration} & \textbf{Model identifier} & \textbf{Route} & \textbf{Difference from the registered request} & \textbf{Date; \S} \\
\midrule
\endhead
\bottomrule
\endfoot
Claude Sonnet 4.6 (registered) & \modelid{eu.anthropic.claude-sonnet-4-6} & Bedrock, EU profile & --- (temperature 0, cache point, extended reasoning off) & 13 Oct (census 19 Oct); registered \\
Claude Opus 4.6 & \modelid{eu.anthropic.claude-opus-4-6-v1} & Bedrock, EU profile & none & 19 Oct; \S4f \\
Claude Sonnet 4.5 & \modelid{eu.anthropic.claude-sonnet-4-5-20250929-v1:0} & Bedrock, EU profile & none & 19 Oct; \S4f \\
Claude Haiku 4.5 & \modelid{eu.anthropic.claude-haiku-4-5-20251001-v1:0} & Bedrock, EU profile & none & 19 Oct; \S4f \\
Amazon Nova Pro & \modelid{eu.amazon.nova-pro-v1:0} & Bedrock, EU profile & no cache point (not supported) & 19 Oct; \S4f \\
NVIDIA Nemotron 3 Super & \modelid{nvidia.nemotron-super-3-120b} & Bedrock, in-region & no cache point & 19 Oct; \S4f \\
Qwen3-235B-A22B & \modelid{qwen.qwen3-235b-a22b-2507-v1:0} & Bedrock, in-region & no cache point & 19 Oct; \S4f \\
Mercury 3, instant & \modelid{mercury-3} & developer's API, hosting region not stated & no cache point; temperature 0.5 (the API's minimum); no reasoning; request translated by an adapter & 30 Oct; \S4g \\
Mercury 3, low & \modelid{mercury-3} & as above & as above, low reasoning, 8,192 completion tokens & 30 Oct; \S4g \\
Claude Sonnet 5 & \modelid{eu.anthropic.claude-sonnet-5} & Bedrock, EU profile & temperature omitted (rejects 0; provider's default sampling); extended thinking disabled & 1 Nov; \S4k \\
Claude Opus 5 & \modelid{eu.anthropic.claude-opus-5} & Bedrock, EU profile & as Sonnet 5 & 1 Nov; \S4k \\
Controls: Sonnet 4.6, Opus 4.6 & as above & Bedrock, EU profile & as Sonnet 5 & 1 Nov; \S4k \\
Claude Opus 5.5 & \modelid{eu.anthropic.claude-opus-5-5} & Bedrock, EU profile & tool chosen freely (rejects a forced choice); adaptive thinking at effort low (rejects disabled thinking); temperature omitted; 8,192 tokens; three executions & 8 Nov; \S4m \\
Opus 5.5, effort high & as above & Bedrock, EU profile & as Opus 5.5, effort high; one execution & 8 Nov; \S4m \\
Control: Opus 5, 5.5 request & \modelid{eu.anthropic.claude-opus-5} & Bedrock, EU profile & as Opus 5.5; one execution & 8 Nov; \S4m \\
Opus 5, repeated & \modelid{eu.anthropic.claude-opus-5} & Bedrock, EU profile & as Sonnet 5; three executions & 8 Nov; \S4m \\
\end{longtable}
\endgroup

\begin{table}[htbp]
\centering\footnotesize
\setlength{\tabcolsep}{4pt}
\begin{threeparttable}
\caption[Precision minus recall per configuration]{Element precision minus element recall of the first execution per configuration against each sample's first coding, with its 2.5--97.5 percentile interval over 2{,}000 article resamples: design-weighted on the evaluation sample, unweighted on the fresh census.}\label{tab:pr_gaps}
\begin{tabular}{@{}>{\raggedright\arraybackslash}p{50mm}>{\raggedright\arraybackslash}p{42mm}>{\raggedright\arraybackslash}p{42mm}@{}}
\toprule
\textbf{Configuration} & \textbf{Evaluation sample (120)} & \textbf{Fresh census (55)} \\
\midrule
Claude Sonnet 4.6 (registered) & $+0.202$ [$+0.134$, $+0.262$] & $+0.164$ [$+0.071$, $+0.245$] \\
Claude Opus 4.6 & $+0.208$ [$+0.151$, $+0.261$] & $+0.111$ [$+0.013$, $+0.187$] \\
Claude Sonnet 4.5 & $+0.297$ [$+0.233$, $+0.364$] & $+0.300$ [$+0.172$, $+0.400$] \\
Claude Haiku 4.5 & $+0.432$ [$+0.359$, $+0.499$] & $+0.416$ [$+0.293$, $+0.519$] \\
Amazon Nova Pro & $+0.223$ [$+0.161$, $+0.293$] & $+0.101$ [$-0.035$, $+0.220$] \\
NVIDIA Nemotron 3 Super & $+0.343$ [$+0.272$, $+0.415$] & $+0.261$ [$+0.135$, $+0.396$] \\
Qwen3-235B-A22B & $+0.301$ [$+0.245$, $+0.353$] & $+0.213$ [$+0.066$, $+0.365$] \\
Mercury 3, instant & $+0.396$ [$+0.294$, $+0.495$] & $+0.431$ [$+0.272$, $+0.559$] \\
Mercury 3, low & $+0.346$ [$+0.272$, $+0.413$] & $+0.329$ [$+0.219$, $+0.422$] \\
Claude Sonnet 5 & $+0.282$ [$+0.221$, $+0.333$] & $+0.223$ [$+0.126$, $+0.306$] \\
Claude Opus 5 & $+0.016$ [$-0.045$, $+0.076$] & $-0.096$ [$-0.191$, $-0.004$] \\
Claude Opus 5.5 & $+0.146$ [$+0.086$, $+0.195$] & $+0.034$ [$-0.039$, $+0.083$] \\
\addlinespace
Control: Sonnet 4.6, Claude~5 request & $+0.215$ [$+0.145$, $+0.276$] & $+0.120$ [$+0.031$, $+0.198$] \\
Control: Opus 4.6, Claude~5 request & $+0.200$ [$+0.147$, $+0.249$] & $+0.135$ [$+0.039$, $+0.209$] \\
Control: Opus 5, Opus~5.5 request & $+0.086$ [$+0.025$, $+0.143$] & $-0.094$ [$-0.204$, $+0.008$] \\
\bottomrule
\end{tabular}
\begin{tablenotes}\footnotesize
\item Point estimates are positive for 11 of the 12 configurations (request controls aside) on both samples. The interval includes zero for Opus~5 on the evaluation sample and for Nova Pro and Opus~5.5 on the census, and lies below zero for Opus~5 on the census.
\end{tablenotes}
\end{threeparttable}
\end{table}

Opus~5.5 rejects the registered request twice, so it ran with the tool chosen freely and reasoning at the lowest effort; Opus~5 under this same request is its control (Table~\ref{tab:model_config}; the results table is electronic, Appendix~\ref{app:electronic}). Against the control, Opus~5.5 shows no consistent difference in element F1, $-0.018$ [$-0.051$, $+0.019$] and $+0.074$ [$+0.006$, $+0.156$], nor against Opus~5 as run; against Sonnet~4.6 and Opus~4.6 it is better on both samples. Its three executions keep precision above recall on the evaluation sample ($+0.146$ to $+0.186$). Opus~5's balance also recurs in three executions a week later, so neither gap is a draw of one execution. At high instead of low effort, Opus~5.5 shows no consistent difference in element F1, lower recall on the evaluation sample ($-0.036$ [$-0.060$, $-0.012$]) and higher precision on the census ($+0.094$ [$+0.033$, $+0.150$]).

\subsection{Translated Input}\label{app:language}\label{subsec:language_results}

The comparison was declared on 19 October 2026 and amended on 5 November 2026 before any paragraph was translated. It sets the original German input against its machine translation into English, both scored against the unchanged first coding. It isolates no language effect, because translation also changes wording, explicitness and length. Translation cost USD 0.02; the extraction run \modelid{a2de_translated_sonnet_20261105} made 188 calls at USD 1.04. The detailed table is electronic (Appendix~\ref{app:electronic}).

Translation leaves agreement inconclusive under the rule declared before the run, showing neither a difference nor equivalence. The 60 German evaluation paragraphs were translated into English once with gpt-oss-120b. The coder checked each and flagged none. Against the first coding, translated minus original is $-0.022$ [$-0.084$, $+0.032$] for first-execution element micro-F1. The paired differences in technology micro-F1 ($+0.017$ [$-0.027$, $+0.069$]), element precision ($-0.034$ [$-0.120$, $+0.043$]) and element recall ($-0.014$ [$-0.068$, $+0.040$]) are descriptive, and their intervals also include zero. On the 55 paragraphs where both first executions succeeded, the element set changed in 30 of 70 technology rows predicted from both inputs. This rate of 43\% compares with 6\% and 3\% between the original's first and repeated executions (all unweighted). The translated input, however, ran 23 days after the original.

Future studies should report precision and recall per label dimension and language, on source text and against a reference whose stability they report.

\appfloatbarrier

\subsection{Instrument Checks: Perturbation Battery and Mandatory Form}\label{app:battery}

Table~\ref{tab:instrument_checks} summarises four descriptive instrument checks and points to their detailed tables across the appendix. The battery rewords single codebook definitions on the 443 frame-completion paragraphs and asks where the wording rather than the content drives the labels (Table~\ref{tab:battery_primary}; its design and conduct record is electronic, Appendix~\ref{app:electronic}). Table~\ref{tab:battery_sharpen} lists the codebook rules it recommends settling. Section~\ref{app:form_variant} reports the mandatory form, which asks for a yes or no on every element.

% 2026-11-24: four rows moved unchanged from tab_added_analyses; hand-maintained; detailed sources remain in the cited appendix tables.
% OUTLINE: Four descriptive instrument checks summarise retry handling, wording sensitivity, mandatory decisions and development iterations; none replaced the registered instrument.
\begin{table}[htbp]
\centering
\scriptsize
\setlength{\tabcolsep}{3pt}
\begin{threeparttable}
\caption[Instrument checks added after the registered result]{Instrument checks added after the registered result, supplementing Table~\ref{tab:added_analyses}. Descriptive; no variant replaced the registered instrument.}
\label{tab:instrument_checks}
\begin{tabular}{@{}>{\raggedright\arraybackslash}p{36mm}>{\raggedright\arraybackslash}p{50mm}>{\raggedright\arraybackslash}p{65mm}@{}}
\toprule
\textbf{Analysis (appendix tables)} & \textbf{Question; material} & \textbf{Result} \\
\midrule
Retry candidate (Tables~\ref{tab:a3_candidate} and~\ref{tab:form_variant}) & Does validator feedback recover failed output? Both samples, later runs & Census: 0 against 23 terminal failures of 165 executions; F1 0.662 against 0.563. Evaluation sample: F1 0.646 against 0.643 \\
Perturbation battery (Tables~\ref{tab:battery_primary} and~\ref{tab:battery_sharpen}) & Do labels follow wording rather than content? 443 paragraphs outside both samples & Labels move with definition wording, frame order and element labels; sensitivity at the documented gap shown; planted gaps not shown, as two computations disagree \\
Mandatory form (Table~\ref{tab:form_variant}) & Does a yes or no per element reduce the deficit? Both samples, against the retry candidate & Recall $+0.019$ [$-0.046$, $+0.078$] (evaluation sample) and $+0.079$ [$+0.031$, $+0.116$] (census); output cap doubled \\
Development iterations (Table~\ref{tab:development_review}) & Do instrument changes help? 32 development paragraphs & None adopted: mean element F1 0.591 to 0.656, registered instrument 0.663 \\
\bottomrule
\end{tabular}
\begin{tablenotes}\footnotesize
\item Evaluation sample design-weighted, census unweighted. F1: element micro-F1. Terminal failure: no valid record after retry. Brackets: 95\% intervals over paired article resamples. Battery analysis code was corrected after its first output.
\end{tablenotes}
\end{threeparttable}
\end{table}


\begingroup\appcaption{Battery: planted and documented gaps}{Primary readings of the battery, route-stratified and as declared; the battery finds gaps of known position only if the intervals of the PR1 mean and the PR2 contrast both lie above zero. The rules were fixed before any call, the route stratification before any variant result was read but after a smoke test had printed interim contrasts.}% Generated by tools/pipeline/battery_analyse.py from tools/runs/battery/analysis.json (specification §4i). Do not edit by hand.
\begin{table}[htbp]
\centering
\scriptsize
\setlength{\tabcolsep}{4pt}
\begin{threeparttable}
\caption[Battery: planted and documented gaps]{Primary readings of the battery: the rules were fixed before any call (specification §4i), the route stratification on 30 October 2026 after the switch of inference route and before any variant result was read, though after a smoke test of the analysis code on partial data had printed interim contrasts. Shifts on the declared instability (1 minus Jaccard). The statement ``the battery finds gaps whose position is known'' requires both intervals above zero: route-stratified not shown at this $n$, as declared both intervals lie above zero. Values without an interval are the components of the contrast; \dag\ marks an element with a union below ten paragraphs in some pair against baseline executions 1--3. Article resampling, $B = 2{,}000$, seed 20260905.}
\label{tab:battery_primary}
\begin{tabular}{@{}>{\raggedright\arraybackslash}p{62mm}>{\raggedright\arraybackslash}p{44mm}>{\raggedright\arraybackslash}p{44mm}@{}}
\toprule
\textbf{Quantity} & \textbf{Route-stratified [interval]} & \textbf{As declared [interval]} \\
\midrule
\multicolumn{3}{@{}l}{\emph{PR1, planted gaps: shift with the element's own rules removed minus with the other half's removed}} \\
\texttt{F2c} & $-0.063$ [$-0.164$, $+0.020$] & $-0.076$ [$-0.136$, $-0.019$] \\
\texttt{F4b} & $+0.061$ [$-0.066$, $+0.195$] & $+0.089$ [$+0.000$, $+0.211$] \\
\texttt{F5a} & $+0.129$ [$-0.032$, $+0.348$] & $+0.153$ [$-0.000$, $+0.312$] \\
\texttt{F8a} & $+0.007$ [$-0.119$, $+0.119$] & $+0.008$ [$-0.105$, $+0.102$] \\
\texttt{F1a} & $-0.003$ [$-0.083$, $+0.091$] & $-0.044$ [$-0.105$, $+0.007$] \\
\texttt{F2b} & $+0.044$ [$-0.065$, $+0.157$] & $+0.040$ [$-0.057$, $+0.141$] \\
\texttt{F6a} & $+0.076$ [$+0.007$, $+0.170$] & $+0.093$ [$+0.016$, $+0.186$] \\
\texttt{F6b} & $+0.096$ [$-0.165$, $+0.314$] & $+0.093$ [$-0.157$, $+0.317$] \\
Mean over the eight & $+0.043$ [$-0.002$, $+0.090$] & $+0.045$ [$+0.002$, $+0.090$] \\
\addlinespace
\multicolumn{3}{@{}l}{\emph{PR2, documented gap: \texttt{general} under K1 minus the mean of hydrogen and nuclear under K1c}} \\
\texttt{general}, K1 & $+0.103$ & $+0.103$ \\
\texttt{hydrogen}, K1c & $-0.010$ & $-0.004$ \\
\texttt{nuclear}, K1c & $+0.001$ & $+0.000$ \\
Contrast & $+0.108$ [$+0.056$, $+0.167$] & $+0.105$ [$+0.059$, $+0.153$] \\
\bottomrule
\end{tabular}
\end{threeparttable}
\end{table}
\endgroup

The battery's one validation on labelled units is weak. On the evaluation sample, the number of variants whose label for a paragraph and element differs from the baseline separates units where model and first coding disagree from the others with an area under the curve (AUC) of 0.62 [0.59, 0.66]. For units where the first coding and the registered re-code disagree, it is 0.56 [0.52, 0.61]. The rules of Table~\ref{tab:battery_sharpen} are therefore candidates for review by coders, not validated diagnoses.

Route stratification was fixed on 30 October 2026, after a smoke test had printed interim contrasts but before any variant result was read. Under it, the mean planted-gap contrast PR1 is $+0.043$ [$-0.002$, $+0.090$] and includes zero. As first declared, it is $+0.045$ [$+0.002$, $+0.090$], just above zero (Table~\ref{tab:battery_primary}). As the two computations disagree, the battery is not shown to find gaps of known position, and the specification makes no gap-detector claim. The documented gap of the general category shows sensitivity under both. This sensitivity is not specific to definitions. Under route stratification, reversing the frame order moves 14 of 24 reported classes beyond repetition noise. Neutral element labels move 13 of 23, and even the negative-control sentence on hydrogen and nuclear moves 6 of 23 (as first declared: 15, 16 and 9).

\begingroup\scriptsize
\setlength{\tabcolsep}{4pt}
\begin{longtable}{@{}>{\raggedright\arraybackslash}p{14mm}>{\raggedright\arraybackslash}p{84mm}>{\raggedright\arraybackslash}p{50mm}@{}}
\caption[Battery: rules to settle before the next freeze]{Codebook rules to settle before the next freeze, from the flips of the perturbation battery (specification §4i), in the order of what they do for comparable frame shares; condensed from the full texts in \texttt{sharpening\_recommendations.json}. Cases decided: of the flipping cases in the unit's packet, those the rule decides, as counted by the verifying agent.}\label{tab:battery_sharpen} \\
\toprule
\textbf{Units} & \textbf{Rule (condensed)} & \textbf{Evidence} \\
\midrule
\endfirsthead
\toprule
\textbf{Units} & \textbf{Rule (condensed)} & \textbf{Evidence} \\
\midrule
\endhead
\bottomrule
\endfoot
\texttt{general} & Tick \texttt{general} for every bundle term, also beside named technologies. Bundle terms include renewables, green or clean power or energy, the energy transition, and collective clean-technology terms. Assign elements from claims whose subject is the bundle. Without a bundle term in the paragraph, create no \texttt{general} row; it is never a residual for fossil, grid or energy-system statements. & PR2 $+0.108$; the residual sentence removed 2 of 58 rows; 7 of 8 rewordings resolved the reading; 37 of 46 cases decided. \\
F3b & F3b requires an anchor in the statement: a named shift of the energy system, a coal, nuclear or fossil phase-out named as such (or its reversal), or a stated substitution between the technology and coal, oil or gas. Subsidies, laws, investments, policy turns, growth figures and ``green'' in a technology's name are not anchors and go to their own elements. & Swap occurrences into F3b 54, 42, 34, 31 (8, 6, 8, 3 relations); the only element displaced upward on both samples ($+12.7$ points on the census, $+3.4$ on the evaluation sample); 116 of 126 cases. \\
F6a, F2b & F6a only for technical possibility or maturity; a decision or commitment about an identified project (announced, planned, financed, construction started, postponed, shelved, cancelled) is F2b, the stated reason fires its own element, and a milestone without decision content (completion, commissioning, grid connection) fires neither. & F6a planted contrast $+0.076$ [$+0.007$, $+0.170$]; swap occurrences F6a--F2b 38 and 33 (8 and 6 relations); 70 of 75 cases. \\
F1a, F2b & What buyers and users pay is F1a; producers' revenue, profit, financing and capital are F2b; F1a and F1b never displace F2b, each firing where its own definition is met. & F2b underdetermined (both admitted one-unit rewordings raise its rate, $+2.4$ and $+2.6$ points) and displaced most on the evaluation sample ($-18.3$); swap occurrences F2b--F1a 31 (8 relations); 22 of 28 cases. \\
\texttt{general} & A bundle term with no claim of its own (a reference quantity, a label for a firm, the power source of something else such as renewable hydrogen) is \texttt{general} without a frame and inherits no element. & 10 flips on two paragraphs under five families; 14 cases. \\
F2c & F2c is a change in the amount deployed, produced or sold, with or without a number; targets, plans, stated needs and static shares do not fire it; aggregate cancellations are a decline. & Coder $\alpha$ 0.37; three of four rewordings move it; 31 of 34 and 14 of 14 cases. \\
F2c, F6b & F2c is the technology's own deployment, production, sales or project pipeline and its change, including a forecast of where deployment is heading; F6b is potential per unit, site or resource, and an aggregate quantity is never F6b. & Swap occurrences F6b--F2c 36 (5 relations); neutral labels shift F6b by $+0.341$; F6b's F1 0.43 on the evaluation sample. \\
F5a & F5a needs a listed trigger or a stated time dimension (power missing or in surplus at hours, days or seasons); overall quantities, capacity factors and supply lost to shutdowns do not fire it unless the text names variability. & Three of four rewordings move it; swap occurrences F5a--F6b 18, F5a--F5b 23 (3 and 6 relations); 63 of 84 cases. \\
\end{longtable}
\endgroup

\subsubsection{Mandatory Form Variant}\label{app:form_variant}

The variant v1.0F mechanically transforms v1.0r3, the registered request under the feedback retry. The list of fired elements becomes seventeen required yes-or-no fields per technology, and four prompt sentences are replaced by their counterparts. A normaliser maps the answers back before the registered validator and evaluator apply. Its output cap of 2,048 tokens instead of 1,024 was set from the longest journalled answers re-encoded in the form, about 1,000 tokens. The comparator is v1.0r3, run beside it on the evaluation sample and, on the census, its run of 26 October.

Table~\ref{tab:form_variant} compares both configurations on the two samples.

% Generated by tools/pipeline/compare_form.py from tools/runs/tables/form_variant/form_comparison.json (specification §4j). Do not edit by hand.
\begin{table}[htbp]
\centering
\footnotesize
\setlength{\tabcolsep}{4pt}
\begin{threeparttable}
\caption[Mandatory form against a list]{Mandatory form against a list (specification §4j): v1.0F answers each of the 17 elements yes or no for every technology, v1.0r3 lists the elements that fire; same model, decoding and retry policy, output cap 2,048 against 1,024. Differences are form minus list under common article resamples ($B = 2{,}000$, seed 20260905). Both samples were inspected before; descriptive, nothing is selected.}
\label{tab:form_variant}
\begin{tabular}{@{}>{\raggedright\arraybackslash}p{40mm}>{\raggedright\arraybackslash}p{28mm}>{\raggedright\arraybackslash}p{28mm}>{\raggedright\arraybackslash}p{28mm}>{\raggedright\arraybackslash}p{22mm}@{}}
\toprule
\textbf{Configuration} & \textbf{Element recall} & \textbf{Element precision} & \textbf{Element micro-F1} & \textbf{Invalid first / failed} \\
\midrule
\multicolumn{5}{@{}l}{\emph{Registered evaluation sample A2 (120 paragraphs, design-weighted; both $k = 1$)}} \\
v1.0F (mandatory form) & 0.582 & 0.744 & 0.653 & 10 / 0 \\
v1.0r3 (list, comparator) & 0.562 & 0.758 & 0.646 & 6 / 0 \\
Difference [interval] & $+0.019$ [$-0.046$, $+0.078$] & $-0.014$ [$-0.088$, $+0.055$] & $+0.007$ [$-0.053$, $+0.065$] & -- \\
\addlinespace
\multicolumn{5}{@{}l}{\emph{Fresh-sample census A3 (55 paragraphs; execution 1 of each)}} \\
v1.0F (mandatory form) & 0.700 & 0.726 & 0.713 & 12 / 0 \\
v1.0r3 (list, comparator) & 0.621 & 0.707 & 0.662 & 6 / 0 \\
Difference [interval] & $+0.079$ [$+0.031$, $+0.116$] & $+0.019$ [$-0.025$, $+0.065$] & $+0.051$ [$+0.030$, $+0.071$] & -- \\
\bottomrule
\end{tabular}
\end{threeparttable}
\end{table}


A unary breakdown into one question per frame or element \parencite[pp.~2479--2480]{ma2025multilabel} would test a remedy, not the deficit's mechanism: prompts and inference budget also change.

\subsection{Checks of the Reference, the Misses and the Context}\label{app:later_checks}

Five computed checks, one of them a comparison-target baseline, and a later-coded re-reading of sampled misses examine the reference, the missed elements and the input context of the registered run. They were declared in the pre-analysis specification (\S4l) on 6 November 2026, before any of them was computed or, for the re-reading of misses, coded. Intervals for the coder's $\alpha$ and a second re-code wave followed on 10 November (\S4n), seven further checks the same day (\S4o) and a last set on 11 November (\S4p), likewise declared before computation and coding. They are descriptive and reuse samples whose results had already been inspected; nothing is selected.

Table~\ref{tab:recode_comparison} scores the model and the first coding against the coder's blind re-code of 40 paragraphs from 28 articles. With the second coding as reference, the model reaches precision 0.662, recall 0.584 and F1 0.620, and the first coding 0.602, 0.720 and 0.656, so the first coding marks more elements than the second. On the same paragraphs, the model's recall against the second coding exceeds that against the first (0.496) by $+0.088$ [$+0.002$, $+0.178$]. The comparison measures one coder's temporal stability after a minimum five-day interval instead of the registered 14 (Table~\ref{tab:deviations}), not an intercoder ceiling.

% Generated by tools/pipeline/review_analyses.py from tools/runs/tables/review_analyses/results.json (specification §4l). Do not edit by hand.
% OUTLINE: The registered re-code supports a descriptive comparison against the second reading after a minimum five-day interval.
\begin{table}[htbp]
\centering
\footnotesize
\setlength{\tabcolsep}{4pt}
\begin{threeparttable}
\caption[Agreement against the blind re-code]{Later element-agreement comparisons with the registered blind re-code of 40 paragraphs (specification §4l (2)). The second coding is the reference in the first two rows of each block. Design-weighted; 95\% intervals over article resamples ($B = 2{,}000$, seed 20260905). One coder's consistency after a minimum five-day interval, not an intercoder ceiling.}
\label{tab:recode_comparison}
\begin{tabular}{@{}>{\raggedright\arraybackslash}p{52mm}>{\raggedright\arraybackslash}p{30mm}>{\raggedright\arraybackslash}p{30mm}>{\raggedright\arraybackslash}p{30mm}@{}}
\toprule
\textbf{Comparison} & \textbf{Precision} & \textbf{Recall} & \textbf{Micro-F1} \\
\midrule
\multicolumn{4}{@{}l}{\emph{Both languages (40 paragraphs, 28 articles)}} \\
Model against the second coding & 0.662 [0.594, 0.732] & 0.584 [0.492, 0.677] & 0.620 [0.548, 0.683] \\
First coding against the second coding & 0.602 [0.504, 0.695] & 0.720 [0.628, 0.810] & 0.656 [0.570, 0.732] \\
Model against the first coding (registered reference) & 0.673 [0.555, 0.785] & 0.496 [0.386, 0.601] & 0.571 [0.460, 0.669] \\
Model recall, second minus first coding & -- & $+0.088$ [0.002, 0.178] & -- \\
\addlinespace
\multicolumn{4}{@{}l}{\emph{German (20 paragraphs, 15 articles)}} \\
Model against the second coding & 0.687 [0.580, 0.841] & 0.630 [0.462, 0.790] & 0.657 [0.533, 0.758] \\
First coding against the second coding & 0.584 [0.453, 0.738] & 0.682 [0.479, 0.849] & 0.630 [0.482, 0.746] \\
Model against the first coding (registered reference) & 0.627 [0.458, 0.790] & 0.492 [0.336, 0.649] & 0.552 [0.395, 0.693] \\
Model recall, second minus first coding & -- & $+0.137$ [-0.034, 0.296] & -- \\
\addlinespace
\multicolumn{4}{@{}l}{\emph{English (20 paragraphs, 13 articles)}} \\
Model against the second coding & 0.649 [0.555, 0.725] & 0.562 [0.446, 0.671] & 0.602 [0.506, 0.678] \\
First coding against the second coding & 0.610 [0.493, 0.735] & 0.739 [0.644, 0.837] & 0.668 [0.565, 0.769] \\
Model against the first coding (registered reference) & 0.697 [0.540, 0.830] & 0.497 [0.347, 0.642] & 0.580 [0.427, 0.713] \\
Model recall, second minus first coding & -- & $+0.064$ [-0.038, 0.170] & -- \\
\bottomrule
\end{tabular}
\end{threeparttable}
\end{table}


Table~\ref{tab:alpha_full_grid} recomputes Krippendorff's $\alpha$ of the re-code on the full grid of 55 rows that either coding created, a missing row counting as unmarked. Pooled over element cells, $\alpha$ is 0.619 on this grid against 0.641 on the 50 rows both codings created, and per element the two grids differ little.

% Generated by tools/pipeline/review_analyses.py from tools/runs/tables/review_analyses/results.json (specification §4l (4), §4n). Do not edit by hand.
% OUTLINE: The coder's alpha with intervals, per population: supports, raw agreement, alpha with article-resampling intervals, q and undefined draws.
\begingroup\scriptsize
\setlength{\tabcolsep}{3pt}
\setlength{\LTpost}{3pt}
\begin{longtable}{@{}r>{\raggedright\arraybackslash}p{22mm}rrrr>{\raggedright\arraybackslash}p{31mm}rrr@{}}
\caption[The coder's alpha with intervals]{The coder's $\alpha$ in the registered blind re-code of 40 paragraphs from 28 articles. Article-resampling intervals were added later (specification §4n (2)). Both codings created 50 shared rows; their union contains 55 rows, 1 exclusive to the first coding and 4 to the second. Shared-row $\alpha$ values reproduce Table~\ref{tab:reliability}.}\label{tab:alpha_full_grid} \\
\toprule
\textbf{Set} & \textbf{Item} & \textbf{Units} & \textbf{1st/2nd} & \textbf{Art.} & \textbf{Raw} & \textbf{$\alpha$ [95\,\% interval]} & \textbf{$q_{.667}$} & \textbf{$q_{.800}$} & \textbf{Undef.} \\
\midrule
\endfirsthead
\toprule
\textbf{Set} & \textbf{Item} & \textbf{Units} & \textbf{1st/2nd} & \textbf{Art.} & \textbf{Raw} & \textbf{$\alpha$ [95\,\% interval]} & \textbf{$q_{.667}$} & \textbf{$q_{.800}$} & \textbf{Undef.} \\
\midrule
\endhead
\bottomrule
\endfoot
\bottomrule
\endlastfoot
\multicolumn{10}{@{}l}{\textit{Technology ticks (paragraphs)}} \\
40 & general & 40 & 8/9 & 7 & 0.925 & 0.779 [0.430, 1.000] & 0.2140 & 0.5715 & 0 \\
40 & solar & 40 & 11/11 & 9 & 1.000 & 1.000 [1.000, 1.000] & 0.0000 & 0.0000 & 0 \\
40 & wind & 40 & 9/9 & 7 & 1.000 & 1.000 [1.000, 1.000] & 0.0000 & 0.0000 & 0 \\
40 & hydrogen & 40 & 10/10 & 5 & 1.000 & 1.000 [1.000, 1.000] & 0.0000 & 0.0000 & 10 \\
40 & nuclear & 40 & 9/9 & 6 & 1.000 & 1.000 [1.000, 1.000] & 0.0000 & 0.0000 & 1 \\
40 & other\_renewable & 40 & 1/1 & 1 & 1.000 & $n<5$ & -- & -- & 734 \\
40 & storage & 40 & 3/5 & 3 & 0.950 & $n<5$ & -- & -- & 11 \\
\multicolumn{10}{@{}l}{\textit{Elements, rows both codings created}} \\
40 & F1a & 50 & 10/9 & 7 & 0.900 & 0.678 [0.242, 1.000] & 0.4825 & 0.7470 & 0 \\
40 & F1b & 50 & 6/5 & 5 & 0.980 & 0.899 [0.543, 1.000] & 0.0852 & 0.2130 & 5 \\
40 & F2a & 50 & 0/1 & 0 & 0.980 & $n<5$ & -- & -- & 741 \\
40 & F2b & 50 & 20/17 & 14 & 0.900 & 0.788 [0.548, 0.955] & 0.1415 & 0.5305 & 0 \\
40 & F2c & 50 & 8/9 & 7 & 0.820 & 0.369 [-0.104, 0.703] & 0.9585 & 0.9940 & 0 \\
40 & F3a & 50 & 4/4 & 4 & 0.960 & $n<5$ & -- & -- & 32 \\
40 & F3b & 50 & 8/5 & 7 & 0.900 & 0.562 [0.175, 0.848] & 0.7680 & 0.9430 & 0 \\
40 & F4a & 50 & 3/2 & 3 & 0.940 & $n<5$ & -- & -- & 28 \\
40 & F4b & 50 & 2/3 & 1 & 0.940 & $n<5$ & -- & -- & 74 \\
40 & F5a & 50 & 4/6 & 3 & 0.920 & $n<5$ & -- & -- & 23 \\
40 & F5b & 50 & 3/2 & 3 & 0.940 & $n<5$ & -- & -- & 34 \\
40 & F6a & 50 & 8/4 & 6 & 0.920 & 0.625 [-0.051, 0.919] & 0.6221 & 0.8378 & 2 \\
40 & F6b & 50 & 5/3 & 2 & 0.920 & 0.462 [-0.043, 0.650] & 1.0000 & 1.0000 & 267 \\
40 & F7a & 50 & 2/3 & 1 & 0.980 & $n<5$ & -- & -- & 709 \\
40 & F7b & 50 & 2/1 & 2 & 0.980 & $n<5$ & -- & -- & 273 \\
40 & F8a & 50 & 12/6 & 9 & 0.880 & 0.598 [0.174, 0.922] & 0.6465 & 0.8760 & 0 \\
40 & F9 & 50 & 0/0 & 0 & 1.000 & $n<5$ & -- & -- & 2{,}000 \\
40 & pooled cells & 850 & 97/80 & 26 & 0.933 & 0.641 [0.555, 0.723] & 0.7545 & 1.0000 & 0 \\
\multicolumn{10}{@{}l}{\textit{Elements, full grid of rows either coding created}} \\
40 & F1a & 55 & 10/9 & 7 & 0.909 & 0.685 [0.251, 1.000] & 0.4670 & 0.7410 & 0 \\
40 & F1b & 55 & 7/5 & 5 & 0.964 & 0.815 [0.484, 1.000] & 0.1764 & 0.4702 & 5 \\
40 & F2a & 55 & 0/2 & 0 & 0.964 & $n<5$ & -- & -- & 268 \\
40 & F2b & 55 & 21/18 & 14 & 0.873 & 0.724 [0.470, 0.909] & 0.3230 & 0.7695 & 0 \\
40 & F2c & 55 & 8/9 & 7 & 0.836 & 0.380 [-0.090, 0.709] & 0.9535 & 0.9940 & 0 \\
40 & F3a & 55 & 4/4 & 4 & 0.964 & $n<5$ & -- & -- & 32 \\
40 & F3b & 55 & 8/6 & 7 & 0.891 & 0.513 [0.158, 0.779] & 0.8835 & 0.9865 & 0 \\
40 & F4a & 55 & 3/2 & 3 & 0.945 & $n<5$ & -- & -- & 28 \\
40 & F4b & 55 & 2/3 & 1 & 0.945 & $n<5$ & -- & -- & 74 \\
40 & F5a & 55 & 4/7 & 3 & 0.909 & $n<5$ & -- & -- & 23 \\
40 & F5b & 55 & 3/2 & 3 & 0.945 & $n<5$ & -- & -- & 34 \\
40 & F6a & 55 & 8/4 & 6 & 0.927 & 0.629 [-0.045, 0.921] & 0.6186 & 0.8333 & 2 \\
40 & F6b & 55 & 5/3 & 2 & 0.927 & 0.466 [-0.037, 0.651] & 1.0000 & 1.0000 & 267 \\
40 & F7a & 55 & 2/3 & 1 & 0.982 & $n<5$ & -- & -- & 709 \\
40 & F7b & 55 & 2/1 & 2 & 0.982 & $n<5$ & -- & -- & 273 \\
40 & F8a & 55 & 12/6 & 9 & 0.891 & 0.605 [0.187, 0.923] & 0.6295 & 0.8745 & 0 \\
40 & F9 & 55 & 0/0 & 0 & 1.000 & $n<5$ & -- & -- & 2{,}000 \\
40 & pooled cells & 935 & 99/84 & 26 & 0.933 & 0.619 [0.530, 0.697] & 0.8910 & 1.0000 & 0 \\
\multicolumn{10}{@{}l}{\textit{Evaluations, rows on which both codings fired the frame}} \\
40 & costs & 12 & 12 & 10 & 1.000 & 1.000 [1.000, 1.000] & 0.0000 & 0.0000 & 35 \\
40 & jobs & 19 & 19 & 13 & 0.895 & 0.787 [0.508, 1.000] & 0.1656 & 0.5238 & 1 \\
40 & climate & 8 & 8 & 8 & 0.750 & 0.492 [-0.118, 1.000] & 0.8064 & 0.9171 & 94 \\
40 & security & 4 & 4 & 4 & 1.000 & $n<5$ & -- & -- & 2{,}000 \\
40 & reliability & 4 & 4 & 3 & 1.000 & $n<5$ & -- & -- & 938 \\
40 & feasibility & 7 & 7 & 4 & 0.857 & 0.755 [-0.167, 1.000] & 0.5138 & 0.6747 & 334 \\
40 & siting & 3 & 3 & 2 & 0.667 & $n<5$ & -- & -- & 709 \\
40 & regulation & 6 & 6 & 5 & 0.833 & 0.756 [-0.222, 1.000] & 0.4531 & 0.6198 & 135 \\
40 & other & 0 & 0 & 0 & -- & $n<5$ & -- & -- & 2{,}000 \\
\multicolumn{10}{@{}l}{\textit{Comparison target, rows both codings created}} \\
40 & all rows & 50 & 50 & 27 & 0.880 & 0.624 [0.260, 0.875] & 0.6250 & 0.9165 & 0 \\
\end{longtable}
% OUTLINE: The note defines units, supports, weighting, the resampling and the suppression rule.
\par\noindent{\footnotesize Set: the population, by its number of paragraphs. Units are paragraphs for technology ticks, technology rows for individual elements and comparison targets, and shared frame-bearing rows for evaluations. Pooled units are technology-row $\times$ element cells. 1st/2nd gives positive counts for binary items and applicable-unit counts for categorical items. Art. counts articles with a first-coding positive for binary items, and articles with an applicable unit otherwise. Raw is the share of units on which both codings agree; it and $\alpha$ are unweighted. Intervals use defined draws from 2{,}000 whole-article resamples (seed 20260905). $q$ is the fraction of defined draws below the stated benchmark. Undef. counts draws with no applicable units or no variation across ratings. An item below five first-coding support units has its $\alpha$ suppressed. Pooled cells summarise heterogeneous elements and certify none of them.\par}
\par\bigskip
\endgroup


The table adds 95\,\% intervals over 2{,}000 resamples of the 28 articles, together with $q$, the fraction of defined resamples below 0.667 or 0.800 \parencite[pp.~237--238]{krippendorff2004contentanalysis}. For a smallest category share of 10\,\%, Krippendorff's planning table lists 86 units, compared with 41 at 25\,\% ($\alpha_{\min}=0.667$, 5\,\% level) \parencite[p.~240]{krippendorff2004contentanalysis}; these planning examples do not account for article clustering. Undefined resamples are excluded from the intervals and $q$, and their counts are reported. Resampling whole articles preserves dependence among their paragraphs and technology rows; it does not guarantee wider intervals. These analyses were declared on 10 November 2026 (\S4n), before computation. All eight reportable element intervals extend below 0.667, as do both pooled-cell intervals. Technical performance lacks variation in 267 resamples. Intervals and $q$ exclude undefined draws and therefore condition on $\alpha$ being defined, not merely on the element occurring. Perfect observed agreement can produce an interval of [1.000, 1.000]; resampling cannot introduce disagreements absent from the observed pairs.

The supplementary re-code (wave 2) covered the remaining 80 paragraphs on 10 November 2026, at least 33 days after the first coding, with the first labels hidden (\S4n). Analyses using it, alone or pooled, are exploratory. Its pooled shared-row $\alpha$ is 0.739 [0.680, 0.792] (Table~\ref{tab:alpha_waves}). Pooling both waves gives 0.711 [0.655, 0.762], mixing different intervals and exposures rather than replacing the registered estimate. Deployment growth, technical performance and regulation have point estimates below 0.667 in both waves. Before the second wave, the coder knew the results and had seen first-coding and model elements for one row in each of 25 paragraphs. Excluding those paragraphs leaves the combined shared-row point estimate at 0.711 after rounding.

\appgrab{tables/tab_alpha_waves.tex}\appkeepheads\appout

Table~\ref{tab:recode_waves} scores the model and the first coding against the second coding of all 120 paragraphs. The model reaches precision 0.695 and recall 0.666, against 0.759 and 0.557 with the first coding as reference; its recall rises by $+0.109$ [$+0.056$, $+0.161$]. The rise appears in German ($+0.135$) and English ($+0.096$) and without the 25 re-read paragraphs ($+0.062$ [$+0.012$, $+0.119$]). The first coding against the second reaches precision 0.640 and recall 0.835, so it marks more elements than the second, as it did in the registered wave.

\begingroup\let\footnotesize\scriptsize% Generated by tools/pipeline/review_analyses.py from tools/runs/tables/review_analyses/results.json (specification §4n (1), (3)). Do not edit by hand.
% OUTLINE: Mixed-wave comparisons report reference sensitivity and paired contrasts without replacing registered results.
\begingroup\footnotesize
\setlength{\tabcolsep}{4pt}
\setlength{\LTpost}{3pt}
\begin{longtable}{@{}>{\raggedright\arraybackslash}p{50mm}>{\raggedright\arraybackslash}p{29mm}>{\raggedright\arraybackslash}p{29mm}>{\raggedright\arraybackslash}p{34mm}@{}}
\caption[Agreement against both re-code waves]{Element agreement against the coder's second coding (specification §4n). Wave 1 is the registered re-code of 40 paragraphs, made before the registered results, at least five days after completion of the first coding. Wave 2 is the supplementary re-code of the remaining 80, made after the results, at least 33 days after completion of the first coding. Comparisons using supplementary labels, alone or pooled, are exploratory. Combined results mix these intervals and exposures. The second coding is the reference in the first two rows of each block. One coder's temporal stability, not an intercoder ceiling.}\label{tab:recode_waves} \\
\toprule
\textbf{Comparison} & \textbf{Precision} & \textbf{Recall} & \textbf{Micro-F1} \\
\midrule
\endfirsthead
\toprule
\textbf{Comparison} & \textbf{Precision} & \textbf{Recall} & \textbf{Micro-F1} \\
\midrule
\endhead
\bottomrule
\endfoot
\bottomrule
\endlastfoot
\multicolumn{4}{@{}l}{\emph{Both waves, both languages (120 paragraphs, 56 articles)}} \\
Model against the second coding & 0.695 [0.639, 0.750] & 0.666 [0.606, 0.724] & 0.680 [0.636, 0.721] \\
First coding against the second coding & 0.640 [0.587, 0.695] & 0.835 [0.782, 0.880] & 0.725 [0.677, 0.767] \\
Model against the first coding & 0.759 [0.691, 0.824] & 0.557 [0.491, 0.631] & 0.643 [0.581, 0.705]$^{\ast}$ \\
Model precision minus recall, against the second coding & \multicolumn{3}{@{}l}{$+0.029$ [-0.048, 0.106]} \\
Model precision minus recall, against the first coding & \multicolumn{3}{@{}l}{$+0.202$ [0.134, 0.262]} \\
Change of that gap, second minus first & \multicolumn{3}{@{}l}{$-0.173$ [-0.228, -0.117]} \\
F1, model minus first coding, against the second & \multicolumn{3}{@{}l}{$-0.044$ [-0.092, 0.002]} \\
\addlinespace
\multicolumn{4}{@{}l}{\emph{Both waves, German (60 paragraphs, 28 articles)}} \\
Model against the second coding & 0.712 [0.620, 0.795] & 0.687 [0.597, 0.756] & 0.700 [0.618, 0.763] \\
First coding against the second coding & 0.607 [0.514, 0.697] & 0.753 [0.668, 0.822] & 0.672 [0.593, 0.736] \\
Model against the first coding & 0.710 [0.608, 0.817] & 0.552 [0.466, 0.647] & 0.621 [0.532, 0.710]$^{\ast}$ \\
Model precision minus recall, against the second coding & \multicolumn{3}{@{}l}{$+0.025$ [-0.060, 0.112]} \\
Model precision minus recall, against the first coding & \multicolumn{3}{@{}l}{$+0.157$ [0.070, 0.240]} \\
Change of that gap, second minus first & \multicolumn{3}{@{}l}{$-0.132$ [-0.243, -0.033]} \\
F1, model minus first coding, against the second & \multicolumn{3}{@{}l}{$+0.027$ [-0.042, 0.096]} \\
\addlinespace
\multicolumn{4}{@{}l}{\emph{Both waves, English (60 paragraphs, 28 articles)}} \\
Model against the second coding & 0.686 [0.618, 0.758] & 0.656 [0.575, 0.739] & 0.670 [0.614, 0.723] \\
First coding against the second coding & 0.655 [0.588, 0.721] & 0.876 [0.815, 0.937] & 0.749 [0.689, 0.801] \\
Model against the first coding & 0.784 [0.695, 0.862] & 0.559 [0.473, 0.651] & 0.653 [0.569, 0.731]$^{\ast}$ \\
Model precision minus recall, against the second coding & \multicolumn{3}{@{}l}{$+0.031$ [-0.076, 0.136]} \\
Model precision minus recall, against the first coding & \multicolumn{3}{@{}l}{$+0.224$ [0.135, 0.308]} \\
Change of that gap, second minus first & \multicolumn{3}{@{}l}{$-0.194$ [-0.263, -0.125]} \\
F1, model minus first coding, against the second & \multicolumn{3}{@{}l}{$-0.079$ [-0.147, -0.018]} \\
\addlinespace
\multicolumn{4}{@{}l}{\emph{Wave 1 alone, before the registered results (40 paragraphs, 28 articles)}} \\
Model against the second coding & 0.662 [0.594, 0.732] & 0.584 [0.492, 0.677] & 0.620 [0.548, 0.683] \\
First coding against the second coding & 0.602 [0.504, 0.695] & 0.720 [0.628, 0.810] & 0.656 [0.570, 0.732] \\
Model against the first coding & 0.673 [0.555, 0.785] & 0.496 [0.386, 0.601] & 0.571 [0.460, 0.669] \\
Model precision minus recall, against the second coding & \multicolumn{3}{@{}l}{$+0.078$ [-0.018, 0.174]} \\
Model precision minus recall, against the first coding & \multicolumn{3}{@{}l}{$+0.177$ [0.098, 0.261]} \\
Change of that gap, second minus first & \multicolumn{3}{@{}l}{$-0.099$ [-0.185, -0.009]} \\
F1, model minus first coding, against the second & \multicolumn{3}{@{}l}{$-0.035$ [-0.129, 0.063]} \\
\addlinespace
\multicolumn{4}{@{}l}{\emph{Wave 2 alone, after the results (80 paragraphs, 45 articles)}} \\
Model against the second coding & 0.706 [0.634, 0.772] & 0.698 [0.626, 0.773] & 0.702 [0.647, 0.749] \\
First coding against the second coding & 0.653 [0.592, 0.710] & 0.881 [0.823, 0.931] & 0.750 [0.701, 0.793] \\
Model against the first coding & 0.789 [0.704, 0.860] & 0.579 [0.491, 0.671] & 0.668 [0.588, 0.743] \\
Model precision minus recall, against the second coding & \multicolumn{3}{@{}l}{$+0.008$ [-0.098, 0.098]} \\
Model precision minus recall, against the first coding & \multicolumn{3}{@{}l}{$+0.211$ [0.118, 0.291]} \\
Change of that gap, second minus first & \multicolumn{3}{@{}l}{$-0.203$ [-0.277, -0.130]} \\
F1, model minus first coding, against the second & \multicolumn{3}{@{}l}{$-0.048$ [-0.100, 0.001]} \\
\addlinespace
\multicolumn{4}{@{}l}{\emph{Both waves without the 25 re-read paragraphs (95 paragraphs, 51 articles)}} \\
Model against the second coding & 0.681 [0.628, 0.736] & 0.686 [0.609, 0.754] & 0.684 [0.629, 0.729] \\
First coding against the second coding & 0.650 [0.586, 0.712] & 0.810 [0.754, 0.865] & 0.721 [0.667, 0.768] \\
Model against the first coding & 0.772 [0.694, 0.841] & 0.625 [0.539, 0.696] & 0.691 [0.615, 0.752] \\
Model precision minus recall, against the second coding & \multicolumn{3}{@{}l}{$-0.005$ [-0.082, 0.082]} \\
Model precision minus recall, against the first coding & \multicolumn{3}{@{}l}{$+0.148$ [0.079, 0.226]} \\
Change of that gap, second minus first & \multicolumn{3}{@{}l}{$-0.152$ [-0.222, -0.078]} \\
F1, model minus first coding, against the second & \multicolumn{3}{@{}l}{$-0.037$ [-0.090, 0.012]} \\
\addlinespace
\multicolumn{4}{@{}l}{\emph{Wave 2 without the 25 re-read paragraphs (55 paragraphs, 36 articles)}} \\
Model against the second coding & 0.692 [0.617, 0.768] & 0.756 [0.657, 0.836] & 0.723 [0.656, 0.773] \\
First coding against the second coding & 0.681 [0.608, 0.750] & 0.870 [0.791, 0.936] & 0.764 [0.703, 0.813] \\
Model against the first coding & 0.827 [0.715, 0.911] & 0.707 [0.589, 0.801] & 0.762 [0.664, 0.837] \\
Model precision minus recall, against the second coding & \multicolumn{3}{@{}l}{$-0.064$ [-0.181, 0.066]} \\
Model precision minus recall, against the first coding & \multicolumn{3}{@{}l}{$+0.120$ [0.013, 0.242]} \\
Change of that gap, second minus first & \multicolumn{3}{@{}l}{$-0.184$ [-0.279, -0.081]} \\
F1, model minus first coding, against the second & \multicolumn{3}{@{}l}{$-0.042$ [-0.102, 0.025]} \\
\end{longtable}
% OUTLINE: The note distinguishes registered intervals from later paired calculations and states the exclusion and undefined-draw conventions.
\par\noindent{\footnotesize Design-weighted; 95\,\% percentile intervals over 2{,}000 resamples of whole articles (seed 20260905), common to the rows of a block, so differences are paired. $^{\ast}$: the registered interval of Table~\ref{tab:primary}; the other intervals of that row are later calculations. Paired rows were added after the results of the wave were inspected (specification §4n (3)). An interval containing zero shows no equivalence. No paired resample was undefined. The blocks without the re-read paragraphs are not random subsamples.\par}
\par\bigskip
\endgroup
\endgroup

Table~\ref{tab:recode_elements} shows that the nearly balanced pooled point estimates mix two directions. Against the second coding, investment, regulation and both siting elements retain recall below precision. Recall exceeds precision for several others, including deployment growth, supply security and technical performance.

% Generated by tools/pipeline/review_analyses.py from tools/runs/tables/review_analyses/results.json (specification §4n (3)). Do not edit by hand.
% OUTLINE: Per-element weighted comparisons and unweighted wave counts show the directions concealed by pooled balance.
\begin{table}[htbp]
\centering
\footnotesize
\setlength{\tabcolsep}{4pt}
\begin{threeparttable}
\caption[Direction per element against both readings]{The model's element precision and recall against the coder's first and second coding of all 120 paragraphs, design-weighted, and the unweighted element-tuple counts (specification §4n (3), added after the results of the supplementary re-code were inspected). Analyses using the second coding are exploratory. Element codes as in Table~\ref{tab:codebook_elements}. A cell needs at least five tuples in its reading. Tuples per wave: first/second/model. Technology rows: 178 (first), 167 (second), 164 (model).}
\label{tab:recode_elements}
\begin{tabular}{@{}lrrrrrrr@{}}
\toprule
 & \multicolumn{3}{c}{\textbf{Tuples}} & \multicolumn{2}{c}{\textbf{Against the first}} & \multicolumn{2}{c}{\textbf{Against the second}} \\
\cmidrule(lr){2-4}\cmidrule(lr){5-6}\cmidrule(lr){7-8}
\textbf{Element} & \textbf{1st/2nd/model} & \textbf{Wave 1} & \textbf{Wave 2} & \textbf{Precision} & \textbf{Recall} & \textbf{Precision} & \textbf{Recall} \\
\midrule
F1a & 45/36/33 & 10/9/8 & 35/27/25 & 0.883 & 0.643 & 0.808 & 0.716 \\
F1b & 30/23/24 & 7/5/6 & 23/18/18 & 0.817 & 0.674 & 0.769 & 0.853 \\
F2a & 3/4/3 & 0/2/1 & 3/2/2 & $n<5$ & -- & $n<5$ & -- \\
F2b & 60/51/30 & 21/18/9 & 39/33/21 & 0.976 & 0.495 & 1.000 & 0.592 \\
F2c & 35/26/36 & 8/9/9 & 27/17/27 & 0.617 & 0.613 & 0.446 & 0.625 \\
F3a & 13/9/9 & 4/4/2 & 9/5/7 & 0.664 & 0.481 & 0.664 & 0.663 \\
F3b & 38/28/31 & 8/6/10 & 30/22/21 & 0.599 & 0.486 & 0.497 & 0.556 \\
F4a & 11/7/6 & 3/2/3 & 8/5/3 & 1.000 & 0.554 & 0.704 & 0.521 \\
F4b & 13/8/14 & 2/3/2 & 11/5/12 & 0.699 & 0.761 & 0.348 & 0.772 \\
F5a & 20/26/23 & 4/7/5 & 16/19/18 & 0.789 & 0.875 & 0.886 & 0.842 \\
F5b & 17/14/15 & 3/2/1 & 14/12/14 & 0.796 & 0.689 & 0.859 & 0.877 \\
F6a & 20/11/14 & 8/4/6 & 12/7/8 & 0.860 & 0.623 & 0.615 & 0.860 \\
F6b & 15/9/15 & 5/3/5 & 10/6/10 & 0.391 & 0.485 & 0.564 & 1.000 \\
F7a & 14/6/5 & 2/3/4 & 12/3/1 & 0.699 & 0.219 & 0.699 & 0.500 \\
F7b & 12/7/4 & 2/1/0 & 10/6/4 & 1.000 & 0.320 & 1.000 & 0.544 \\
F8a & 28/15/11 & 12/6/3 & 16/9/8 & 0.788 & 0.320 & 0.714 & 0.530 \\
F9 & 2/10/3 & 0/0/1 & 2/10/2 & $n<5$ & -- & 0.517 & 0.142 \\
\midrule
All elements & 376/290/276 & 99/84/75 & 277/206/201 & \multicolumn{4}{@{}l}{pooled: Table~\ref{tab:recode_waves}} \\
\bottomrule
\end{tabular}
\end{threeparttable}
\end{table}


Table~\ref{tab:fn_causes} splits by cause the 168 first-coding elements that execution one of the registered run missed. Of these, 140 fall within a matched technology row, 15 in technology rows the model did not create and 13 in executions without a valid record, with design-weighted shares of 86\%, 9\% and 6\%. The shares are 73\%, 9\% and 18\% in German and 91\%, 9\% and 0\% in English. Counting an element as marked when any of the three executions marks it raises recall from 0.557 to 0.596 (precision 0.749 against 0.759), so the misses are largely systematic rather than draws of one execution.

% Generated by tools/pipeline/review_analyses.py from tools/runs/tables/review_analyses/results.json (specification §4l). Do not edit by hand.
\begin{table}[htbp]
\centering
\footnotesize
\setlength{\tabcolsep}{4pt}
\begin{threeparttable}
\caption[Missed elements by cause]{Missed elements of the registered run by cause (specification §4l (3)). Counts of first-coding elements the model's execution 1 does not mark, with their design-weighted shares. Below, recall and precision when an element counts as marked if any of the three executions marks it.}
\label{tab:fn_causes}
\begin{tabular}{@{}>{\raggedright\arraybackslash}p{62mm}>{\raggedright\arraybackslash}p{27mm}>{\raggedright\arraybackslash}p{27mm}>{\raggedright\arraybackslash}p{27mm}@{}}
\toprule
\textbf{Cause} & \textbf{Both languages} & \textbf{German} & \textbf{English} \\
\midrule
Execution without a valid record & 13 (6\%) & 13 (18\%) & 0 (0\%) \\
Valid record, technology row missing & 15 (9\%) & 6 (9\%) & 9 (9\%) \\
Row present, element missing & 140 (86\%) & 54 (73\%) & 86 (91\%) \\
All missed elements & 168 (100\%) & 73 (100\%) & 95 (100\%) \\
\midrule
Element recall, execution 1 & 0.557 & 0.552 & 0.559 \\
Element recall, union of executions 1--3 & 0.596 & 0.571 & 0.608 \\
Element precision, execution 1 & 0.759 & 0.710 & 0.784 \\
Element precision, union of executions 1--3 & 0.749 & 0.701 & 0.772 \\
\bottomrule
\end{tabular}
\end{threeparttable}
\end{table}


A further check re-reads misses of the third cause. From the 140 elements missed within matched rows, 45 were drawn at random with the seed fixed in \S4l; they come from 35 paragraphs of 25 articles. For each, the coder read the paragraph, the missed element's codebook definition and the elements both sides had marked on the row, and assigned one of six categories fixed in advance (Table~\ref{tab:missed_elements}). Two categories contradicted the coder's own note; on the author's instruction, the coding agent set both to borderline, following the note and the codebook text (as returned: 17 clear misses, 8 borderline, 16 possible over-marks). No reference label was changed. The model's output was visible and the coder judged their own reference. The counts are therefore non-blind judgements, not error rates; category (c) establishes neither the prevalence of reference over-marking nor a bound on it.

\begin{table}[htbp]
\centering
\footnotesize
\setlength{\tabcolsep}{4pt}
\begin{threeparttable}
\caption[Re-reading of sampled missed elements]{The coder's re-reading of 45 elements that execution 1 of the registered run missed within matched rows (specification \S4l (6)): one category per element, fixed before coding; two categories reconciled with the coder's note. Not blind to the model's output; one coder's judgement of their own first coding; shares rounded.}
\label{tab:missed_elements}
\begin{tabular}{@{}>{\raggedright\arraybackslash}p{36mm}>{\raggedright\arraybackslash}p{56mm}>{\raggedright\arraybackslash}p{17mm}>{\raggedright\arraybackslash}p{40mm}@{}}
\toprule
\textbf{Category} & \textbf{Reading} & \textbf{Elements} & \textbf{Most frequent elements} \\
\midrule
(a) Clear miss & the paragraph states the element's cue explicitly & 16 (36\%) & transition pathway 5, investment 3 \\
(b) Borderline & the cue supports the element only under a generous reading of the codebook & 10 (22\%) & investment 5 \\
(c) Coder possibly over-marked & on re-reading, the paragraph does not support the element & 15 (33\%) & investment 5, regulation 2, technical performance 2 \\
(d) Cue outside the paragraph & the element rests on the headline, the outlet or another paragraph & 2 (4\%) & subsidies 1, landscape 1 \\
(e) Assigned elsewhere & the model expressed the content through another element or row & 1 (2\%) & emissions, climate 1 \\
(f) Other & with a note; here ecological harm that neither emissions nor landscape covers & 1 (2\%) & landscape 1 \\
\midrule
All & & 45 (100\%) & investment 13, transition pathway 6 \\
\bottomrule
\end{tabular}
\end{threeparttable}
\end{table}

Clear misses and possible over-marks occur about equally often, 16 and 15 of the 45. In every clear miss the model marked another element on the same row. Investment is the most frequent missed element, with 3 clear misses, 5 borderline cases and 5 possible over-marks. Four of these over-marks report deployment growth or decline, which codebook version 1.0 had moved out of investment, and the model marked deployment growth on each. Transition pathway accounts for 5 of the clear misses. Possible over-marks are more frequent among the German elements (8 of 18) than among the English ones (7 of 27), a split the declaration did not fix.

The topic description suggested the paragraph ``potentially combined with the headline and limited surrounding context'' as a starting point, and the context-parity run tests that suggestion for the model. Instrument v1.0C places the headline, outlet and language the coder saw before the unchanged paragraph; it ran once on both samples on 6 November 2026, compared with execution one of the runs of 13 and 19 October. Element recall is 0.539 against 0.557 on the evaluation sample ($-0.019$ [$-0.067$, $+0.026$]) and 0.500 against 0.493 on the census ($+0.007$ [$-0.065$, $+0.113$]), no consistent difference under the declared rule (Table~\ref{tab:context_variant}). Only precision on the census is higher ($+0.065$ [$+0.013$, $+0.111$]), and the precision--recall gap does not narrow (0.238 against 0.202; 0.222 against 0.164). The run tests only the model's side of the asymmetry, because the coder also met other paragraphs of an article, and it is a single execution, made weeks after its comparators.

% Generated by tools/pipeline/compare_context.py from tools/runs/tables/context_variant/context_comparison.json (specification §4l (5)). Do not edit by hand.
\begin{table}[htbp]
\centering
\footnotesize
\setlength{\tabcolsep}{4pt}
\begin{threeparttable}
\caption[Context parity]{Context parity (specification §4l (5)): v1.0C gives the model the headline, outlet and language the coder saw before the paragraph; v1.0 gives the paragraph alone. Same model, prompt, decoding and retry policy. Differences are context minus paragraph only under common article resamples ($B = 2{,}000$, seed 20260905). The comparator runs date from October 2026; descriptive, nothing is selected.}
\label{tab:context_variant}
\begin{tabular}{@{}>{\raggedright\arraybackslash}p{40mm}>{\raggedright\arraybackslash}p{28mm}>{\raggedright\arraybackslash}p{28mm}>{\raggedright\arraybackslash}p{28mm}>{\raggedright\arraybackslash}p{22mm}@{}}
\toprule
\textbf{Configuration} & \textbf{Element recall} & \textbf{Element precision} & \textbf{Element micro-F1} & \textbf{Invalid first / failed} \\
\midrule
\multicolumn{5}{@{}l}{\emph{Registered evaluation sample A2 (120 paragraphs, design-weighted; execution 1 of each)}} \\
v1.0C (headline, outlet, language) & 0.539 & 0.777 & 0.636 & 6 / 6 \\
v1.0 (paragraph only) & 0.557 & 0.759 & 0.643 & 8 / 5 \\
Difference [interval] & $-0.019$ [$-0.067$, $+0.026$] & $+0.018$ [$-0.022$, $+0.062$] & $-0.006$ [$-0.048$, $+0.034$] & -- \\
\addlinespace
\multicolumn{5}{@{}l}{\emph{Fresh-sample census A3 (55 paragraphs; execution 1 of each)}} \\
v1.0C (headline, outlet, language) & 0.500 & 0.722 & 0.591 & 4 / 3 \\
v1.0 (paragraph only) & 0.493 & 0.657 & 0.563 & 8 / 8 \\
Difference [interval] & $+0.007$ [$-0.065$, $+0.113$] & $+0.065$ [$+0.013$, $+0.111$] & $+0.027$ [$-0.034$, $+0.106$] & -- \\
\bottomrule
\end{tabular}
\end{threeparttable}
\end{table}


For the comparison target, the constant guess ``none'' reaches a matched-row accuracy of 0.674 on the 158 matched rows, against the model's 0.794 (German 0.887 against 0.881; English 0.579 against 0.755).

Seven further checks on existing data were declared on 10 November 2026 (\S4o) before computation (Table~\ref{tab:checks_4o}). (1) The German--English difference in element micro-F1 is $-0.032$ [$-0.145$, $+0.091$], too wide to show a difference or its absence; for technology detection it is $-0.077$ [$-0.189$, $-0.002$]. (2) Terminal failures are almost all violations of the relational rules after two attempts (\texttt{contract\_invalid}; 12 of 13 in the registered run), not truncations. They fall on German paragraphs in the registered run (11 of 13) but split evenly on the fresh census (12 and 11). (3) On paragraphs with a valid first execution, the registered model's precision exceeds its recall by $0.187$ [$+0.119$, $+0.252$] on the evaluation sample and by $0.057$ [$-0.006$, $+0.108$] on the census, against $0.202$ and $0.164$ end to end. Failed outputs thus explain little of the gap on the evaluation sample and most of it on the census. The point estimate stays positive in the same eleven of twelve configurations. Over the 575 eligible-frame paragraphs with a valid record, the model marks investment on 20.9\%, against 19.4\% over all 618. (4) Without the rows of the general category, the element micro-F1 is 0.658 [0.589, 0.727], against the registered 0.643. (5) Per-element intervals are wide: that of technical feasibility, [0.491, 0.878], overlaps that of technical performance, [0.087, 0.681]. (6) With the second coding of all 120 paragraphs as reference, the adjusted investment share is 33.1\% [27.1, 39.0] instead of 37.7\% [31.1, 44.3]. The adjusted shares of technical feasibility, regulation, transition pathway and land requirement fall by 5 to 9 points. The second coding mixes two waves with different intervals and exposures, so this is a sensitivity, not a better estimate. (7) The per-language element macros use 16 classes each, 15 of them shared (Table~\ref{tab:conventions}).

\appgrab{tables/tab_checks_4o.tex}\appreplace{p\cB\{50mm\cE\}}{p\cB\{46mm\cE\}}\appreplace{p\cB\{24mm\cE\}}{p\cB\{28mm\cE\}}\appkeepheads\appout

A last set of checks was declared on 11 November 2026 (\S4p), before computation, and is likewise descriptive. Table~\ref{tab:checks_4p} gathers eight separate checks in nine panels, each with its own column headings, so it is best read one panel at a time. The main text cites it mainly for panels (a1), (a2), (b), (c) and (e) in Sections~\ref{subsec:agreement} to~\ref{subsubsec:categorical}, (f) in Section~\ref{subsec:displacement_results}, (d) in Section~\ref{subsec:recurrence}, (g) in Section~\ref{subsec:mechanisms} and the speaker counts of (h) in Section~\ref{subsec:annotation}. (a1), (a2) The coder's first coding is scored against the second coding in the model's own metrics, first on the 40 paragraphs re-coded before the run, then on all 120. (b) The model's hits and misses are split by whether the second coding keeps the tuple: of 168 missed tuples, 102 are absent from the second coding, against 36 of 208 found ones. (c) The prefilter's keyword matches serve as a comparator for technology detection, and the modal value as one for evaluation. The keyword pattern for other renewables and storage is joint, so the comparison is given on five and on seven technologies; every sampled paragraph was selected by these keywords. (d) The twelve configurations are scored against the second coding. (e) Technology detection by language is repeated on paragraphs with a valid first execution, a selected subset. (f) The share differences are given against both codings and with both interval schemes; a cell below five supporting paragraphs is suppressed. The first-coding article intervals are recomputed for this table; for intermittency the upper limit prints $+5.2$ where the registered Table~\ref{tab:displacement} prints 5.3. (g) and (h) report the no-frame decision and margins of the codings. The second coding mixes a wave coded before the run with one coded after the results were known (Section~\ref{subsec:annotation}), so none of these figures replaces a registered one. Agreement between the coder's two codings is neither an intercoder ceiling nor evidence of validity.

\appgrab{tables/tab_checks_4p.tex}\appkeepheads\appout

\subsection{Development Iterations and Codebook Review}\label{app:development}\label{app:codebook_review}

Two development iterations tested changes to the instrument after the registered result, and a retrospective codebook review asked whether a model would have found the codebook's documented gaps. Table~\ref{tab:development_review} records what was tried, the outcome and whether it was adopted. Both iterations were declared on 14 October 2026, after the registered tables had been inspected. They were sealed before their runs and ran on the 32 development paragraphs only, with three executions per candidate. The review replayed each coding session before the evaluation sample: a model received the codebook version the coder had, in a second arm with that session's paragraphs, and named at most 15 issues per run. None of these modifications entered the registered instrument or result.

\appneedspace{0.15\textheight}
\begingroup\scriptsize
\setlength{\tabcolsep}{4pt}
\begin{longtable}{@{}>{\raggedright\arraybackslash}p{24mm}>{\raggedright\arraybackslash}p{62mm}>{\raggedright\arraybackslash}p{46mm}>{\raggedright\arraybackslash}p{18mm}@{}}
\caption[Development iterations and codebook review]{Development iterations and the retrospective codebook review: what was tried, the outcome, and whether it was adopted; $\bar{F}$ is the mean unweighted element micro-F1 over the executions.}\label{tab:development_review} \\
\toprule
\textbf{Candidate} & \textbf{What was tried} & \textbf{Outcome} & \textbf{Adopted} \\
\midrule
\endfirsthead
\toprule
\textbf{Candidate} & \textbf{What was tried} & \textbf{Outcome} & \textbf{Adopted} \\
\midrule
\endhead
\bottomrule
\endfoot
\multicolumn{4}{@{}l}{\emph{First development iteration (W3, specification \S4); 96 executions per candidate}} \\
v1.0 & The registered instrument as contemporaneous comparator & $\bar{F} = 0.663$; 0 invalid first responses; 0 executions lost & retained \\
v1.1S & Strict structured-output mode; one object per fired frame; one labelled demonstration; feedback retry & $\bar{F} = 0.646$; 20 invalid first responses; 4 executions lost & no (not eligible) \\
v1.1E & As v1.1S, plus a required short explanation before the labels & $\bar{F} = 0.656$; 9 invalid first responses; 1 execution lost & no (not eligible) \\
\addlinespace
\multicolumn{4}{@{}l}{\emph{Second development iteration (W4, \S4a), designed after W3}} \\
v1.2S & An element only under its own frame; a required subtype field; field order as the keys reach the client & $\bar{F} = 0.591$ against 0.663; 2 invalid first responses (recovered) against 1 & no \\
\addlinespace
\multicolumn{4}{@{}l}{\emph{Retrospective codebook review (\S4h, 30 October); 18 runs per configuration, 84 documented-gap slots}} \\
Opus 4.6 & Temperature 0, Bedrock & 38 of 84 gaps found (45\%); 16\% of issues matched a gap; USD 2.42 & no change \\
Sonnet 4.6 & Temperature 0, Bedrock & 29 (35\%); 10\%; USD 1.44 & no change \\
Mercury 3, low & Temperature 0.5, developer's API & 25 (30\%); 11\%; USD 0.08 & no change \\
Mercury 3, instant & Temperature 0.5, developer's API & 18 (21\%); 7\%; USD 0.03 & no change \\
\multicolumn{4}{@{}p{154mm}@{}}{\scriptsize A documented gap counts as found only if two blinded judges (Codex and a Claude judge) agree; they agreed on 95\% of the 336 run--gap pairs ($\kappa = 0.88$). The review is retrospective and descriptive: an unmatched issue may be a gap nobody documented, and the models may know common frame typologies.} \\
\end{longtable}
\endgroup

\appfloatbarrier
\subsection{The Rendered Instrument}\label{app:instrument}

{\emergencystretch=3em The system prompt of the registered instrument v1.0, rendered from the frozen codebook (sha256 \texttt{1edb4601}\ldots), is reproduced verbatim from the pinned file (\modelid{system-prompt_v1.0.txt}, sha256 \texttt{138cade3}\ldots), as every run of that instrument sent it. The paragraph itself is the only user message. The tool configuration that forces the single extraction tool (\modelid{tool-config_v1.0.json}, sha256 \texttt{a1e82607}\ldots) is in the electronic appendix (Table~\ref{tab:electronic}). Three sources of its sixteen quoted illustrations are established. The phrase on weak wind conditions and the figure of 132 TWh follow calibration paragraph A011\_p02; article A011 is excluded from the eligible frame for that reason. The illustration of deployment growth follows calibration paragraph A010\_p15; article A010 was not excluded and has four paragraphs in the eligible frame, two of them in the evaluation sample. Both are among the 15 paragraphs that the sensitivity check of Section~\ref{subsec:agreement} removes; the first coding marks deployment growth on one of them and the model on neither, which does not measure the illustration's influence. The one German illustration paraphrases article A002, all of whose paragraphs belong to the calibration split. The origin of the other illustrations was not traced.\par}

The prompt is not a verbatim copy of the codebook file. The renderer keeps every technology, element and evaluation definition and the frame mapping word for word, and treats the eleven exclusion rules of the codebook as its decision record of 2 September 2026 prescribes (Table~\ref{tab:rule_concordance}). The coder worked with the coding sheets generated from the codebook, the codebook file and a hand-written one-page summary of version 1.0; which of these was open in a given sitting was not recorded. Model and coder therefore share the definitions, not every rule text. One difference is substantive: only the prompt tells the model to record \emph{mixed} on a single relation when several voices evaluate in opposite directions. The first coding marks an attributed speaker on 77 of its 178 rows (Table~\ref{tab:checks_4p}); no field records opposing voices, so the effect of this difference on the evaluation figures is unknown.

\begin{table}[htbp]
\centering\scriptsize
\setlength{\tabcolsep}{4pt}
\caption[Exclusion rules in the codebook and in the prompt]{The codebook's exclusion rules and their treatment by the renderer of the prompt (decision record of 2 September 2026). The coder's summary sheet gives every rule in one line, as in the codebook.}\label{tab:rule_concordance}
\begin{tabular}{@{}>{\raggedright\arraybackslash}p{15mm}>{\raggedright\arraybackslash}p{20mm}>{\raggedright\arraybackslash}p{117mm}@{}}
\toprule
\textbf{Rule} & \textbf{Treatment} & \textbf{What differs for the model} \\
\midrule
X1, X2, X5, X6 & as written & Nothing. \\
X3 & transformed & Codebook: code an evaluation voiced by a quoted source, mark the speaker field, and do not average the source's and the outlet's evaluation into one row. Prompt: code quoted evaluations as asserted in the paragraph without inferring the outlet's endorsement; where voices oppose each other on one frame, record \emph{mixed} on the single relation. The speaker field is not asked of the model. \\
X4 & transformed & Codebook: use the paragraph only; if a referent is resolvable only outside it, set the lowest confidence and say so in the note. Prompt: use the paragraph only; the confidence field is not asked of the model. \\
X7, X11 & transformed & The operative rule is kept; a literature figure, a prevalence from calibration and an anecdote are dropped. \\
X8 & excluded & Spreadsheet hygiene; the output schema enforces the types. \\
X9 & excluded & A cost of rule X4 measured in calibration, not a rule. \\
X10 & excluded & Superseded history of the deployment-growth element, whose own definition carries the rule; the codebook's clause that a growth figure tied to the transition also fires the transition-pathway element is not in the prompt. \\
Abstention & adapted & Codebook: an empty relation list when no technology--frame relation is present. Prompt and coding sheet alike keep a row for a technology that is stated but not framed. \\
\bottomrule
\end{tabular}
\end{table}

\makeatletter
\long\def\app@instparse#1\subsubsection#2\subsubsection#3\app@stop{\subsubsection#2}
\appgrab{tables/instrument_appendix.tex}
\begingroup\let\footnotesize\scriptsize
\expandafter\app@instparse\app@content\subsubsection\app@stop
\endgroup
\makeatother

\subsection{GUIDE-LLM Reporting Checklist}\label{app:guide_llm}

The GUIDE-LLM checklist \parencite{feuerriegel2026guidellm} reports how large language models were used in this thesis, separately for the five uses U1 to U5 named at its head.

\begingroup\renewcommand{\theHtable}{guidellm.\arabic{table}}
\documentclass[12pt]{article}
\usepackage{fancyhdr}
\usepackage{graphicx}
\usepackage[
    a4paper,
    left=1.5cm,
    right=1.5cm,
    top=2cm,
    bottom=1.5cm
]{geometry}

\usepackage[table]{xcolor}
\usepackage{array}
\usepackage{pifont}
\usepackage{longtable}
\usepackage{ragged2e}

\usepackage[colorlinks=true, urlcolor=blue]{hyperref}

\providecommand{\checked}{\ding{110}}   % filled box
\providecommand{\unchecked}{\ding{111}} % empty box
\DeclareUrlCommand\modelid{\urlstyle{tt}\def\UrlBreaks{\do\.\do\-\do\:\do\/\do\_}}

\pagestyle{fancy}
\fancyhf{} % clear all fields
\fancyhead[L]{\url{http://www.llm-checklist.com}} % Left header
\fancyhead[R]{Version: 1.1} % Right header
\fancyfoot[C]{\thepage}
\renewcommand{\headrulewidth}{0pt}

\renewcommand{\maketitle}{
        \vspace{-1.5cm} \noindent
        \textbf{\Large{A checklist for reporting studies with large language models in behavioral and social science}} \\
}

\begin{document}

\providecommand{\gopen}{\textcolor{red}{\textbf{[OPEN]}}}
\providecommand{\guse}[1]{\textbf{U#1}}
\providecommand{\guideexpl}[1]{}
\newcolumntype{J}{>{\RaggedRight\arraybackslash}p{\dimexpr0.19\linewidth-2\tabcolsep-1.5\arrayrulewidth\relax}}
\newcolumntype{K}{>{\RaggedRight\arraybackslash}p{\dimexpr0.81\linewidth-2\tabcolsep-1.5\arrayrulewidth\relax}}

\ifdefined\inthesis\else
\maketitle

\noindent The GUIDE-LLM checklist provides a standardized framework for reporting studies that use large language models (LLMs) in the behavioral and social sciences. It aims to promote transparency, reproducibility, and ethical accountability across all stages of LLM-based research.

\vspace{0.2cm}
\fi

{\footnotesize
\setlength{\tabcolsep}{3pt}
\renewcommand{\arraystretch}{1.0}

\noindent\textbf{Uses reported separately (Item A.1).} \guse{1}: the registered extraction and re-runs of the same instrument; \guse{2}: robustness comparisons of further extraction models (specification \S4f, \S4g, \S4k, \S4m); \guse{3}: instrument checks --- codebook review (\S4h), perturbation battery with paraphrasers, judges and sharpening agents (\S4i), mandatory-form variant (\S4j); \guse{4}: the language arm (\S4e(a) amendment); \guse{5}: AI assistants for literature, corpus collection, code and writing. ``Specification'' is the pre-analysis specification (\modelid{tools/docs/pre-analysis-specification.md}). No label of any registered figure comes from \guse{2}--\guse{5}.

\begin{longtable}{|J|K|}
\hline
\rowcolor{green!20}
\textbf{Scope of LLM use} & \textbf{Answer} \\*
\noalign{\hrule height \arrayrulewidth}% not \hline: longtable's \hline allows a page break below itself, which

\textbf{Item A.1:} LLMs were used in this project for:
\guideexpl{Briefly describe how and for what purposes LLMs were used in the study (research design, data processing, analysis, LLM as research object, participant-facing settings, communication). Depending on the use case, later items may have to be reported separately for each use case.}
&
\textbf{Analysis} (\guse{1}): extracting technology $\times$ frame element $\times$ evaluation $\times$ comparison target per paragraph under the frozen codebook v1.0; the LLM is the measuring instrument. The exposé's relevance step was never a model call (a deterministic keyword prefilter). \textbf{LLM as research object} (\guse{2}, \guse{3}): benchmarking further models on the task; model review of the codebook; perturbing the instrument, with model paraphrases rated by model judges; codebook-sharpening rules derived by model agents. \textbf{Data processing} (\guse{4}): machine translation of the 60 German evaluation paragraphs. \textbf{Research support and communication} (\guse{5}): coding assistants for literature work, code, drafting and revising the text, adversarial reviews; a claude.ai Artifact as the coding interface (versions in the thesis's \emph{Übersicht verwendeter Hilfsmittel}).
\\
\hline
\textbf{Item A.2:} Degree of automation (human-in-the-loop vs. fully automated):
\guideexpl{Indicate how much human oversight was involved, whether each output was reviewed, and who provided oversight.}
&
\guse{1}, \guse{2}, \guse{4}: fully automated; individual outputs are not reviewed (revised 2026-10-08): the 120 evaluation paragraphs, and on 2026-10-19 the 443 frame-completion and 55 fresh-sample paragraphs under the unchanged request. Humans author the codebook and code the reference and the blind re-code (Item E.1); one coder, no adjudication. Translations are never edited; all 60 pairs are checked before any agreement figure. \guse{3}: automated; paraphrases and review issues are accepted or matched by two blinded model judges; sharpening rules are recommendations, not adopted into the codebook. \guse{5}: interactive and agentic Claude Code and Codex CLI sessions, some of them autonomous handover sessions (e.g.\ the completion plan of 2026-09-05). How the output was checked and adopted: direct proofreading of every adopted passage, a cross-check by GPT-6 Astra, and comparison versions showing removed text in red and new text in blue.
\\[2pt]
& \textbf{Who performed the human coding} (reference, re-code, fresh sample), qualifications and training: the author only, for every set; trained by calibration on development material before evaluation coding.
\\
\hline
\rowcolor{green!20}
\textbf{Model/system details} & \textbf{Answer} \\*
\noalign{\hrule height \arrayrulewidth}% not \hline: longtable's \hline allows a page break below itself, which

\textbf{Item B.1:} Model name, including provider, model size, exact version/ID, date of access, and source link (if possible):
\guideexpl{Report exact model names (provider, version, date accessed); for local models a source link; if several models were tested, which one was used in the final study and why; roles consistent with Item A.1.}
&
\guse{1}: \textbf{Claude Sonnet 4.6} (Anthropic; proprietary, size not disclosed), \modelid{eu.anthropic.claude-sonnet-4-6}, Amazon Bedrock, \modelid{eu-central-1}; verified 2026-08-31 (\modelid{tools/pipeline/verifications/api_contract_20260831T170812Z.json}, 12 probes, plus two \modelid{cache_ttl_longgap} records); first certified run 2026-09-05. Defined in \modelid{tools/pipeline/instrument.py} (\texttt{MODELS}) and digested by the preflight gate. Selection W1 (Sonnet 4.6, Opus 4.6, Sonnet 5, Opus 5; registered 2026-09-06, when Sonnet 5 and Opus 5 were listed but denied): fork W1-a chosen 2026-09-06, lapsed 2026-09-15 unrun; the pin stands. W2 (Opus 4.6, registered 2026-09-15) closed 2026-09-24, reopened as \S4f.
\\[2pt]
& \guse{2}, \S4f (Bedrock, \modelid{eu-central-1}, accessed 2026-10-19): Claude Opus 4.6 \modelid{eu.anthropic.claude-opus-4-6-v1} (probed 2026-08-31), Claude Sonnet 4.5 \modelid{eu.anthropic.claude-sonnet-4-5-20250929-v1:0}, Claude Haiku 4.5 \modelid{eu.anthropic.claude-haiku-4-5-20251001-v1:0}; Amazon Nova Pro \modelid{eu.amazon.nova-pro-v1:0}; Qwen3-235B-A22B-Instruct-2507 \modelid{qwen.qwen3-235b-a22b-2507-v1:0} (Alibaba, open weights, 235B, 22B active); NVIDIA Nemotron 3 Super \modelid{nvidia.nemotron-super-3-120b} (open weights, about 120B). Unreachable that day: Claude 5 and 5.5 series, Opus 4.7/4.8; excluded (always return a reasoning block): gpt-oss-120b, MiniMax M2.5. \S4g: \textbf{Mercury 3} \modelid{mercury-3} (Inception Labs, proprietary diffusion model, size not disclosed), \modelid{https://api.inceptionlabs.ai/v1/chat/completions}, accessed 2026-10-30. \S4k: \textbf{Claude Sonnet 5} \modelid{eu.anthropic.claude-sonnet-5}, \textbf{Claude Opus 5} \modelid{eu.anthropic.claude-opus-5} (size not disclosed), Bedrock \modelid{eu-central-1}, accessed 2026-11-01 (denied 2026-09-06 and 2026-10-19); Sonnet 4.6 and Opus 4.6 as configuration controls (agent's addition, ratified 2026-11-02). \S4m: \textbf{Claude Opus 5.5} \modelid{eu.anthropic.claude-opus-5-5} (size not disclosed), Bedrock \modelid{eu-central-1}, accessed 2026-11-08; Opus 5 as configuration control under the same request and repeated under the \S4k request (agent's additions within the author's request).
\\[2pt]
& \guse{3}, 2026-10-30: \S4h Opus 4.6, Sonnet 4.6 (Bedrock EU) and Mercury 3; \S4i Sonnet 4.6, from 15:11 UTC via \modelid{global.anthropic.claude-sonnet-4-6}; paraphrasers Kimi K3 \modelid{global.moonshotai.kimi-k3} (Moonshot AI; the only profile the account can call) and Qwen3-235B (EU); judges Codex CLI (OpenAI; 0.158.0), \modelid{gpt-6-astra}, reasoning effort xhigh, read-only, and a Claude judge as Claude Code subagents on Claude Opus 5.5 (\modelid{claude-opus-5-5}) (both read from the local session records of that day); sharpening agents: Claude Code subagents, Claude Opus 5.5. \S4j: Sonnet 4.6 under v1.0F and v1.0r3. \S4l (2026-11-06): Sonnet 4.6 under v1.0C.
\\[2pt]
& \guse{4}: \textbf{gpt-oss-120b} \modelid{openai.gpt-oss-120b-1:0} (OpenAI, open weights, 120B, mixture of experts), Bedrock on demand in \modelid{eu-central-1}, accessed 2026-11-05: the only OpenAI model reachable without leaving the EU (proprietary OpenAI models only via \texttt{global.} profiles, not used); another developer than the extractor (registered Sonnet 4.6 request).
\\[2pt]
& \guse{5}: Claude Code (Anthropic), models Claude Fable 5, Opus 5, Fable 5.1, Opus 5.5 and Fable 5.2. Codex CLI (OpenAI), versions 0.151.0--0.158.0, was read-only. It used GPT-5.6 Sol \modelid{gpt-5.6-sol}, the CLI 0.151.0 default model, from 2026-08-29 to 2026-09-02 for prototype plan and design reviews (\modelid{tools/docs/PROTOTYPE-PLAN-REVIEW-LOG.md}). It then used \modelid{gpt-6-astra} from 2026-09-05. GPT-6 Astra \modelid{global.openai.gpt-6-astra} (OpenAI) ran on Amazon Bedrock's global inference profile with reasoning effort xhigh. From 2026-11-09 to 2026-11-12, it reviewed the whole thesis, debated cuts and re-code integration with Claude Code (three to five rounds), and contributed a measured cut pass. On 2026-11-24, it debated the changes following the supervisor's final feedback and the cut of the main text to below 70,000 characters with Claude Code (two and five rounds). It used read-only repository tools without corpus text, and none in the 2026-11-12 cut pass. Claude Opus 5.5 \modelid{eu.anthropic.claude-opus-5-5} ran on Bedrock's EU inference profile with effort high or xhigh from 2026-11-09. Its work included whole-thesis reviews and the final optimisation pass. Claude Fable 5.2 \modelid{global.anthropic.claude-fable-5-2} ran on Bedrock's global inference profile, possibly outside the EU, with effort high from 2026-11-10. It contributed reviews, debates and chapter drafts without receiving corpus text. Perplexity web search served online research, without corpus text. The overview of tools used itemises these activities. The claude.ai model that built the coding interface was Claude Opus~5.
\\
\hline
\textbf{Item B.2:} Model access (e.g., API, web interface, local) and context mode (e.g., chat mode or separate calls):
\guideexpl{How the models were accessed and whether in chat mode or stateless mode; the exact API name and version.}
&
\guse{1}--\guse{4}: \textbf{Amazon Bedrock Runtime \texttt{Converse} API} (\modelid{eu-central-1}) via \texttt{boto3}/\texttt{botocore} 1.43.83 (pinned in \modelid{pyproject.toml}, \modelid{uv.lock}); \textbf{stateless separate calls}, \texttt{build\_request()} sends one single-turn message holding one paragraph; no web interface, no local model. A \texttt{cachePoint} (\texttt{ttl: 1h}) after the system block changes transmission, not behaviour. Mercury 3 through Inception's OpenAI-compatible chat-completions API via a Converse-shaped adapter (\modelid{tools/pipeline/inception_client.py}), one stateless execution per paragraph (439 calls for 350 executions); adapter digest bound into later contracts, the four runs via commit \texttt{e7e049a}. Retry policy v3 (W3 development, v1.0r3): one correction call carrying paragraph, invalid call and validator messages, within one execution; feedback journalled privately. Judges and agents (\guse{3}) and \guse{5}: interactive or agentic Codex CLI and Claude Code sessions; also the Perplexity web interface, claude.ai for the coding-interface Artifact, and Amazon Bedrock for the reviewers named in B.1.
\\
\hline
\textbf{Item B.3:} Relevant LLM configurations reported (as applicable), such as temperature, max tokens, seed, and number of runs:
\guideexpl{Temperature, sampling parameters, penalties, stop sequences, number of completions or runs, quantization level, reasoning-related settings.}
&
\guse{1}: \texttt{build\_request()} (2026-09-02): \textbf{temperature 0}, no \texttt{topP} (rejected together); \textbf{no seed} (not offered; identical output not claimed); no reasoning mode, compute budget or quantization; \texttt{maxTokens} 1{,}024; no stop sequences, \texttt{stopReason == "tool\_use"} asserted; retry policy v2. \textbf{$k = 3$} executions on identical bytes (a convention), execution 1 primary, 1--3 for stability (registered run \texttt{a2\_sonnet\_20261013}, 2026-10-13; A3 run); frame-completion run $k = 1$. One English instrument, the paragraph in its original language (\texttt{INSTRUMENT-DESIGN.md} D6). v1.0r3 (A3, 2026-10-26): retry policy v3, $k = 3$.
\\[2pt]
& \guse{2}: \S4f the registered request, $k = 1$ on A2 and A3. \S4g \texttt{mercury3\_instant}: \texttt{reasoning\_effort} instant, \texttt{max\_completion\_tokens} 1{,}024; \texttt{mercury3\_low}: low, 8{,}192 (reasoning counts against it); temperature 0.5, the API minimum (0 refused, out-of-range reset to 1); no seed; $k = 1$; gate check 13 against \texttt{instrument.REQUEST\_ADAPTATION}. \S4k: temperature omitted (``\texttt{temperature} is deprecated for this model''; provider default), \modelid{additionalModelRequestFields: {"thinking": {"type": "disabled"}}}, otherwise pinned (1{,}024, cache point, policy v2); no seed; $k = 1$; controls under the same request; the tokenizer counts about 40\,\% more prompt tokens; no truncation. \S4m: Opus 5.5 rejects disabled thinking and a forced tool choice, so its request omits temperature, sets \modelid{additionalModelRequestFields: {"thinking": {"type": "adaptive"}, "output_config": {"effort": "low"}}} (also \texttt{"high"} as a contrast), \texttt{toolChoice auto} and \texttt{maxTokens} 8{,}192 (reasoning counts against it); otherwise pinned (cache point, policy v2); no seed; $k = 3$ at effort low, $k = 1$ at high; Opus 5 under the same request $k = 1$; Opus 5 under the \S4k request again, $k = 3$.
\\[2pt]
& \guse{3}: \S4h Opus and Sonnet 4.6 temperature 0, Mercury 0.5 instant and low; caps 8{,}192 (Mercury low 32{,}768), raised before the main runs after a smoke truncation at 8{,}192; 3 runs per cell, 72 runs; one identical retry. \S4i registered decoding, one execution per variant and paragraph (baseline three, most variants two); Kimi K3 (no temperature accepted), Qwen3 temperature 0.7; one call per text and set, at most three draws per text. \S4j v1.0F temperature 0, cap 2{,}048 (estimate about 1{,}000, longest 1{,}121; set from journalled answers, a disclosed deviation), policy v3; v1.0r3 1{,}024. \S4l v1.0C: the registered decoding (temperature 0, 1{,}024, policy v2), $k = 1$ on A2 and A3. Codex judge xhigh; Claude Code defaults.
\\[2pt]
& \guse{4}, translator: temperature 0, \texttt{maxTokens} 4{,}096, \texttt{reasoning\_effort} medium (reasoning discarded, never stored); one translation per paragraph, never edited or regenerated; up to two identical repeats on an exception, a stop other than \texttt{end\_turn} or empty text. Extractor: registered request on the stored translations, $k = 3$.
\\
\hline
\textbf{Item B.4:} Customization:
\guideexpl{Fine-tuning, RAG, automated prompt optimization, web search, agentic workflows, post-training refinements, other provider-specific features.}
&
\checked~Base model \quad \unchecked~Fine-tuning \quad \unchecked~RAG \quad \unchecked~Automated prompt optimization \quad \checked~Tool/function calling \quad \checked~Web search (\guse{5} only; none in \guse{1}--\guse{4}) \quad \checked~Agentic workflows (\guse{3} judges and agents, \guse{5}) \quad \unchecked~Other
\\[2pt]
& \guse{1}, \guse{2}, \guse{4}: \textbf{JSON-schema-constrained output through one forced tool} (\texttt{toolChoice}; 2026-09-02), \modelid{tools/pipeline/instrument/tool-config_v1.0.json} (sha256 \texttt{a1e82607f50e8c04}\ldots), generated from the codebook. Forcing does not guarantee a valid call (validation, Item E.2); the model retrieves and acts on nothing; one call, no loop. \S4f: Nova Pro, Qwen3, Nemotron without the cache point, which they refuse (gate-certified, no cache traffic); Claude models byte-identical. \S4g: one function with the pinned schema (\texttt{strict} not set), forced by name, no cache point; finish reason \texttt{stop} on a forced call read as tool use; malformed envelopes are response failures (hardened after two code reviews, 2026-10-30). \S4m: the one pinned tool offered with \texttt{toolChoice auto}, since Opus 5.5 rejects a forced choice; a reply without the call is a parse failure under the registered retry; a \texttt{reasoningContent} block beside the call is accepted for the thinking configurations only, never read or stored publicly (gate check 13 certifies the declared choice). W3 (development; seven probe calls on invented paragraphs, 2026-10-13): \texttt{strict: true}, nested shape, variant E an \texttt{explanation} field; neither eligible (lost 4 of 96 executions with 20 invalid first responses, and 1 with 9; v1.0 none), not carried forward. \guse{3}: \S4h forced tool \texttt{report\_codebook\_issues} (at most 15 issues); \S4j 17 booleans \texttt{element\_checks}, normalised to the registered list; judges and agents work in multi-step sessions with file access, no web search. \guse{5}: Claude Code web search and page fetching for literature and formal rules.
\\
\hline
\textbf{Item B.5:} Did the LLM session(s) include persistent memory across interactions?
\guideexpl{Whether the LLM could remember previous conversations; spillover effects.}
&
\guse{1}--\guse{4} pipeline calls: \unchecked~Yes \quad \checked~No \quad \unchecked~N/A --- no history, no memory feature; a policy-v3 correction stays within one execution. \guse{5} and the Claude Code sessions of \guse{3}: \checked~Yes --- Claude Code loads the repository's instruction files and a project memory directory.
\\
\hline
\rowcolor{green!20}
\textbf{Prompts} & \textbf{Answer} \\*
\noalign{\hrule height \arrayrulewidth}% not \hline: longtable's \hline allows a page break below itself, which

\textbf{Item C.1:} Exact prompt(s) reported:
\guideexpl{The exact text of prompts, including in-context examples; if not shareable, a representative example or a link.}
&
\guse{1}, \guse{2}: \modelid{tools/pipeline/instrument/system-prompt_v1.0.txt} (194 lines, sha256 \texttt{138cade3add5dcc4}\ldots; pinned 2026-09-02, unchanged since), \textbf{reproduced in full in the thesis appendix on the rendered instrument} (generated by \modelid{render_instrument_appendix.py}, 2026-10-29); the tool configuration (\modelid{tool-config_v1.0.json}, sha256 \texttt{a1e82607}\ldots) is in the electronic appendix. Rendered from \modelid{codebook-v1.0.yaml} (sha256 \texttt{1edb4601519417de}\ldots), both pinned in \modelid{pins.json}. No input-to-label demonstration; 16 quoted illustrations (15 English, one German), about 13 marking when a label fires; one phrase and the figure ``132 TWh'' derive from calibration paragraph \texttt{A011\_p02}, so article A011 is excluded from the frame; the deployment-growth illustration follows calibration paragraph \texttt{A010\_p15}, whose article stays in the frame (four paragraphs, two sampled, both removed by the clean-domain sensitivity); the German illustration paraphrases A002 (calibration only); the other illustrations are untraced. Exclusion rules: X1, X2, X5, X6 as written, X3, X4, X7, X11 transformed, X8--X10 excluded (\modelid{tools/pipeline/INSTRUMENT-DESIGN.md}; thesis Table of rule concordance); only the prompt prescribes \emph{mixed} for opposing voices. Development prompts, not reported: \modelid{system-prompt_v1.1S.txt}, \modelid{system-prompt_v1.1E.txt} (one demonstration on an invented paragraph), \modelid{system-prompt_v1.2S.txt} (W4: development F1 0.591 against 0.663, outcome ``retained''), each with its pins file. \guse{3}: \texttt{SYSTEM}, \texttt{ARM\_B\_NOTE} (\modelid{codebook_review.py}); \texttt{SYSTEM}, \texttt{SYSTEM\_ROUND2} (\modelid{battery_paraphrase.py}); battery variants hashed in the journal; \modelid{system-prompt_v1.0F.txt} (four replacements); judge and agent prompts in session workflow scripts, not in the repository.
\\[2pt]
& \guse{4}: fixed on two invented paragraphs: ``Translate the user's German news paragraph into English. Translate faithfully and completely: keep every statement, hedge, negation, number, unit and name, keep who says what (attribution and reported speech, e.g. render German reported speech with `said', `according to' or `reportedly'), and keep the order of information. Do not add, explain, summarise, correct or resolve ambiguity. Output only the English translation.''
\\[2pt]
& \guse{3}, \S4l v1.0C (2026-11-06): the registered prompt and schema; the user turn carries three lines before the unchanged paragraph block, ``Headline: \ldots'', ``Outlet: \ldots'' and ``Language: German'' or ``English'' (\modelid{tools/pipeline/instrument_v10c.py}).
\\
\hline
\textbf{Item C.2:} System-wide instructions (if any):
\guideexpl{System-level instructions that guide the model's general behavior.}
&
The task prompt is the system instruction (instrument or translation prompt); the user turn holds the paragraph alone (\S4l v1.0C: after the headline, outlet and language). No persona preamble; what the provider adds beneath is not visible and not claimed absent.
\\
\hline
\rowcolor{green!20}
\textbf{Data inputs \& privacy} & \textbf{Answer} \\*
\noalign{\hrule height \arrayrulewidth}% not \hline: longtable's \hline allows a page break below itself, which

\textbf{Item D.1:} Handling of personal or sensitive data (if any) (e.g., consent for data processing):
\guideexpl{Personal or sensitive data, consent, de-identification, provider safeguards, storage location, cross-border transfers.}
&
\textbf{Material:} paragraphs of published news articles, which name and quote persons as the outlets published them; no further personal data was collected and no special categories were sought. The corpus was assembled from online editions and the author's browser prints, under no database licence; the coding assistant (\guse{5}) searched for and retrieved batches 1--4 and transcribed batch 5, so all article text passed through the assistant, by a route not recorded; the author authorised sending it to Bedrock (\modelid{eu-central-1}). Complete article text is never committed; the exceptions (phrases in the prompt; short quotations of calibration paragraphs in the committed codebook files; free-text notes in six legacy label files) are in \modelid{tools/README.md} or named in the thesis appendix. No basis for storing, externally processing or redistributing the hand-collected texts is documented in the study records; the routing approvals are the author's own. The coding interface (\guse{5} tooling) is not archived in any build. Exported sheets with paragraph text were handled by the coding assistant (census pass 2, supplementary re-code) and local files with paragraph text were read by it in the checks of November 2026; inputs and routes of those sessions were not recorded, so further exposure to the assistant's provider is likely and not itemised. \textbf{Routes:} \texttt{eu.} profiles or in-region models, except the remaining \S4i calls from 2026-10-30, 15:11 UTC, via \modelid{global.anthropic.claude-sonnet-4-6} (possibly outside the EU) after the EU quota was exhausted, on the author's written authorisation, journalled per call; Kimi K3 received codebook text only, corpus quotes masked. Headlines (\S4l) are corpus text and went via the EU profile only. The 120 evaluation and 55 census paragraphs (\S4g) and the \S4h calibration and development paragraphs (67 occurrences, 64 distinct) went to Inception's API (hosting region not stated). Sharpening packets were read by Claude Code subagents; the public rules share no five-word run with any of the 563 paragraphs. \guse{4}: in-region. \guse{5}: on 2026-10-14 A3 text with provisional labels and the A2 labels went to a Claude Code session (consistency question), disclosed with the labels.
\\
\hline
\rowcolor{green!20}
\textbf{Validation \& interpretation} & \textbf{Answer} \\*
\noalign{\hrule height \arrayrulewidth}% not \hline: longtable's \hline allows a page break below itself, which

\textbf{Item E.1:} Human validation of LLM outputs:
\guideexpl{Whether and how human reviewers examined outputs, their roles and expertise, selection of reviewed outputs, training, criteria, disagreement resolution, reliability statistics, and use of feedback.}
&
\checked~Yes \quad \unchecked~No \quad \unchecked~N/A. \textbf{Who coded:} the author only (Item A.2).\par\smallskip
\guse{1}: Calibration comprises A0 and A0b (32 paragraphs, codebook drafts v0.1--v0.3); development comprises A1 (32 paragraphs, v1.0). The first development session on 1 September 2026 coded 35 paragraphs, including 6 off-topic paragraphs, and produced 53 scoring rows. Calibration and development are excluded from the evaluation. A2 contains 120 evaluation paragraphs, 60 German and 60 English, drawn from the frozen frame by language and outlet, with a 24-paragraph cap and design weights. The registered blind re-code (wave 1) covers 40 paragraphs, 20 per language, before the model run. The washout was registered at 14 days; the realised minimum was five days, with an evaluator minimum of four (deviation of 12 October 2026). This measures temporal stability, not intercoder reliability; a single coder was specified on 6 September 2026. Reference finalised 2026-09-30 (120 paragraphs) and 2026-10-08 (178 technology rows); re-code 2026-10-13. A first attempt displayed keyword-derived technology suggestions and was discarded; its export was opened only by the coder. No labels were carried over (2026-09-25). Metrics fixed 2026-09-05 (Chapter~5 of the thesis): design-weighted tuple micro-F1, technology detection, per-element F1, scores A and B, 2{,}000 article resamples, change rates, displacement. \textbf{Result:} tuples 0.643 [0.581, 0.705], technology 0.935, evaluation B 0.617 (A 0.874 on 208 units), comparison target B 0.742; 13 terminal failures of 360 executions as omissions; intra-coder $\alpha$ technology 0.78--1.00, elements 0.37--0.90, comparison target 0.62. Later article-resampling intervals for all eight reportable registered elements extend below 0.667 (\S4n (2), 2026-11-10). Pooled-cell $\alpha$ is 0.641 [0.555, 0.723] on shared rows and 0.619 [0.530, 0.697] on the full grid. The supplementary re-code (wave 2) covered the remaining 80 paragraphs on 10 November, at least 33 days after the first coding (\S4n (1)). Analyses using it, alone or pooled, are exploratory. First labels were hidden, but the coder knew the results; 25 paragraphs had direct item exposure through the re-reading of misses. Results are reported with and without those paragraphs. The agent set eight blank comparison targets on unframed rows to \texttt{none}, as in A2R. Combining waves gives shared-row pooled-cell $\alpha$ 0.711 [0.655, 0.762], mixing different intervals and exposures. Against the second coding, model precision and recall are 0.695 and 0.666; against the first, 0.759 and 0.557. The paired recall increase is $+0.109$ [$+0.056$, $+0.161$]. A smaller marking shift preceded the registered results; neither reading is established as correct. No borrowed benchmark or ceiling; below the more than 1{,}300 items that Song et al.'s simulation advises at an intercoder reliability of at least 0.7 (declared). \textbf{A3} (census of 55, 27/28, ten articles, \S4b): passes 2026-10-14 and 2026-10-19 (ten blank cells set to \texttt{none} by the agent on written instruction); reliability unmeasured; run after the reference commit \texttt{3f633bb}: technology 0.865, elements 0.563 [0.511, 0.622], 23 of 165 executions failed, 19 of 25 reportable displacements downward; v1.0r3 (20 invalid first responses, all recovered) 0.662; descriptive.
\\[2pt]
& \guse{2}: same references and evaluator (optional comparison item). \guse{3}: \S4h, nine gaps documented beforehand; found if both judges agree (95\,\% of 336 run--gap pairs, $\kappa = 0.88$); gaps found in 18 runs each (of 84): Opus 4.6 38, Sonnet 4.6 29, Mercury low 25, instant 18; none diffusion-only. Paraphrases: 35 of 90 accepted in round 1 (Codex 61, Claude 39), 55 regenerated once. \S4i on A2: AUC 0.62 [0.59, 0.66] (model--reference), 0.56 [0.52, 0.61] (first vs re-code); rules unvalidated. \S4j: element recall $+0.019$ [$-0.046$, $+0.078$] A2, $+0.079$ [$+0.031$, $+0.116$] A3: no consistent difference. \S4p (2026-11-11, existing data, no model call): first against second coding technology 0.942, elements 0.725, matched-frame evaluation 0.866 (model against first coding 0.935, 0.643, 0.874); 102 of 168 model misses absent from the second coding; keyword lexicon 0.961 against model 0.940 on five technologies, 0.895 against 0.935 on seven; precision above recall against the second coding in 10 of 12 configurations; German minus English technology detection on valid output $-0.016$ [$-0.077$, $+0.033$]. \S4o (2026-11-10, existing data, no model call): German minus English element micro-F1 $-0.032$ [$-0.145$, $+0.091$]; 12 of 13 terminal failures of the registered run violate the relational rules, none truncated; precision minus recall on valid first executions $+0.187$ (A2) and $+0.057$ (A3); adjusted investment share 33.1\% with the second coding as reference. \S4l (2026-11-06): v1.0C element recall $-0.019$ [$-0.067$, $+0.026$] A2, $+0.007$ [$-0.065$, $+0.113$] A3, no consistent difference. Against the blind re-code (40 paragraphs) the model reaches precision 0.662, recall 0.584 and F1 0.620, the first coding 0.602, 0.720 and 0.656; element $\alpha$ on the full row grid 0.619 (shared rows 0.641); on all 120 paragraphs, 86\,\% (design-weighted) of the missed elements lie within matched rows. \S4l (6), tallied 2026-11-09: the coder re-read 45 sampled misses within matched rows against the codebook, with the model's output visible. Counts: clear miss 16, borderline 10, reference possibly over-marked 15, cue outside the paragraph 2, assigned elsewhere 1, other 1. Two categories contradicted their notes; the coding agent (\guse{5}) set them to borderline on the author's instruction. The returned counts were 17 / 8 / 16 / 2 / 1 / 1. No reference label changed; the analysis is descriptive.
\\[2pt]
& \guse{4}: the author checks all 60 pairs before any agreement figure (sealed); primary: execution-1 element F1 translated minus the registered German 0.621 (reproduced before the first call), paired over 2{,}000 resamples of 28 articles. \textit{Result (2026-11-05):} the author checked all 60 pairs individually and flagged none (sealed); translated 0.599 against original 0.621, difference $-0.022$ [$-0.084$, $+0.032$], inconclusive; precision, recall and technology F1 differences include zero; 2 against 5 failed first executions; the translation changed the element sets of 43\,\% of matched units, against 3--6\,\% between native repeats.
\\
\hline
\textbf{Item E.2:} Describe any relevant post-processing (e.g., filtering in case of format mismatches, unit conversions, etc.):
\guideexpl{Cleaning or reformatting outputs, handling of inconsistent or unusable outputs, automated or manual corrections, exclusions.}
&
\textbf{Validation} (\modelid{run_extract.py}, policy v2): one \texttt{toolUse} block, pinned schema, relational rules. \texttt{refusal} terminal; \texttt{truncated}, \texttt{no\_tool\_use}, \texttt{schema\_invalid}, \texttt{contract\_invalid}, \texttt{unexpected\_content}, \texttt{unexpected\_stop\_reason}: one identical retry, then \emph{parse\_failed}; \texttt{throttled} six retries, \texttt{error} two (at most two journalled recovery episodes). Nothing is repaired; public files carry codes, not model text. Policy v3 retries with feedback; \S4j normaliser; \S4m thinking configurations: a \texttt{reasoningContent} block beside the one call is accepted and discarded (elsewhere \texttt{unexpected\_content}); \guse{4} text blocks joined and stripped, a missing translation is an omission. \textbf{Abstention:} \texttt{off\_topic} and an empty \texttt{elements} list are classes, not missing data. Codes mapped by script; frames derived from elements. \textbf{Funnel:} 95 articles $\to$ 74 with prose $\to$ 1{,}047 paragraphs $\to$ 696 candidates $\to$ 618 eligible $\to$ 120 sampled. Registered run: off topic 3 / 2 / 3, parse failures 5 / 5 / 3 per execution (omissions in execution 1, out of stability pairs). \textbf{Reporting:} cells need reference support $\geq s = 5$ (sealed 2026-09-07, commit \texttt{7c61b3d}, digest \texttt{4d4fbd87}\ldots); displacement in percentage points; exploratory DSL (\S4d) beside it (e.g.\ the investment element 37.7\,\% against the model's 19.4\,\%). \textbf{A3 embargo:} runs on A3 refused until both passes were committed (\modelid{A3_pass2_v1.0.csv}, \texttt{0257e3b0e081a766}\ldots).
\\
\hline
\rowcolor{green!20}
\textbf{Reproducibility} & \textbf{Answer} \\*
\noalign{\hrule height \arrayrulewidth}% not \hline: longtable's \hline allows a page break below itself, which

\textbf{Item F.1:} Code/\allowbreak notebooks/\allowbreak scripts for LLM calls shared:
\guideexpl{Whether code, prompts, logs or transcripts are shared, with sensitive information removed and a README.}
&
\unchecked~Yes \quad \checked~No \quad \unchecked~N/A. The repository stays private (author's decision of 2026-11-16).\par\smallskip
Prepared: pipeline code with README and tests; codebook, prompt, tool schema and pins; run, per-call (no free text), terminal and prediction files; sample, frame and class manifests; label files (identifiers, labels). Withheld: raw responses, private journals, article text, keys. Nothing is published while the repository stays private.
\\
\hline
\rowcolor{green!20}
\textbf{Competing interests} & \textbf{Answer} \\*
\noalign{\hrule height \arrayrulewidth}% not \hline: longtable's \hline allows a page break below itself, which

\textbf{Item G.1:} Funding, support, or other relevant relationships (including in-kind access to compute or models, or professional affiliations):
\guideexpl{Funding, in-kind access, affiliations, investments or family relationships with entities that have a financial interest in LLMs.}
&
\checked~Yes \quad \unchecked~No. In-kind access: Amazon Bedrock use (all Bedrock models listed in B.1) was paid from AWS promotional credits. Claude Code and OpenAI Codex were paid privately by the author, as business plans chosen for stronger data protection; a data-processing agreement is in place with AWS. The author also paid privately for the Inception API (Mercury~3). No other funding, affiliation or relationship with AWS, Anthropic, OpenAI or Inception exists.
\\
\hline
\rowcolor{yellow!30}
\textbf{Optional items} & \textbf{Answer} \\*
\noalign{\hrule height \arrayrulewidth}% not \hline: longtable's \hline allows a page break below itself, which

\textbf{Justification for the LLM choice}
\guideexpl{Performance, transparency, reproducibility, ethical considerations, cost or ease-of-use.}
&
\unchecked~Performance: no criterion (W1 lapsed unrun). \checked~Transparency: proprietary, undisclosed size; open weights only in \guse{2}. \checked~Reproducibility: pinned ID, request digest and gate certificate per run; no seed. \checked~Ethics: EU route for every registered call. \checked~Cost: see below.
\\
\hline
\textbf{Rationale for the prompt design}
\guideexpl{Structured format, guidelines or literature, iteration, few-shot selection, standardization across models.}
&
The frozen codebook renders both the human instrument and the prompt, so coder and model get the same definitions: exhaustive descriptions, no demonstrations, explicit abstention states, schema-forced output. No automated optimisation; frozen before any evaluation label existed; identical across \guse{1} and \guse{2}.
\\
\hline
\textbf{Comparison against other methods/LLMs}
\guideexpl{How comparisons were conducted, metrics, fairness of inputs, and exclusions.}
&
\checked~Yes. Comparisons use the registered evaluator's quantities and paired differences under common article resamples. Better or worse requires the paired 95\,\% element-F1 difference interval to exclude zero on both samples (\S4g, \S4k). Requests differ only through declared adaptations. First executions are reported; repeated executions are identified in Item B.3. The prompt was developed with Claude runs. \textbf{Element micro-F1, A2 / A3:} Opus 4.6 0.674 / 0.680, Sonnet 4.5 0.578 / 0.623, Haiku 4.5 0.413 / 0.529, Nova Pro 0.350 / 0.416, Qwen3 0.515 / 0.536, Nemotron 0.425 / 0.422, Mercury instant 0.428 / 0.442, low 0.580 / 0.684, Sonnet 5 0.641 / 0.650, Opus 5 0.740 / 0.684 (controls Sonnet 4.6 0.622 / 0.604, Opus 4.6 0.667 / 0.682; registered 0.643 / 0.563); \S4m Opus 5.5 0.726 / 0.774 (execution 1 of 3; control Opus 5 under the 5.5 request 0.744 / 0.700; effort high 0.709 / 0.824; Opus 5 repeated 0.741 / 0.662). Mercury instant worse than Opus 4.6; Opus 5 better than Sonnet 4.6 on both samples (A3 margin rests on that run's failed first executions), the only model without precision distinguishably above recall on A2 (0.75 / 0.73) and with recall above precision on A3 (0.64 / 0.74), reproduced in three further executions on 2026-11-08; Opus 5.5 better than Sonnet 4.6 and Opus 4.6 on both samples, no consistent difference against Opus 5 or its control, precision above recall on A2 (0.81 / 0.66); effort high against low no consistent difference (recall lower on A2, precision higher on A3); all other pairings no consistent difference; no selection. Prompt variants: W3 and W4 on the development split only; the later variants (\S4i paraphrases, \S4j v1.0F, \S4l v1.0C) ran on inspected samples, descriptive and without selection.
\\
\hline
\textbf{Training data leakage}
&
\unchecked~Yes \quad \checked~No \quad \unchecked~N/A: not examined; the corpus is published online news; A3 controls exposure to this study, not to pre-training.
\\
\hline
\textbf{Risk of bias}
&
\checked~Yes: tuples German 0.621 [0.532, 0.710] vs English 0.653 [0.569, 0.731], technology 0.884 vs 0.960; 11 of 13 parse failures German. Language is collinear with country and outlet; no mitigation beyond reporting (motivation of \guse{4}).
\\
\hline
\textbf{Conversation transcripts}
&
\unchecked~Yes \quad \unchecked~No \quad \checked~N/A for \guse{1}--\guse{4} (single journalled calls). \guse{5} transcripts are not in the repository.
\\
\hline
\textbf{Ethical implications}
&
No human participants; published news; licensing and routing in Item D.1.
\\
\hline
\textbf{Computational resources}
&
USD estimated from observed tokens at provisional list rates (run records \modelid{tools/runs/}; Bedrock price map of 2026-08-14); probes excluded; no GPU hours. Dry run 2026-09-05: 96 calls, 48{,}660 input, 12{,}353 output, 414{,}624 cache-read tokens, USD 0.50 (superseded run 0.51; re-check 2026-10-09 0.51). Registered run: 379 attempts, 0 throttles, 200{,}646 input, 59{,}296 output, 1{,}636{,}901 cache-read tokens, USD 2.18, about 12 minutes (cumulative 3.70). W3 1.85; W4 1.08; frame completion 2.76; A3 189 calls, 1.09; \S4f about 4.2 (+0.13 superseded); v1.0r3 185 calls, 1.07 (cumulative about 15.9); \S4g 439 calls, about 0.34; \S4h 3.97; \S4i 19{,}206 executions, 22{,}643 attempts, 121.01 (\S4i + \S4j about 126); \S4k 5.07; \S4m 14.98 for the eight reported runs (1{,}400 executions, 1{,}413 calls) plus 2.19 for four superseded journals; \guse{4} translation 61 calls, 14{,}709 / 29{,}542 tokens, 0.02, extraction 188 calls, 1.04. Sum about USD 169; \guse{5} subscriptions not measured.
\\
\hline
\end{longtable}
}

\ifdefined\inthesis\else
\noindent
\fcolorbox{gray}{gray!20}{
  \parbox{.95\linewidth}{
    \textbf{How to cite:} Feuerriegel et al. (2026). \textit{A reporting checklist for large language models in behavioural science}. In: \textit{Nature Human Behaviour}. DOI: 10.1038/s41562-026-02492-7
  }
}
\fi

\end{document}

\endgroup
\addtocounter{table}{-1}
\endgroup


\newpage
\section{Electronic Appendix}\label{app:electronic}
The electronic appendix is the submission repository, a copy of the project repository at the commit named below without its version history. It holds the material behind every figure of this thesis and the tables not printed in Appendix~A. Table~\ref{tab:electronic} maps each retained claim to its files; paths are relative to the repository root, and the README describes the set-up, the pipeline and the reproduction check.

The files support two levels of reproduction. Every reported figure can be recomputed from the committed run records, labels and manifests without a model call; most generators re-render and compare their floats in a \texttt{-{}-check} mode. Re-running an extraction needs the article text, which is not in the repository, and access to Amazon Bedrock, and it may return different outputs, because the provider offers no seed.

\begin{description}
\item[Commit] {\small\texttt{0b77ebdc831508fe8c5cd4ab1b517899841c4643}} of the private project repository. The submitted PDF is built from this state; only this identifier was added afterwards.
\item[Withheld] Article and paragraph text, raw model responses, private journals and coding sheets, the stored translations, credentials. The texts are withheld as a precaution: the basis for storing, externally processing and redistributing the hand-collected articles was not clarified. The coding interface, a browser tool built outside the repository, is not archived in any build; its behaviour is documented by its written specification only.
\end{description}

\begingroup\footnotesize
\setlength{\tabcolsep}{4pt}
\begin{longtable}{@{}>{\raggedright\arraybackslash}p{48mm}>{\raggedright\arraybackslash}p{26mm}>{\raggedright\arraybackslash}p{78mm}@{}}
\caption[Index of the electronic appendix]{Index of the electronic appendix: retained claims and material, where the thesis uses them, and the repository files that back them.}\label{tab:electronic} \\
\toprule
\textbf{Claim or material} & \textbf{In the thesis} & \textbf{Repository files} \\
\midrule
\endfirsthead
\toprule
\textbf{Claim or material} & \textbf{In the thesis} & \textbf{Repository files} \\
\midrule
\endhead
\bottomrule
\endfoot
Pipeline: gate, runner, validator, evaluator, estimators and tests; plans and review logs & Chapters~\ref{sec:system_design} and~\ref{sec:methodology} & \modelid{tools/pipeline/}, \modelid{tools/README.md}, \modelid{tools/docs/} \\
Codebook versions v0.1 to v1.0 & Table~\ref{tab:codebook_elements} & \modelid{tools/annotation/codebook/} \\
Rendered instrument: system prompt (printed), tool configuration and pins (not printed; tool configuration sha256 \texttt{a1e82607}\ldots) & Appendix~\ref{app:instrument} & \modelid{tools/pipeline/instrument/system-prompt_v1.0.txt}, \modelid{tools/pipeline/instrument/tool-config_v1.0.json}, \modelid{tools/pipeline/pins.json}, \modelid{paper/tables/instrument_appendix.tex} \\
Corpus metadata (no text) & Table~\ref{tab:corpus_facts} & \modelid{corpus/manifest.csv} \\
Figures not printed: corpus paragraphs by publication month; element support in the sessions before the evaluation sample & Appendices~\ref{app:corpus_details}, \ref{app:reference_disclosures} & \modelid{paper/images/corpus_timeline.tex}, \modelid{paper/images/element_support.tex}, generated by \modelid{tools/annotation/descriptives.py} \\
Registry and split assignments & Table~\ref{tab:corpus_facts} & \modelid{tools/annotation/split-assignment.json} \\
Requirements of the topic description: status, reason and place & Appendix~\ref{app:topic_requirements} & the files indexed for the tables it cites \\
Output records and aggregation keys; registry corpus versions & Appendix~\ref{app:data_dictionary} & \modelid{tools/runs/<run_id>_predictions.csv}, \modelid{<run_id>_terminal.csv}, \modelid{corpus/manifest.csv}, \modelid{tools/annotation/archive-manifest.json} \\
Frame, sample, re-code, census and completion manifests; threshold seal & Figure~\ref{fig:sample_flow}, Table~\ref{tab:corpus} & \modelid{tools/annotation/labels/A2_frame_manifest.json}, \modelid{A2_sample_manifest.json}, \modelid{A2R_manifest.json}, \modelid{A3_sample_manifest.json}, \modelid{A2_frame_completion_manifest.json}, \modelid{A2_threshold_seal.json} \\
First coding: registered evaluation reference, 120 paragraphs (identifiers and labels only) & Table~\ref{tab:primary} & \modelid{tools/annotation/labels/A2_pass1_v1.0.csv}, \modelid{A2_pass2_v1.0.csv} \\
Registered re-code: 40 paragraphs, wave 1 & Table~\ref{tab:reliability} & \modelid{tools/annotation/labels/A2R_pass1_v1.0.csv}, \modelid{A2R_pass2_v1.0.csv}, \modelid{A2R_manifest.json} \\
Supplementary re-code: 80 paragraphs, wave 2; exploratory analyses & Tables~\ref{tab:alpha_waves} and~\ref{tab:recode_waves} & \modelid{tools/annotation/labels/A2R2_pass1_v1.0.csv}, \modelid{A2R2_pass2_v1.0.csv}, \modelid{A2R2_manifest.json} \\
Other reference labels (identifiers and labels only), sessions and deviation records & Table~\ref{tab:disclosures} & \modelid{tools/annotation/labels/} (\modelid{*_pass1_*.csv}, \modelid{*_pass2_*.csv}, \modelid{sessions.csv}, \modelid{README.md}) \\
Run records, gate records, per-call, terminal and prediction files & all results & \modelid{tools/runs/<run_id>.json}, \modelid{preflight_<run_id>.json}, \modelid{<run_id>_calls.jsonl}, \modelid{_terminal.csv}, \modelid{_predictions.csv} \\
Registered metrics, unweighted figures and all evaluator tables & Table~\ref{tab:primary}, Appendix~\ref{app:detailed_results} & \modelid{tools/runs/tables/a2_sonnet_20261013/} \\
Registered tables not printed: run record (tokens, gate, environment); per-class F1 of every dimension, incl.\ comparison-target classes; both categorical scores with excluded units; classes defined across resamples; clean-domain sensitivity & Appendix~\ref{app:detailed_results} & \modelid{paper/tables/tab_run_record.tex}, \modelid{tab_per_class.tex}, \modelid{tab_categorical.tex}, \modelid{tab_limits.tex}, \modelid{tab_clean_domain.tex}; the evaluator's output files \modelid{table2_per_class.csv}, \modelid{table3_categorical.csv}, \modelid{table6_limits.csv} and \modelid{table7_clean_domain.csv} in \modelid{tools/runs/tables/a2_sonnet_20261013/} \\
Adjusted prevalences, incl.\ the classes below the reporting threshold & Figure~\ref{fig:dsl}, Appendix~\ref{app:dsl_details} & \modelid{tools/runs/tables/dsl/}, \modelid{paper/tables/tab_dsl.tex} \\
Fresh census and candidate, incl.\ the replicate tables not printed and the agreement by dimension and language & Appendix~\ref{app:fresh} & \modelid{tools/runs/tables/a3_sonnet_20261019/}, \modelid{a3_candidate/}, \modelid{paper/tables/fresh_sample/} (\modelid{tab_primary.tex} and the replicates), \modelid{paper/images/fresh_sample/} \\
Model comparison: results table with the configuration controls, Opus~5.5 table, complete and per-element tables (not printed) & Appendix~\ref{app:model_comparison} & \modelid{tools/runs/tables/model_comparison/}, \modelid{model_comparison_diffusion/}, \modelid{model_comparison_gen5/}, \modelid{paper/tables/tab_model_comparison.tex}, \modelid{tab_opus55.tex}, \modelid{tab_model_comparison_full.tex}, \modelid{tab_model_comparison_elements.tex} \\
Translated input: comparison, translation record and check seal; detailed table and figure of the four differences (not printed) & Appendix~\ref{app:language} & \modelid{tools/runs/tables/language_arm/}, \modelid{tools/runs/language_arm/}, \modelid{paper/tables/tab_language_arm.tex}, \modelid{paper/images/fig_language_arm.tex} \\
Perturbation battery: analysis, accepted paraphrases, sharpening rules; design and conduct record, results by variant and paraphrase reading per unit (not printed) & Appendix~\ref{app:battery} & \modelid{tools/runs/battery/}, \modelid{paper/tables/tab_battery_design.tex}, \modelid{tab_battery_families.tex}, \modelid{tab_battery_units.tex} \\
Mandatory form variant & Appendix~\ref{app:form_variant} & \modelid{tools/runs/tables/form_variant/} \\
Checks of 10 and 11 November (\S4o, \S4p): language contrast, failures by code, valid-output agreement, per-element intervals, adjustment under the second coding; the coder's retest, keyword and majority comparators, configurations against the second coding, the no-frame decision & Appendix~\ref{app:later_checks} & \modelid{tools/runs/tables/review_checks_4o/}, \modelid{review_checks_4p/}; \modelid{tools/pipeline/review_checks_4o.py}, \modelid{review_checks_4p.py}; \modelid{tools/annotation/labels/A2_prefilter_matches.json} (keyword matches, no text) \\
Later checks: blind re-code comparison, $\alpha$ on the full grid, the coder's $\alpha$ without the 25 re-read paragraphs (key \modelid{recode\_wave2}, not printed), missed elements by cause, comparison-target baseline; context-parity run (instrument v1.0C) & Appendix~\ref{app:later_checks} & \modelid{tools/runs/tables/review_analyses/}, \modelid{tools/runs/tables/context_variant/}; \modelid{tools/pipeline/review_analyses.py}, \modelid{compare_context.py}, \modelid{instrument_v10c.py}, \modelid{missed_elements_sheet.py} \\
Development iterations: seals, selections, full table & Appendix~\ref{app:development} & \modelid{tools/runs/w3_seal.json}, \modelid{w3_selection.json}, \modelid{w4_seal.json}, \modelid{w4_selection.json}, \modelid{paper/tables/tab_development.tex} \\
Codebook review: documented gaps, runs, judge matches, full table & Appendix~\ref{app:codebook_review} & \modelid{tools/runs/codebook_review/}, \modelid{paper/tables/tab_codebook_review.tex} \\
Headline by dimension and language as a figure; former allocation table & --- & \modelid{paper/images/fig_primary.tex}, \modelid{paper/tables/tab_a2_allocation.tex} \\
Pre-analysis specification with dated amendments & Appendix~\ref{app:protocol} & \modelid{tools/docs/pre-analysis-specification.md} \\
GUIDE-LLM source & Appendix~\ref{app:guide_llm} & \modelid{paper/GUIDE-LLM.tex} \\
Generators of the tables and figures (most with a \texttt{-{}-check} mode that re-renders and compares) & all floats & \modelid{tools/pipeline/figures.py}, \modelid{compare_models.py}, \modelid{error_breakdown.py}, \modelid{dsl.py}, \modelid{language_arm.py}, \modelid{battery_analyse.py}, \modelid{compare_form.py}, \modelid{codebook_review.py}, \modelid{render_instrument_appendix.py}, \modelid{review_analyses.py}, \modelid{compare_context.py}; \modelid{tools/annotation/descriptives.py} \\
\end{longtable}
\endgroup

\newpage
\section*{Übersicht verwendeter Hilfsmittel}
\addcontentsline{toc}{section}{Übersicht verwendeter Hilfsmittel}
\RaggedRight
This is the overview of tools used (\foreignlanguage{ngerman}{Übersicht verwendeter Hilfsmittel}) to which the statutory declaration at the end of this thesis refers. It names all generative models used for this thesis by product name and version, and states how, to what extent and for what purpose each was used. Generated wording is not marked passage by passage in the text. Its origin is declared here, and the author's statement after the table says how adopted passages were checked. Part~A lists the models that served as research instruments in the extraction pipeline; their configuration, prompts, data handling and validation are reported in the GUIDE-LLM checklist (\guideref) and in Chapters~\ref{sec:system_design} to~\ref{sec:results}. Part~B lists the AI assistants used in producing the thesis text, the code and the corpus; the author's own statement follows the table.

{\footnotesize
\setlength{\tabcolsep}{4pt}
\renewcommand{\arraystretch}{1.15}
\begin{longtable}{@{}>{\RaggedRight\arraybackslash}p{0.15\textwidth}>{\RaggedRight\arraybackslash}p{0.25\textwidth}>{\RaggedRight\arraybackslash}p{0.14\textwidth}>{\RaggedRight\arraybackslash}p{0.405\textwidth}@{}}
\toprule
\textbf{Tool} & \textbf{Version} & \textbf{Provider} & \textbf{Purpose and extent} \\
\midrule
\endfirsthead
\toprule
\textbf{Tool} & \textbf{Version} & \textbf{Provider} & \textbf{Purpose and extent} \\
\midrule
\endhead
\midrule
\multicolumn{4}{r@{}}{\textit{continued on the next page}} \\
\endfoot
\bottomrule
\endlastfoot
\multicolumn{4}{@{}l}{\textbf{A. Research instruments in the extraction pipeline}} \\*
Registered extraction model &
Claude Sonnet 4.6, \modelid{eu.anthropic.claude-sonnet-4-6}; Amazon Bedrock &
Anthropic &
The measuring instrument: codes technologies, frame elements, evaluations and comparison targets per paragraph under the frozen codebook v1.0 (Section~\ref{subsec:extraction_call}, Chapter~\ref{sec:methodology}). 360 registered executions on the 120 evaluation paragraphs; further runs on 443 and 55 paragraphs and in the later analyses (Section~\ref{subsec:added_methods}). \\
\addlinespace
Comparison models (Anthropic) &
Claude Opus 4.6 (\modelid{eu.anthropic.claude-opus-4-6-v1}), Claude Sonnet 4.5 (\modelid{eu.anthropic.claude-sonnet-4-5-20250929-v1:0}), Claude Haiku 4.5 (\modelid{eu.anthropic.claude-haiku-4-5-20251001-v1:0}); Amazon Bedrock &
Anthropic &
Robustness comparison of the extraction (Section~\ref{subsec:recurrence}, Appendix~\ref{app:model_comparison}); no model was selected and no registered figure comes from them. One execution per paragraph on the 120 evaluation and the 55 fresh-census paragraphs. \\
\addlinespace
Comparison models (other providers) &
Amazon Nova Pro (\modelid{eu.amazon.nova-pro-v1:0}); Qwen3-235B-A22B-Instruct-2507 (\modelid{qwen.qwen3-235b-a22b-2507-v1:0}); NVIDIA Nemotron 3 Super (\modelid{nvidia.nemotron-super-3-120b}); Amazon Bedrock &
Amazon; Alibaba (open weights); NVIDIA (open weights) &
As in the previous row. \\
\addlinespace
Claude 5 models &
Claude Sonnet 5 (\modelid{eu.anthropic.claude-sonnet-5}), Claude Opus 5 (\modelid{eu.anthropic.claude-opus-5}), Claude Opus 5.5 (\modelid{eu.anthropic.claude-opus-5-5}); Amazon Bedrock &
Anthropic &
Robustness comparison with Claude Sonnet 4.6 and Claude Opus 4.6 as configuration controls, and Opus 5.5 with Opus 5 under its request as control (Section~\ref{subsec:recurrence}, Appendix~\ref{app:model_comparison}); one execution per paragraph on the evaluation sample and the fresh census, three for Opus 5.5 and an Opus 5 repeat. \\
\addlinespace
Diffusion model &
Mercury 3 (\modelid{mercury-3}), configurations \emph{instant} and \emph{low}; Inception API &
Inception Labs &
Robustness comparison of the extraction and, second, the codebook review (next row); 439 extraction calls, one execution per paragraph and configuration on the evaluation sample and the fresh census. \\
\addlinespace
Codebook review &
Claude Opus 4.6, Claude Sonnet 4.6 and Mercury 3 as above &
Anthropic; Inception Labs &
Retrospective benchmark: finding gaps in the codebook versions v0.1, v0.2 and v1.0 (Appendix~\ref{app:codebook_review}); 72 review runs. No label of any reported figure comes from this use. \\
\addlinespace
Paraphrasers &
Kimi K3 (\modelid{global.moonshotai.kimi-k3}); Qwen3-235B-A22B-Instruct-2507 as above; Amazon Bedrock &
Moonshot AI; Alibaba &
Rewording the codebook definitions for the perturbation battery (Appendix~\ref{app:battery}); one call per definition text and set, at most three draws per text. Codebook text only; corpus quotations masked for Kimi K3. \\
\addlinespace
Perturbation battery and form variant &
Claude Sonnet 4.6 as above; later through the global profile \modelid{global.anthropic.claude-sonnet-4-6} &
Anthropic &
The registered request with one declared change at a time, and the mandatory-form variant v1.0F (Appendices~\ref{app:battery} and~\ref{app:form_variant}); 19,206 battery executions. No label of any registered figure comes from either. \\
\addlinespace
Judges &
Codex CLI 0.158.0, model \modelid{gpt-6-astra} (part~B); Claude Code workflow subagents, model Claude Opus 5.5 &
OpenAI; Anthropic &
Blinded matching of model-reported issues to documented codebook gaps (Appendix~\ref{app:codebook_review}) and blinded judgement of whether a paraphrase keeps a definition's meaning (Appendix~\ref{app:battery}); 336 run--gap pairs judged by both ($\kappa = 0.88$) and 90 paraphrase candidates. \\
\addlinespace
Sharpening agents &
Claude Code workflow subagents, model Claude Opus 5.5 &
Anthropic &
Deriving candidate rules for sharpening the codebook from the battery's results (Appendix~\ref{app:battery}); one reader, one adversarial verifier and one merging agent per unit. The rules are not validated on new labels. \\
\addlinespace
Translator (language comparison) &
gpt-oss-120b (\modelid{openai.gpt-oss-120b-1:0}); Amazon Bedrock &
OpenAI (open weights) &
Translating the 60 German paragraphs of the evaluation sample into English (Appendix~\ref{app:language}); one translation per paragraph, followed by three executions of the registered extraction model on the stored translations (188 calls). \\
\addlinespace
Access probes &
MiniMax M2.5 and Mistral Pixtral Large; Amazon Bedrock &
MiniMax; Mistral AI &
Compatibility probes on invented paragraphs, without article text, on 19 October 2026. MiniMax M2.5 always returned a reasoning block that the registered response contract refuses, and Pixtral Large rejected the forced tool choice, so both were excluded from the model comparison. No reported figure comes from this use. \\
\addlinespace[4pt]
\multicolumn{4}{@{}l}{\textbf{B. AI assistants used in producing the thesis, the code and the corpus}} \\*
Claude Code &
Models Claude Fable 5, Claude Opus 5, Claude Fable 5.1, Claude Opus 5.5 and Claude Fable 5.2 &
Anthropic &
Writing and revising the pipeline and annotation code and its tests; drafting and revising the thesis text, tables and figure code; literature search and organisation; corpus collection: searching for, selecting and retrieving the 35 articles of batches 1--4 against the collection brief, and screening and transcribing the author's browser prints of batch 5; review workflows with subagents; drafting this overview. One consistency question during the first coding pass of the fresh sample. \\
\addlinespace
OpenAI Codex &
Codex CLI 0.151.0 to 0.158.0; models GPT-5.6 Sol (\modelid{gpt-5.6-sol}) until 2026-09-02, then \modelid{gpt-6-astra} &
OpenAI &
Adversarial review of the prototype plan, research-design decisions, plans, code and thesis text in a read-only sandbox; blinded judge (part~A). \\
\addlinespace
OpenAI GPT on Amazon Bedrock &
GPT-6 Astra (\modelid{global.openai.gpt-6-astra}), reasoning effort xhigh, read-only repository tools; Amazon Bedrock, global inference profile &
OpenAI &
Critical review in debates with Claude Code, on the author's decisions: the proposals for shortening the main text, the integration of the coder's second re-code and the $\alpha$ intervals, and the changes following the supervisor's final feedback. Reviews of the whole thesis: first-reader and expert passes, cross-checks of the Claude models' findings and proposals, grading passes on a pre-set scheme and a check of every cited source against its PDF. Independent proposals for wording and cuts, and consistency checks. Claude Code checked every finding against the files, measured the cuts in the built PDF and alone edited the files. No corpus text was sent. \\
\addlinespace
Claude on Amazon Bedrock &
Claude Opus 5.5 (\modelid{eu.anthropic.claude-opus-5-5}), effort high or xhigh; Amazon Bedrock, EU inference profile. Claude Fable 5.2 (\modelid{global.anthropic.claude-fable-5-2}), effort high; Amazon Bedrock, global inference profile &
Anthropic &
Opus 5.5: independent reviews of the whole thesis (first-reader, expert, cross-check, grading and source passes); the register of findings; proposals on form, style and cuts; reader and consistency checks; the final optimisation pass of the thesis text under the author's brief, in Claude Code workflows of subagents. Fable 5.2: proposals on method, statistics and theory, the Claude side of cross-reviews and debates, grading passes, and drafts of revised chapters. Fable 5.2 received no corpus text. The findings were checked against the files before any change. \\
\addlinespace
Perplexity &
Perplexity web search (perplexity.ai) &
Perplexity AI &
Online research. No corpus text was sent. \\
\addlinespace
Coding interface &
Web coding interface built as a claude.ai Artifact; built with Claude Opus~5 on claude.ai; source outside the repository &
Anthropic &
Displaying the paragraphs and recording the reference codes (Section~\ref{subsec:annotation}). The build used has no suggestion, highlighting, default or pre-selection function. An earlier build with keyword-derived technology suggestions was withdrawn and its attempt discarded; a development-set sheet carrying model labels was quarantined and enters no figure (Appendix~\ref{app:reference_disclosures}). \\
\addlinespace
Other software (not generative) &
TeX Live 2025 (pdfLaTeX, biber, biblatex-apa, TikZ/pgfplots); Python 3.13 (pipeline) and 3.14 (annotation scripts); Git &
Open-source projects &
Typesetting, figures, pipeline and analysis code, version control; own scripts for checking citations against the source PDFs. Grammarly occasionally for spelling. No reference manager or translation tool was used for writing the thesis. \\
\end{longtable}
}

\textbf{Statement of extent by the author.} I wrote the basis of the thesis text largely myself, partly as full drafts and partly as bullet points. In several iterations, the assistants in part~B shortened passages, formulated bullet points as prose and helped me explore lines of argument. The final text was written almost exclusively by me and Claude. I checked every adopted passage by proofreading it directly, had GPT-6 Astra cross-check it, and read revised versions as comparison documents showing removed text in red and new text in blue.

\textbf{How, why and with what experience.} I used the assistants to brainstorm the structure of the thesis and its methods, to formulate and check text, to implement changes and correct errors, to analyse parts of the thesis, to aggregate information from the literature research and to organise it in an LLM-maintained wiki. I chose them because they let me check results quickly and rebuild the thesis quickly when results came in. Keeping information consistent was difficult despite the wiki, small tasks sometimes used many tokens, and the generated wording was at times too mechanical and hard to read.

\newpage
\section*{Eidesstattliche Erklärung}
\addcontentsline{toc}{section}{Eidesstattliche Erklärung}
\begin{otherlanguage}{ngerman}
\emergencystretch=3em
\noindent
Hiermit versichere ich, dass ich die vorliegende Arbeit selbständig und ohne Benutzung anderer als der angegebenen Hilfsmittel angefertigt, noch nicht einer anderen Prüfungsbehörde vorgelegt und noch nicht veröffentlicht habe.

\noindent
Im Falle der Nutzung von generativen Modellen zur Erstellung von Texten, Abbildungen, Berechnungen und anderen Leistungen verantworte ich die Auswahl, Übernahme und sämtliche Ergebnisse des von mir verwendeten generierten Outputs vollumfänglich selbst. Im Verzeichnis \glqq Übersicht verwendeter Hilfsmittel\grqq{} habe ich alle verwendeten generativen Modelle mit ihrem Produktnamen benannt und angegeben, wie, in welchem Umfang und zu welchem Zweck diese benutzt wurden.

\noindent
\foreignlanguage{english}{I am aware that this thesis in digital form can be examined for the use of unauthorised aid and in order to determine whether the thesis as a whole or parts incorporated in it may be deemed to constitute plagiarism. For the comparison of my work with existing sources, I agree that it shall be entered in a database where it shall also remain after examination, to enable comparison with future theses submitted. Further rights of reproduction and usage, however, are not granted here.}

\vspace{2cm}

\noindent
\begin{minipage}[t]{0.42\textwidth}
München, den \mycompletiondate\par
\vspace{0.1cm}
\rule{\linewidth}{0.4pt}\par
(Datum)
\end{minipage}\hfill
\begin{minipage}[t]{0.42\textwidth}
\mbox{}\par
\vspace{0.1cm}
\rule{\linewidth}{0.4pt}\par
(Unterschrift)\\
\myname
\end{minipage}
\end{otherlanguage}

\vspace{2.5cm}

\noindent
{\small
\textit{English translation as printed on the ISC form:}

\noindent
I hereby declare that I have prepared this thesis independently and without the use of aids other than those specified, that I have not yet submitted it to another examination authority and that it has not yet been published.

\noindent
In the case of the use of generative models for the creation of texts, illustrations, calculations and other services, I am fully responsible for the selection, adoption and all results of the generated output used by me. In the list ``Overview of tools used'' I have named all generative models used with their product name and indicated how, to what extent and for what purpose they were used.
}

\end{document}
